% Coordonnées d'un vecteur — Mathématiques 3ème — Cours signé OnBuch+
% Programme officiel MINESEC. Compiler avec : tectonic N08-coordonnees-vecteur.tex
\newcommand{\DOCMATIERE}{Mathématiques}
\newcommand{\DOCNIVEAU}{3ème}
\newcommand{\DOCMODULE}{Mathématiques – classe de 3ème}
\newcommand{\DOCLECON}{N08}
\newcommand{\DOCTITRE}{Coordonnées d'un vecteur}
% OnBuch+ — préambule « cours signés » ENRICHI (compilé avec tectonic / XeLaTeX)
% Système visuel « Pop » d'OnBuch+ : fond crème (jamais blanc), bordures encre
% épaisses, ombres « dures » décalées, titres Archivo Black en capitales.
% Version MATHÉMATIQUES : graphes (pgfplots/tikz), boîtes pédagogiques
% (définition, propriété, méthode, attention, le savais-tu, exercice résolu,
% exercice, TP, bilan), page de garde et signature OnBuch+ ancrée en bas.
%
% Macros attendues dans main.tex AVANT \input{preamble} :
%   \DOCMATIERE  \DOCNIVEAU  \DOCMODULE  \DOCLECON (numéro)  \DOCTITRE
\documentclass[11pt]{article}

\usepackage{fontspec}
\usepackage[french]{babel}
\usepackage[a4paper,margin=2cm,top=3.4cm,bottom=2.8cm,headheight=1.4cm,footskip=1.3cm]{geometry}
\usepackage{amsmath,amssymb,amsfonts}
\usepackage{mathtools} % \xmapsto, \coloneqq... souvent utilisés par les modèles
\usepackage{cancel} % simplifications barrées
% \abs{z} (module, valeur absolue) et \norm{u} (norme), utilisés par les modèles
\providecommand{\abs}{}\let\abs\relax\DeclarePairedDelimiter{\abs}{\lvert}{\rvert}
\providecommand{\norm}{}\let\norm\relax\DeclarePairedDelimiter{\norm}{\lVert}{\rVert}
\usepackage[table]{xcolor} % [table] : fournit aussi \rowcolors (alternance de lignes)
\usepackage{pagecolor}
\usepackage{tikz}
\usetikzlibrary{arrows.meta,positioning,calc,shapes.geometric,shapes.misc,shapes.arrows,decorations.pathmorphing,decorations.pathreplacing,decorations.markings,patterns,shadows,mindmap,trees,fit,backgrounds,matrix,intersections,angles,quotes}
\usepackage{pgfplots}
\usepackage{pgf-pie} % diagrammes circulaires \pie (statistiques)
\pgfplotsset{compat=1.18}
\usepgfplotslibrary{fillbetween,groupplots} % aires entre courbes (calcul intégral)
\usepackage{enumitem}
\usepackage{fancyhdr}
% [most] charge toutes les bibliothèques tcolorbox (listings, minted, raster,
% documentation...) et gonfle inutilement la mémoire TeX sur un document long
% (~300 boîtes) : on ne charge que ce qui est réellement utilisé.
\usepackage[skins,breakable]{tcolorbox}
\usepackage{graphicx}
\usepackage{colortbl}
\usepackage{booktabs}
\usepackage{tabularx}
\usepackage{diagbox} % case d'en-tête diagonale des tableaux à double entrée
% Colonnes extensibles centrée / gauche / droite (C, L, R), souvent utilisées par les modèles
\newcolumntype{C}{>{\centering\arraybackslash}X}
\newcolumntype{L}{>{\raggedright\arraybackslash}X}
\newcolumntype{R}{>{\raggedleft\arraybackslash}X}
\usepackage{array}
\usepackage{multicol}
\usepackage{siunitx}
% parse-numbers=false : accepte aussi \SI{2,5 \times 10^{-3}}{...}
\sisetup{locale=FR,output-decimal-marker={,},group-minimum-digits=4,parse-numbers=false}
\DeclareSIUnit\ohm{\ensuremath{\symup{\Omega}}} % U+2126 absent de la police
\usepackage{qrcode}
% \degree : commande fréquemment utilisée par les modèles pour un angle ou une
% température en dehors de \SI{}{} (non fournie par les paquets chargés).
\providecommand{\degree}{^{\circ}}
% \qed (fin de démonstration), utilisé par les modèles sans amsthm
\providecommand{\qed}{\hfill$\square$}
% \barre : double barre des valeurs interdites dans un tableau de variations (tkz-tab)
\providecommand{\barre}{\ensuremath{\|}}
% environnement proof (amsthm non chargé) utilisé par les modèles
\ifdefined\proof\else\newenvironment{proof}[1][Démonstration]{\par\noindent\textit{#1.}\ }{\qed\par}\fi
\usepackage{lastpage}
\usepackage{hyperref}
\hypersetup{hidelinks}

% ============================================================
% Palette « Pop » OnBuch+ — mêmes hex que la classe P (plus_ui.dart)
% ============================================================
\definecolor{poporange}{HTML}{F59321}
\definecolor{popdark}{HTML}{E07A0C}
\definecolor{popcream}{HTML}{FFF6E8}
\definecolor{popink}{HTML}{1C1714}
\definecolor{poppurple}{HTML}{7A5AE0}
\definecolor{popgreen}{HTML}{2E9E6A}
\definecolor{poppink}{HTML}{E0457B}
\definecolor{popblue}{HTML}{2F7DE1}
\definecolor{popgold}{HTML}{C98A12}
\definecolor{popmuted}{HTML}{8A7F76}
\definecolor{popline}{HTML}{EADFD2}
\definecolor{popgray}{HTML}{8A7F76}
\colorlet{popgrey}{popgray}
% Alias tolérants (le modèle nomme parfois les couleurs différemment)
\colorlet{popwhite}{white}
\colorlet{popblack}{popink}
\colorlet{popred}{poppink}
\colorlet{popyellow}{popgold}
% Teintes claires pour fonds de tableaux / zones
\colorlet{popblueL}{popblue!10!white}
\colorlet{poporangeL}{poporange!14!white}
\colorlet{popgreenL}{popgreen!12!white}
\colorlet{poppurpleL}{poppurple!12!white}
\colorlet{poppinkL}{poppink!12!white}
\colorlet{popgoldL}{popgold!14!white}

\pagecolor{popcream}

% ============================================================
% Typographie
% ============================================================
\defaultfontfeatures{Ligatures=TeX}
\setmainfont{PlusJakartaSans-Medium.ttf}[Path=../../../fonts/,BoldFont=PlusJakartaSans-Bold.ttf]
% Maths en sans-serif (Fira Math) avec chiffres et lettres droites de Plus
% Jakarta, pour que les formules se fondent dans le texte.
\usepackage[math-style=ISO,bold-style=ISO]{unicode-math}
\setmathfont{FiraMath-Regular.otf}[Path=../../../fonts/]
\setmathfont{PlusJakartaSans-Medium.ttf}[Path=../../../fonts/,range={up/num,up/{Latin,latin},\mathup}]
\usepackage{icomma} % virgule décimale sans espace en mode math
% Caractères mathématiques Unicode absents des polices de texte (Plus Jakarta,
% Archivo Black) : rendus par leur équivalent mathématique, en texte comme en
% formule — sinon ils s'affichent en carré vide (ex. « Arithmétique dans ℤ »).
\usepackage{newunicodechar}
\newunicodechar{ℝ}{\ensuremath{\mathbb{R}}}
\newunicodechar{ℤ}{\ensuremath{\mathbb{Z}}}
\newunicodechar{ℂ}{\ensuremath{\mathbb{C}}}
\newunicodechar{ℕ}{\ensuremath{\mathbb{N}}}
\newunicodechar{ℚ}{\ensuremath{\mathbb{Q}}}
\newunicodechar{𝔻}{\ensuremath{\mathbb{D}}}
\newunicodechar{α}{\ensuremath{\alpha}}
\newunicodechar{β}{\ensuremath{\beta}}
\newunicodechar{γ}{\ensuremath{\gamma}}
\newunicodechar{δ}{\ensuremath{\delta}}
\newunicodechar{ε}{\ensuremath{\varepsilon}}
\newunicodechar{θ}{\ensuremath{\theta}}
\newunicodechar{λ}{\ensuremath{\lambda}}
\newunicodechar{μ}{\ensuremath{\mu}}
\newunicodechar{σ}{\ensuremath{\sigma}}
\newunicodechar{τ}{\ensuremath{\tau}}
\newunicodechar{φ}{\ensuremath{\varphi}}
\newunicodechar{ψ}{\ensuremath{\psi}}
\newunicodechar{ω}{\ensuremath{\omega}}
\newunicodechar{Σ}{\ensuremath{\Sigma}}
% majuscules produites par \MakeUppercase dans les titres (α-aminés -> Α)
\newunicodechar{Α}{\ensuremath{\alpha}}
\newunicodechar{Β}{\ensuremath{\beta}}
\newunicodechar{Γ}{\ensuremath{\Gamma}}
\newunicodechar{Δ}{\ensuremath{\Delta}}
\newunicodechar{Ε}{\ensuremath{\varepsilon}}
\newunicodechar{Θ}{\ensuremath{\Theta}}
\newunicodechar{Λ}{\ensuremath{\Lambda}}
\newunicodechar{Μ}{\ensuremath{\mu}}
\newunicodechar{Τ}{\ensuremath{\tau}}
\newunicodechar{Φ}{\ensuremath{\Phi}}
\newunicodechar{Ψ}{\ensuremath{\Psi}}
\newunicodechar{Ω}{\ensuremath{\Omega}}
\newunicodechar{↦}{\ensuremath{\mapsto}}
\newunicodechar{∈}{\ensuremath{\in}}
\newunicodechar{∉}{\ensuremath{\notin}}
\newunicodechar{∧}{\ensuremath{\wedge}}
\newunicodechar{⇔}{\ensuremath{\Leftrightarrow}}
\newunicodechar{⟺}{\ensuremath{\Longleftrightarrow}}
\newunicodechar{⇒}{\ensuremath{\Rightarrow}}
\newunicodechar{⟹}{\ensuremath{\Longrightarrow}}
\newunicodechar{∘}{\ensuremath{\circ}}
\newunicodechar{∪}{\ensuremath{\cup}}
\newunicodechar{∩}{\ensuremath{\cap}}
\newunicodechar{≡}{\ensuremath{\equiv}}
\newunicodechar{⊂}{\ensuremath{\subset}}
\newunicodechar{⊥}{\ensuremath{\perp}}
\newunicodechar{∃}{\ensuremath{\exists}}
\newunicodechar{∀}{\ensuremath{\forall}}
\newunicodechar{∛}{\ensuremath{\sqrt[3]{\,}}}
\newunicodechar{∜}{\ensuremath{\sqrt[4]{\,}}}
\newunicodechar{⃗}{\ensuremath{{}^{\to}}}
\newunicodechar{ⁿ}{\ifmmode^{n}\else\textsuperscript{\ensuremath{n}}\fi}
\newunicodechar{ˣ}{\ifmmode^{x}\else\textsuperscript{\ensuremath{x}}\fi}
\newunicodechar{ᵇ}{\ifmmode^{b}\else\textsuperscript{\ensuremath{b}}\fi}
\newunicodechar{ʳ}{\ifmmode^{r}\else\textsuperscript{\ensuremath{r}}\fi}
\newunicodechar{ᵘ}{\ifmmode^{u}\else\textsuperscript{\ensuremath{u}}\fi}
\newunicodechar{ʸ}{\ifmmode^{y}\else\textsuperscript{\ensuremath{y}}\fi}
\newunicodechar{ᵃ}{\ifmmode^{a}\else\textsuperscript{\ensuremath{a}}\fi}
\newunicodechar{ᵉ}{\ifmmode^{e}\else\textsuperscript{\ensuremath{e}}\fi}
\newunicodechar{ᵏ}{\ifmmode^{k}\else\textsuperscript{\ensuremath{k}}\fi}
\newunicodechar{ᵖ}{\ifmmode^{p}\else\textsuperscript{\ensuremath{p}}\fi}
\newunicodechar{ᴺ}{\ifmmode^{N}\else\textsuperscript{\ensuremath{N}}\fi}
\newunicodechar{ᶜ}{\ifmmode^{c}\else\textsuperscript{\ensuremath{c}}\fi}
\newunicodechar{ʰ}{\ifmmode^{h}\else\textsuperscript{\ensuremath{h}}\fi}
\newunicodechar{ᵠ}{\ifmmode^{q}\else\textsuperscript{\ensuremath{q}}\fi}
\newunicodechar{ₙ}{\ifmmode_{n}\else\textsubscript{\ensuremath{n}}\fi}
\newunicodechar{ₐ}{\ifmmode_{a}\else\textsubscript{\ensuremath{a}}\fi}
\newunicodechar{ᵢ}{\ifmmode_{i}\else\textsubscript{\ensuremath{i}}\fi}
\newunicodechar{ᵣ}{\ifmmode_{r}\else\textsubscript{\ensuremath{r}}\fi}
\newunicodechar{ₖ}{\ifmmode_{k}\else\textsubscript{\ensuremath{k}}\fi}
\newunicodechar{ᵦ}{\ifmmode_{\beta}\else\textsubscript{\ensuremath{\beta}}\fi}
\newunicodechar{ₓ}{\ifmmode_{x}\else\textsubscript{\ensuremath{x}}\fi}
\newfontfamily\popdisplay{ArchivoBlack-Regular.ttf}[Path=../../../fonts/]
\newfontfamily\poplogo{SpaceGrotesk-Bold.ttf}[Path=../../../fonts/]
\newfontfamily\popsemi{PlusJakartaSans-SemiBold.ttf}[Path=../../../fonts/]
\newfontfamily\popxbold{PlusJakartaSans-ExtraBold.ttf}[Path=../../../fonts/]
\newfontfamily\popxboldls{PlusJakartaSans-ExtraBold.ttf}[Path=../../../fonts/,LetterSpace=3.0]

\setlist[enumerate]{leftmargin=1.7em,itemsep=5pt,topsep=4pt}
\setlist[itemize]{leftmargin=1.7em,itemsep=4pt,topsep=4pt,label={\color{poporange}\textbullet}}
\setlength{\parindent}{0pt}
\setlength{\parskip}{5pt}
\renewcommand{\arraystretch}{1.25}

% pgfplots : style Pop par défaut
\pgfplotsset{
  every axis/.append style={axis line style={popink,line width=1pt},tick style={popink},
    grid style={popline,line width=0.5pt},label style={font=\small},tick label style={font=\footnotesize}},
  popcurve/.style={poporange,line width=1.8pt,no markers,smooth},
}
% Même style disponible dans un \draw TikZ simple (hors pgfplots)
\tikzset{popcurve/.style={poporange,line width=1.8pt}}

% ============================================================
% Pill / squiggle
% ============================================================
\newcommand{\poppill}[3]{%
  \tikz[baseline=-0.55ex]\node[draw=popink,line width=1.2pt,rounded corners=2.6pt,%
    fill=#2,inner xsep=6.5pt,inner ysep=3pt]{%
    \popxboldls\fontsize{7}{7}\selectfont\color{#3}\MakeUppercase{#1}};%
}
\newcommand{\popsquiggle}[1]{%
  \tikz[baseline=-0.4ex]{
    \draw[#1,line width=2.6pt,line cap=round]
      plot [smooth] coordinates {(0,0.09) (0.34,0.01) (0.68,0.21) (1.02,0.01) (1.36,0.10)};
  }%
}
% Mot-clé surligné (marqueur orange sous le texte)
\newcommand{\cle}[1]{{\popxbold\color{popdark}#1}}

% ============================================================
% En-tête / pied
% ============================================================
\pagestyle{fancy}
\fancyhf{}
\renewcommand{\headrulewidth}{0pt}
\renewcommand{\footrulewidth}{0pt}
\lhead{\raisebox{-0.15ex}{\poplogo\fontsize{13}{13}\selectfont\color{popink}OnBuch\color{poporange}+}}
\rhead{\poppill{\DOCMATIERE\ \textbullet\ \DOCNIVEAU\ \textbullet\ Leçon \DOCLECON}{white}{popink}}
\lfoot{\popxbold\fontsize{7.6}{7.6}\selectfont\color{popmuted}\MakeUppercase{Cours signé OnBuch+}\ \textbullet\ {\popsemi\DOCTITRE}}
\rfoot{\tikz[baseline=-0.6ex]\node[rounded corners=8.5pt,fill=popink,minimum height=17pt,minimum width=17pt,inner xsep=5pt,inner sep=0]{\popxbold\fontsize{8}{8}\selectfont\color{popcream}\thepage\,/\,\pageref*{LastPage}};}

% ============================================================
% Titres
% ============================================================
\newcommand{\chaptitre}[2]{%
  \begin{tcolorbox}[enhanced,colback=poporange,colframe=popink,boxrule=2.6pt,arc=11pt,
      drop shadow=popink,left=15pt,right=15pt,top=12pt,bottom=11pt,before skip=3pt,after skip=18pt]
    \poppill{#2}{popink}{popcream}\par\vspace{9pt}
    {\raggedright\hyphenpenalty=10000\exhyphenpenalty=10000\popdisplay\fontsize{22}{24}\selectfont\color{popcream}\MakeUppercase{#1}\par}
  \end{tcolorbox}
}
% Section numérotée : pastille orange avec le numéro + titre Archivo
\newcounter{popsec}
\newcommand{\coursec}[1]{%
  \stepcounter{popsec}\setcounter{popsub}{0}%
  \par\vspace{16pt}\noindent
  \begin{minipage}{\linewidth}
  \tikz[baseline=-0.7ex]\node[draw=popink,line width=1.6pt,fill=poporange,rounded corners=4pt,minimum size=22pt,inner sep=0]{\popdisplay\fontsize{12}{12}\selectfont\color{popcream}\thepopsec};%
  \hspace{8pt}{\popdisplay\fontsize{15}{17}\selectfont\color{popink}\MakeUppercase{#1}}\par
  \vspace{4pt}\noindent{\color{poporange}\rule{36pt}{3.6pt}}
  \end{minipage}\par\nopagebreak\vspace{8pt}
}
\newcounter{popsub}
\newcommand{\courssub}[1]{%
  \stepcounter{popsub}%
  \par\vspace{9pt}\noindent{\popxbold\fontsize{11.5}{13}\selectfont\color{popink}{\color{poporange}\thepopsec.\thepopsub}\hspace{6pt}#1}\par\nopagebreak\vspace{4pt}
}

% ============================================================
% Boîtes pédagogiques — même recette : fond blanc, bordure épaisse colorée,
% ombre dure encre, badge plein. Titre optionnel [#1].
% ============================================================
\tcbset{popbase/.style={breakable,enhanced,colback=white,boxrule=2.3pt,arc=9pt,
  shadow={3pt}{-3pt}{0pt}{popink},left=14pt,right=14pt,top=11pt,bottom=11pt,before skip=11pt,after skip=15pt,
  fontupper=\popsemi\fontsize{10.6}{14.5}\selectfont\color{popink},
  fontlower=\popsemi\fontsize{10.6}{14.5}\selectfont\color{popink}}}
% Coche dessinée (le glyphe ✓ n'existe pas dans les polices du document)
\newcommand{\popcheck}[1]{\tikz[baseline=-0.5ex]\draw[#1,line width=1.6pt,line cap=round,line join=round](-0.13,0.0)--(-0.04,-0.09)--(0.13,0.1);}
\newcommand{\popboxhead}[3]{\poppill{#1}{#2}{white}\ifx\relax#3\relax\else\hspace{6pt}{\popxbold\fontsize{10.6}{12}\selectfont\color{popink}#3}\fi\par\vspace{7pt}}

\newtcolorbox{aretenir}[1][]{popbase,colframe=popgreen,before upper={\popboxhead{À retenir}{popgreen}{#1}}}
\newtcolorbox{exemplebox}[1][]{popbase,colframe=poporange,before upper={\popboxhead{Exemple}{poporange}{#1}}}
\newtcolorbox{definition}[1][]{popbase,colframe=popblue,before upper={\popboxhead{Définition}{popblue}{#1}}}
\newtcolorbox{propriete}[1][]{popbase,colframe=poppurple,before upper={\popboxhead{Propriété}{poppurple}{#1}}}
\newtcolorbox{methode}[1][]{popbase,colframe=popgold,before upper={\popboxhead{Méthode}{popgold}{#1}}}
\newtcolorbox{attention}[1][]{popbase,colframe=poppink,before upper={\popboxhead{Attention}{poppink}{#1}}}
\newtcolorbox{experience}[1][]{popbase,colframe=popblue,colback=popblueL,before upper={\popboxhead{Expérience}{popblue}{#1}}}
% « Le savais-tu ? » : carte violette pleine (contexte, culture, Cameroun)
\newtcolorbox{savaistu}[1][]{popbase,colback=poppurple,colframe=popink,
  fontupper=\popsemi\fontsize{10.4}{14.2}\selectfont\color{white},
  before upper={\poppill{Le savais-tu\,?}{white}{poppurple}\ifx\relax#1\relax\else\hspace{6pt}{\popxbold\color{white}#1}\fi\par\vspace{7pt}}}
% Exercice résolu : énoncé en haut, \tcblower puis la solution sur fond crème
\newcounter{popexo}\newcounter{popexr}
\newtcolorbox{exoresolu}[1][]{popbase,colframe=poporange,colbacklower=popcream,
  segmentation style={popink,line width=1.2pt,dashed},
  before upper={\stepcounter{popexr}\popboxhead{Exercice résolu \thepopexr}{poporange}{#1}},
  before lower={\poppill{Solution}{popink}{popcream}\par\vspace{6pt}}}
% Exercice (non résolu en place) — la correction est dans la partie Corrigés
\newtcolorbox{exercice}[1][]{popbase,colframe=popink,
  before upper={\stepcounter{popexo}\popboxhead{Exercice \thepopexo}{popink}{#1}}}
% Corrigé
\newtcolorbox{corrige}[1][]{popbase,colframe=popgreen,colback=popgreenL,
  before upper={\popboxhead{Corrigé}{popgreen}{#1}}}
% Objectifs / prérequis / bilan
\newtcolorbox{objectifs}{popbase,colframe=popink,colback=poporangeL,
  before upper={\popboxhead{Objectifs de la leçon}{popink}{}}}
\newtcolorbox{prerequis}{popbase,colframe=popink,colback=popblueL,
  before upper={\popboxhead{Prérequis}{popblue}{}}}
\newtcolorbox{bilan}{popbase,colframe=popink,colback=white,boxrule=2.8pt,
  before upper={\popboxhead{Fiche bilan}{poporange}{L'essentiel à connaître pour l'examen}}}
% Figure encadrée avec légende
\newtcolorbox{popfigure}[1][]{enhanced,breakable=false,colback=white,colframe=popline,boxrule=1.4pt,arc=8pt,
  left=10pt,right=10pt,top=10pt,bottom=8pt,before skip=10pt,after skip=12pt,halign=center,
  lower separated=false,halign lower=center,
  fontlower=\popsemi\fontsize{9}{11.5}\selectfont\color{popmuted},
  #1}
\newcommand{\legende}[1]{\par\vspace{6pt}{\popsemi\fontsize{9}{11.5}\selectfont\color{popmuted}#1}}

% ============================================================
% Signature OnBuch+
% ============================================================
\newcommand{\poplogoinline}[1]{{\poplogo\fontsize{#1}{#1}\selectfont\color{popink}OnBuch\color{poporange}+}}
\newcommand{\signatureonbuch}{%
  \begin{tcolorbox}[enhanced,colback=popink,colframe=popink,boxrule=2.2pt,arc=10pt,
      drop shadow=poporange,left=14pt,right=14pt,top=10pt,bottom=10pt,before skip=0pt,after skip=0pt]
    \tikz[baseline=-0.7ex]\node[circle,fill=popgreen,draw=popcream,line width=1.3pt,minimum size=22pt,inner sep=0]{\popcheck{white}};%
    \hspace{9pt}\begin{minipage}[c]{0.62\linewidth}
      {\popxbold\fontsize{12}{13}\selectfont\color{popcream}Signé\ \textbullet\ {\poplogo OnBuch\color{poporange}+}}\par\vspace{2pt}
      {\raggedright\popsemi\fontsize{8}{10.5}\selectfont\color{popline}\MakeUppercase{Équipe pédagogique OnBuch+}\\\MakeUppercase{Conforme au programme officiel MINESEC}\par}
    \end{minipage}\hfill
    {\poplogo\fontsize{22}{22}\selectfont\color{popcream}OnBuch\color{poporange}+}
  \end{tcolorbox}%
}

% ============================================================
% Page de garde
% #1 titre · #2 matière · #3 niveau · #4 module · #5 n° leçon · #6 durée ·
% #7 description / objectifs (texte ou itemize)
% ============================================================
\newcommand{\pagegarde}[7]{%
  \thispagestyle{empty}
  \newgeometry{a4paper,margin=1.7cm,top=1.6cm,bottom=1.4cm}%
  \begin{tikzpicture}[remember picture,overlay]
    \fill[poporange!18] ([xshift=-3.2cm,yshift=-3.6cm]current page.north east) circle (4.6cm);
    \fill[poppurple!14] ([xshift=2.4cm,yshift=7.6cm]current page.south west) circle (3.8cm);
  \end{tikzpicture}%
  {\poplogo\fontsize{30}{30}\selectfont\color{popink}OnBuch\color{poporange}+}\hfill
  \raisebox{0.6ex}{\poppill{Cours signé \textbullet\ Premium}{popink}{popcream}}\par
  \vspace{1pt}\popsquiggle{poppurple}\par
  \vspace{8pt}
  \begin{tcolorbox}[enhanced,colback=poporange,colframe=popink,boxrule=2.8pt,arc=13pt,
      drop shadow=popink,left=18pt,right=18pt,top=12pt,bottom=13pt,after skip=10pt]
    \poppill{#2 \textbullet\ #3}{popink}{popcream}\hspace{5pt}\poppill{#4}{white}{popink}\par\vspace{8pt}
    {\popdisplay\fontsize{14}{14}\selectfont\color{popink}LEÇON #5}\par\vspace{4pt}
    {\raggedright\hyphenpenalty=10000\exhyphenpenalty=10000\popdisplay\fontsize{25}{27}\selectfont\color{popcream}\MakeUppercase{#1}\par}
    \vspace{8pt}
    {\popxbold\fontsize{9.5}{9.5}\selectfont\color{popink}\MakeUppercase{Durée conseillée : #6}}
  \end{tcolorbox}
  \begin{tcolorbox}[enhanced,colback=white,colframe=popink,boxrule=2.2pt,arc=10pt,
      drop shadow=popink,left=16pt,right=16pt,top=10pt,bottom=10pt,after skip=8pt]
    \poppill{Dans cette leçon}{poppurple}{white}\par\vspace{5pt}
    \popsemi\fontsize{9.3}{12.6}\selectfont\color{popink}#7
  \end{tcolorbox}
  \noindent\begin{minipage}[t]{0.487\linewidth}
    \begin{tcolorbox}[enhanced,colback=popgreen,colframe=popink,boxrule=2.2pt,arc=10pt,
        drop shadow=popink,left=11pt,right=11pt,top=10pt,bottom=10pt,height=3.1cm,valign=center]
      \tcbox[colback=white,colframe=popink,boxrule=1.3pt,arc=5pt,boxsep=0pt,nobeforeafter,
        left=3pt,right=3pt,top=3pt,bottom=3pt,valign=center]{\qrcode[height=1.9cm]{https://play.google.com/store/apps/details?id=com.luvvix.onbuch&hl=fr}}%
      \hspace{8pt}\begin{minipage}[c]{0.5\linewidth}
        {\popdisplay\fontsize{11}{12.5}\selectfont\color{white}\MakeUppercase{Télécharge OnBuch}}\par\vspace{4pt}
        {\raggedright\popsemi\fontsize{8.4}{11}\selectfont\color{white}Gratuit \textbullet\ Google Play\\Tuteur IA, annales, concours, résultats\par}
      \end{minipage}
    \end{tcolorbox}
  \end{minipage}\hfill
  \begin{minipage}[t]{0.487\linewidth}
    \begin{tcolorbox}[enhanced,colback=poppurple,colframe=popink,boxrule=2.2pt,arc=10pt,
        drop shadow=popink,left=11pt,right=11pt,top=10pt,bottom=10pt,height=3.1cm,valign=center]
      \tikz[baseline=-0.7ex]\node[circle,fill=white,draw=popink,line width=1.3pt,minimum size=1.6cm,inner sep=0]{\popdisplay\fontsize{24}{24}\selectfont\color{poppurple}+};%
      \hspace{8pt}\begin{minipage}[c]{0.6\linewidth}
        {\popdisplay\fontsize{11}{12.5}\selectfont\color{white}\MakeUppercase{Passe à OnBuch+}}\par\vspace{4pt}
        {\raggedright\popsemi\fontsize{8.4}{11}\selectfont\color{white}Tous les cours signés, Tuteur IA illimité, fiches PDF — onglet OnBuch+ de l'app\par}
      \end{minipage}
    \end{tcolorbox}
  \end{minipage}
  \vspace{8pt}
  \popcontenu
  \vfill
  \signatureonbuch
  \restoregeometry
  \newpage
}

% Rangée « Ce cours contient » (page de garde)
\newcommand{\poptile}[3]{% #1 couleur #2 gros texte #3 légende
  \begin{tcolorbox}[enhanced,colback=#1,colframe=popink,boxrule=2pt,arc=9pt,shadow={2.5pt}{-2.5pt}{0pt}{popink},
      left=6pt,right=6pt,top=9pt,bottom=9pt,halign=center,valign=center,height=1.75cm,width=\linewidth,nobeforeafter]
    {\popdisplay\fontsize{12.5}{13}\selectfont\color{white}\MakeUppercase{#2}}\par\vspace{3pt}
    {\centering\popsemi\fontsize{7.8}{9.5}\selectfont\color{white}#3\par}
  \end{tcolorbox}}
\newcommand{\popcontenu}{%
  \par\noindent{\popxboldls\fontsize{7.5}{7.5}\selectfont\color{popmuted}CE COURS SIGNÉ CONTIENT}\par\vspace{7pt}
  \noindent\begin{minipage}[t]{0.235\linewidth}\poptile{popblue}{Cours}{complet, illustré, pas à pas}\end{minipage}\hfill
  \begin{minipage}[t]{0.235\linewidth}\poptile{poporange}{Méthodes}{et exercices résolus}\end{minipage}\hfill
  \begin{minipage}[t]{0.235\linewidth}\poptile{poppink}{Exercices}{type examen + corrigés}\end{minipage}\hfill
  \begin{minipage}[t]{0.235\linewidth}\poptile{popgreen}{Fiche bilan}{l'essentiel à retenir}\end{minipage}\par}

% Sommaire
\newtcolorbox{sommaire}{enhanced,colback=white,colframe=popink,boxrule=2.2pt,arc=10pt,
  drop shadow=popink,left=18pt,right=18pt,top=13pt,bottom=13pt,before skip=6pt,after skip=20pt,
  before upper={\poppill{Sommaire}{poppurple}{white}\par\vspace{9pt}\popsemi\fontsize{10.6}{14.5}\selectfont\color{popink}}}

% Signature de fin de cours — poussée tout en bas de la dernière page
\newcommand{\findecours}{%
  \par\vfill\nopagebreak
  \begin{center}\popsquiggle{poporange}\end{center}\vspace{4pt}
  \signatureonbuch
}
\AtBeginDocument{\def\llbracket{\mathopen{[\mkern-3.5mu[}}\def\rrbracket{\mathclose{]\mkern-3.5mu]}}} % intervalles d'entiers (glyphe ⟦ absent de la police)
% Glyphes absents de Fira Math : redéfinis à partir de glyphes disponibles
\AtBeginDocument{%
  \def\setminus{\mathbin{\backslash}}\let\smallsetminus\setminus
  \def\vdots{\mathord{\vbox{\baselineskip3.2pt\lineskiplimit0pt\kern4pt\hbox{.}\hbox{.}\hbox{.}}}}%
  \def\ddots{\mathinner{\mkern1mu\raise6.5pt\hbox{.}\mkern2mu\raise3.6pt\hbox{.}\mkern2mu\raise0.7pt\hbox{.}\mkern1mu}}%
  \def\triangle{\mathord{\scalebox{0.9}{$\symup{\Delta}$}}}%
  \def\nabla{\mathord{\raisebox{\depth}{\scalebox{1}[-1]{$\symup{\Delta}$}}}}%
  \def\checkmark{\popcheck{popdark}}%
  \def\star{\mathbin{\ast}}\def\bigstar{\mathord{\ast}}%
  \def\diamond{\mathbin{\scalebox{0.6}{$\Diamond$}}}%
  \def\bigcup{\mathop{\mathchoice{\raisebox{-0.3em}{\scalebox{1.7}{$\cup$}}}{\scalebox{1.25}{$\cup$}}{\cup}{\cup}}\displaylimits}%
  \def\bigcap{\mathop{\mathchoice{\raisebox{-0.3em}{\scalebox{1.7}{$\cap$}}}{\scalebox{1.25}{$\cap$}}{\cap}{\cap}}\displaylimits}%
  \def\lesseqqgtr{\mathrel{\lessgtr}}\def\bigoplus{\mathop{\oplus}\displaylimits}%
}
\providecommand{\ssi}{si et seulement si} % abréviation fréquente
\AtBeginDocument{\providecommand{\wideparen}{\overparen}} % arc de cercle

%% FIGFIX-BEGIN v5 — correctifs globaux des figures (docs/onbuch-generation/tools/figfix)
\usepackage{adjustbox}\usepackage{etoolbox}\tcbuselibrary{raster}
% toute figure TikZ/pgfplots plus large que la ligne est réduite à la largeur disponible
\BeforeBeginEnvironment{tikzpicture}{\begin{adjustbox}{max width=\linewidth}}
\AfterEndEnvironment{tikzpicture}{\end{adjustbox}}
% légendes pgfplots lisibles par-dessus les courbes
\pgfplotsset{every axis/.append style={legend style={fill=white,fill opacity=0.92,text opacity=1,draw=black!15,inner sep=3pt}}}
% étiquettes d'arêtes (syntaxe edge["..."]) avec fond blanc
\tikzset{every edge quotes/.append style={fill=white,inner sep=1.2pt}}
%% FIGFIX-END
\begin{document}

\pagegarde{Coordonnées d'un vecteur}{Mathématiques}{3ème}{Mathématiques – classe de 3ème}{N08}{4 séances (fiche de progression harmonisée)}{Cette leçon établit le lien fondamental entre la géométrie vectorielle et le calcul algébrique dans un repère orthonormé. Elle permet de traduire toute propriété géométrique (longueur, parallélisme, perpendicularité) en égalités numériques sur les coordonnées.
\vspace{4pt}
\begin{multicols}{2}\small
\begin{itemize}
\item À la fin de cette leçon, je sais définir les coordonnées d'un vecteur dans un repère orthonormé et les lire sur une figure.
\item Je sais calculer les coordonnées d'un vecteur AB à partir des coordonnées des points A et B.
\item Je sais calculer la distance entre deux points (longueur d'un vecteur) à partir de leurs coordonnées.
\item Je sais utiliser la condition de colinéarité de deux vecteurs pour prouver que des points sont alignés ou que des droites sont parallèles.
\item Je sais utiliser la condition d'orthogonalité de deux vecteurs (produit scalaire nul) pour prouver que des droites sont perpendiculaires.
\item Je sais rédiger une démonstration structurée en justifiant chaque étape par une propriété vectorielle ou algébrique.
\end{itemize}\end{multicols}}

\begin{sommaire}
\begin{multicols}{2}
\begin{enumerate}
\item 1. Coordonnées d'un vecteur dans un repère orthonormé
\item 2. Norme d'un vecteur et Distance entre deux points
\item 3. Condition de colinéarité de deux vecteurs
\item 4. Condition d'orthogonalité de deux vecteurs (Produit scalaire)
\item 5. Synthèse et Résolution de problèmes complexes
\item Activité d'intégration
\item Méthodes et exercices résolus
\item Exercices
\item Corrigés des exercices
\item Fiche bilan
\end{enumerate}
\end{multicols}
\end{sommaire}

\begin{objectifs}
\begin{itemize}
\item À la fin de cette leçon, je sais définir les coordonnées d'un vecteur dans un repère orthonormé et les lire sur une figure.
\item Je sais calculer les coordonnées d'un vecteur AB à partir des coordonnées des points A et B.
\item Je sais calculer la distance entre deux points (longueur d'un vecteur) à partir de leurs coordonnées.
\item Je sais utiliser la condition de colinéarité de deux vecteurs pour prouver que des points sont alignés ou que des droites sont parallèles.
\item Je sais utiliser la condition d'orthogonalité de deux vecteurs (produit scalaire nul) pour prouver que des droites sont perpendiculaires.
\item Je sais rédiger une démonstration structurée en justifiant chaque étape par une propriété vectorielle ou algébrique.
\item Je sais modéliser un problème concret (repérage, construction, déplacement) en choisissant un repère adapté et en exploitant les coordonnées.
\item Je sais identifier et corriger l'erreur fréquente consistant à confondre coordonnées d'un point et composantes d'un vecteur.
\end{itemize}
\end{objectifs}

\begin{prerequis}
\begin{itemize}
\item Notion de vecteur (direction, sens, norme) et opérations (addition, multiplication par un nombre) - 4ème/Debut 3ème.
\item Repérage dans le plan : repère orthonormé, coordonnées d'un point - 4ème.
\item Théorème de Pythagore et réciproque - 5ème/4ème.
\item Calcul algébrique : développement, factorisation, résolution d'équations - 4ème/3ème.
\item Définition du produit scalaire (approche géométrique : ||u|| ||v|| cos(θ)) - Debut 3ème (Leçon précédente).
\end{itemize}
\end{prerequis}

\newpage

\coursec{Coordonnées d'un vecteur dans un repère orthonormé}

Pour tracer le terrain du nouveau stade omnisports de Yaoundé, le géomètre-topographe doit repérer précisément les poteaux d'angle. Il utilise un système de repérage (un repère orthonormé) et note les coordonnées de chaque point : par exemple $A(2\,;\,-3)$ et $B(-1\,;\,4)$. Mais ces nombres, à eux seuls, ne disent pas tout. Comment, à partir de ces coordonnées, peut-il calculer la longueur exacte des côtés, vérifier que les angles sont droits ou que les lignes de touche sont parallèles ?

La clé, c'est de transformer les \emph{points} en \emph{vecteurs} : un déplacement d'un point à un autre devient alors un couple de nombres que l'on peut additionner, multiplier, comparer. C'est toute la puissance des \cle{coordonnées vectorielles}, et c'est ce que nous allons découvrir dans cette section.

\textbf{Problématique :} comment associer à tout vecteur du plan un couple de nombres, et comment calculer ces nombres à partir des coordonnées des points ?

\courssub{Rappel : le repère orthonormé $(O\,;\vec{i},\vec{j})$}

\begin{definition}[Repère orthonormé]
Un \cle{repère orthonormé} du plan est la donnée d'un point $O$ appelé \cle{origine} et de deux vecteurs $\vec{i}$ et $\vec{j}$ tels que :
\begin{itemize}
\item $\vec{i}$ et $\vec{j}$ ne sont pas colinéaires (ils ne portent pas la même direction) ;
\item les vecteurs $\vec{i}$ et $\vec{j}$ sont \textbf{orthogonaux} (repère \emph{orthogonal}) ;
\item $\vec{i}$ et $\vec{j}$ ont la \textbf{même longueur}, égale à $1$ (repère \emph{normé}) : on dit que ce sont des \cle{vecteurs unitaires}.
\end{itemize}
On note ce repère $(O\,;\vec{i},\vec{j})$. La droite $(O,\vec{i})$ est l'\cle{axe des abscisses}, la droite $(O,\vec{j})$ est l'\cle{axe des ordonnées}.
\end{definition}

\begin{aretenir}[Ce qu'il faut retenir du rappel]
\begin{itemize}
\item Tout point $M$ du plan est repéré par un unique couple $(x_M\,;\,y_M)$ : son \textbf{abscisse} et son \textbf{ordonnée}.
\item Les vecteurs de base ont pour coordonnées $\vec{i}(1\,;\,0)$ et $\vec{j}(0\,;\,1)$.
\item Le mot « orthonormé » résume tout : \emph{ortho} (axes perpendiculaires) + \emph{normé} (même unité de longueur sur les deux axes).
\end{itemize}
\end{aretenir}

\courssub{Définition : coordonnées d'un vecteur}

\textbf{Intuition.} Imagine que tu te déplaces sur le quadrillage du cahier : pour aller d'un point à un autre, tu avances d'un certain nombre de carreaux vers la droite (ou la gauche), puis tu montes d'un certain nombre de carreaux (ou tu descends). Tout déplacement du plan se décrit ainsi : une composante \textbf{horizontale} et une composante \textbf{verticale}. Ces deux nombres sont les coordonnées du vecteur.

\begin{definition}[Coordonnées d'un vecteur]
Soit $(O\,;\vec{i},\vec{j})$ un repère orthonormé. Pour tout vecteur $\vec{u}$ du plan, il existe un \textbf{unique} couple de nombres réels $(x\,;\,y)$ tel que :
\[ \vec{u} = x\,\vec{i} + y\,\vec{j}. \]
Le couple $(x\,;\,y)$ est appelé les \cle{coordonnées} du vecteur $\vec{u}$ dans le repère $(O\,;\vec{i},\vec{j})$. On note :
\[ \vec{u}(x\,;\,y). \]
Le nombre $x$ est la \cle{composante horizontale} (ou abscisse) de $\vec{u}$, et $y$ est sa \cle{composante verticale} (ou ordonnée).
\end{definition}

\begin{exemplebox}[Lecture graphique des coordonnées d'un vecteur]
Sur la figure ci-dessous, le vecteur $\vec{u}$ correspond à un déplacement de $3$ unités vers la droite et de $2$ unités vers le haut. On écrit donc :
\[ \vec{u} = 3\vec{i} + 2\vec{j} \qquad \text{c'est-à-dire} \qquad \vec{u}(3\,;\,2). \]
\end{exemplebox}

\begin{popfigure}
\begin{center}
\begin{tikzpicture}[scale=0.9]
\draw[popline!40, very thin] (-1.5,-1.5) grid (5.5,3.5);
\draw[->, thick, popdark] (-1.5,0) -- (5.5,0) node[below left] {$x$};
\draw[->, thick, popdark] (0,-1.5) -- (0,3.5) node[below left] {$y$};
\draw[->, very thick, popblue] (0,0) -- (1,0) node[midway, below] {$\vec{i}$};
\draw[->, very thick, popblue] (0,0) -- (0,1) node[midway, left] {$\vec{j}$};
\node[below left, popdark] at (0,0) {$O$};
\draw[->, very thick, poporange] (1,-1) -- (4,1) node[midway, above left] {$\vec{u}$};
\draw[->, thick, popgreen, dashed] (1,-1) -- (4,-1) node[midway, below] {$3\vec{i}$};
\draw[->, thick, poppurple, dashed] (4,-1) -- (4,1) node[midway, right] {$2\vec{j}$};
\fill[popink] (1,-1) circle (1.5pt);
\fill[popink] (4,1) circle (1.5pt);
\end{tikzpicture}
\end{center}
\legende{Le vecteur $\vec{u}(3\,;\,2)$ décomposé en ses composantes horizontale $3\vec{i}$ et verticale $2\vec{j}$.}
\end{popfigure}

\begin{aretenir}[Vecteurs remarquables]
\begin{itemize}
\item Le \cle{vecteur nul} a pour coordonnées $\vec{0}(0\,;\,0)$ : aucun déplacement.
\item Les vecteurs de base : $\vec{i}(1\,;\,0)$ et $\vec{j}(0\,;\,1)$.
\end{itemize}
\end{aretenir}

\courssub{Propriété fondamentale : coordonnées du vecteur $\overrightarrow{AB}$}

Voici le résultat central de cette leçon : il permet de passer des coordonnées des \emph{points} aux coordonnées du \emph{vecteur}.

\begin{propriete}[Coordonnées du vecteur $\overrightarrow{AB}$]
Dans un repère orthonormé $(O\,;\vec{i},\vec{j})$, soient $A(x_A\,;\,y_A)$ et $B(x_B\,;\,y_B)$ deux points. Alors le vecteur $\overrightarrow{AB}$ a pour coordonnées :
\[ \overrightarrow{AB}\begin{pmatrix} x_B - x_A \\ y_B - y_A \end{pmatrix} \qquad \text{que l'on note aussi} \qquad \overrightarrow{AB}(x_B - x_A\,;\,y_B - y_A). \]
Autrement dit : \textbf{coordonnées de l'arrivée moins coordonnées du départ}.
\end{propriete}

\begin{experience}[Démonstration guidée (à savoir refaire)]
On veut montrer que $\overrightarrow{AB}(x_B - x_A\,;\,y_B - y_A)$.
\begin{enumerate}
\item \textbf{Relation de Chasles.} En passant par l'origine $O$, on peut écrire :
\[ \overrightarrow{AB} = \overrightarrow{AO} + \overrightarrow{OB} = \overrightarrow{OB} - \overrightarrow{OA}. \]
\item \textbf{Vecteurs position.} Par définition des coordonnées d'un point, dire que $A(x_A\,;\,y_A)$ signifie exactement $\overrightarrow{OA} = x_A\,\vec{i} + y_A\,\vec{j}$. De même $\overrightarrow{OB} = x_B\,\vec{i} + y_B\,\vec{j}$.
\item \textbf{Substitution.} On remplace :
\begin{align*}
\overrightarrow{AB} &= \overrightarrow{OB} - \overrightarrow{OA} \\
&= \left(x_B\,\vec{i} + y_B\,\vec{j}\right) - \left(x_A\,\vec{i} + y_A\,\vec{j}\right) \\
&= (x_B - x_A)\,\vec{i} + (y_B - y_A)\,\vec{j}.
\end{align*}
\item \textbf{Conclusion.} Par définition des coordonnées d'un vecteur, $\overrightarrow{AB}(x_B - x_A\,;\,y_B - y_A)$. $\blacksquare$
\end{enumerate}
Géométriquement, cela revient à projeter le vecteur $\overrightarrow{AB}$ sur chaque axe : la composante horizontale mesure l'écart des abscisses, la composante verticale l'écart des ordonnées.
\end{experience}

\begin{methode}[Calculer les coordonnées d'un vecteur $\overrightarrow{AB}$]
\begin{enumerate}
\item Relever les coordonnées des deux points : $A(x_A\,;\,y_A)$ (départ) et $B(x_B\,;\,y_B)$ (arrivée).
\item Calculer la première coordonnée : $x = x_B - x_A$ (abscisse de l'\textbf{arrivée} moins abscisse du \textbf{départ}).
\item Calculer la deuxième coordonnée : $y = y_B - y_A$.
\item Écrire le résultat : $\overrightarrow{AB}(x_B - x_A\,;\,y_B - y_A)$.
\item (Contrôle) Vérifier graphiquement le signe de chaque coordonnée : un déplacement vers la droite donne $x>0$, vers le haut donne $y>0$.
\end{enumerate}
\end{methode}

\begin{exemplebox}[Exemple chiffré avec vérification graphique]
Soient $A(2\,;\,-3)$ et $B(-1\,;\,4)$. Calculons les coordonnées de $\overrightarrow{AB}$ :
\begin{align*}
x &= x_B - x_A = -1 - 2 = -3, \\
y &= y_B - y_A = 4 - (-3) = 4 + 3 = 7.
\end{align*}
Donc $\overrightarrow{AB}(-3\,;\,7)$.
\textbf{Vérification :} pour aller de $A$ à $B$, on recule de $3$ unités (abscisse négative, cohérent avec $x=-3$) et on monte de $7$ unités (coordonnée positive, cohérent avec $y=7$). La figure ci-dessous confirme ce résultat.
\end{exemplebox}

\begin{popfigure}
\begin{center}
\begin{tikzpicture}[scale=0.62]
\draw[popline!40, very thin] (-3.5,-4.5) grid (3.5,5.5);
\draw[->, thick, popdark] (-3.5,0) -- (3.5,0) node[below left] {$x$};
\draw[->, thick, popdark] (0,-4.5) -- (0,5.5) node[below left] {$y$};
\draw[->, very thick, popblue] (0,0) -- (1,0) node[midway, below] {$\vec{i}$};
\draw[->, very thick, popblue] (0,0) -- (0,1) node[midway, left] {$\vec{j}$};
\node[below left, popdark] at (0,0) {$O$};
\fill[popink] (2,-3) circle (2pt) node[below right] {$A(2\,;\,-3)$};
\fill[popink] (-1,4) circle (2pt) node[above left] {$B(-1\,;\,4)$};
\draw[->, very thick, poporange] (2,-3) -- (-1,4) node[midway, above right] {$\overrightarrow{AB}$};
\draw[->, thick, popgreen, dashed] (2,-3) -- (-1,-3) node[midway, below] {$-3$};
\draw[->, thick, poppurple, dashed] (-1,-3) -- (-1,4) node[midway, left] {$+7$};
\end{tikzpicture}
\end{center}
\legende{De $A(2\,;\,-3)$ à $B(-1\,;\,4)$ : on recule de $3$ et on monte de $7$, donc $\overrightarrow{AB}(-3\,;\,7)$.}
\end{popfigure}

\begin{attention}[Erreurs fréquentes à éviter absolument]
\begin{itemize}
\item \textbf{Inverser l'ordre.} L'erreur n°1 consiste à calculer $x_A - x_B$ au lieu de $x_B - x_A$. Retiens le slogan : \cle{Arrivée moins Départ}. Si tu inverses, tu obtiens les coordonnées de $\overrightarrow{BA}$, l'opposé du vecteur demandé !
\item \textbf{Confondre point et vecteur.} Le point $A(2\,;\,-3)$ et le vecteur $\overrightarrow{OA}(2\,;\,-3)$ portent les mêmes nombres, mais ce ne sont \textbf{pas les mêmes objets} : le point est une position fixe, le vecteur est un déplacement. En revanche, le vecteur $\overrightarrow{AB}(-3\,;\,7)$ est un vecteur \emph{libre} : ses coordonnées ne dépendent pas de l'endroit où on le dessine.
\item \textbf{Oublier les signes.} Avec $y_A = -3$, on a $y_B - y_A = 4 - (-3) = 7$, et non $4 - 3 = 1$. Les doubles signes demandent la plus grande vigilance.
\end{itemize}
\end{attention}

\courssub{Égalité de deux vecteurs et opérations sur les coordonnées}

\begin{propriete}[Égalité de deux vecteurs]
Deux vecteurs $\vec{u}(x\,;\,y)$ et $\vec{v}(x'\,;\,y')$ sont égaux si et seulement si leurs coordonnées sont égales :
\[ \vec{u} = \vec{v} \iff x = x' \ \text{et}\ y = y'. \]
\end{propriete}

Cette propriété est très utile pour démontrer qu'un quadrilatère est un parallélogramme : $ABCD$ est un parallélogramme si et seulement si $\overrightarrow{AB} = \overrightarrow{DC}$, ce qui se traduit par deux égalités numériques.

\begin{propriete}[Opérations sur les coordonnées]
Soient $\vec{u}(x\,;\,y)$ et $\vec{v}(x'\,;\,y')$ deux vecteurs et $k$ un nombre réel.
\begin{itemize}
\item \textbf{Somme :} $\vec{u} + \vec{v}$ a pour coordonnées $(x + x'\,;\,y + y')$.
\item \textbf{Produit par un réel :} $k\vec{u}$ a pour coordonnées $(kx\,;\,ky)$.
\end{itemize}
\end{propriete}

\begin{exemplebox}[Calculs avec les coordonnées]
Soient $\vec{u}(3\,;\,-2)$ et $\vec{v}(-1\,;\,5)$.
\begin{itemize}
\item $\vec{u} + \vec{v}$ a pour coordonnées $(3 + (-1)\,;\,-2 + 5) = (2\,;\,3)$.
\item $4\vec{u}$ a pour coordonnées $(4 \times 3\,;\,4 \times (-2)) = (12\,;\,-8)$.
\item $-2\vec{v}$ a pour coordonnées $(2\,;\,-10)$.
\end{itemize}
\end{exemplebox}

\courssub{Application : trouver un point par translation}

\begin{exoresolu}[Déterminer les coordonnées d'un point image]
Dans un repère orthonormé, on donne le point $A(1\,;\,2)$ et le vecteur $\vec{u}(4\,;\,-3)$. Déterminer les coordonnées du point $M$ tel que $\overrightarrow{AM} = \vec{u}$.
\tcblower
\textbf{Solution détaillée.}

Notons $M(x_M\,;\,y_M)$ les coordonnées cherchées.

\textbf{Étape 1.} Exprimons les coordonnées de $\overrightarrow{AM}$ :
\[ \overrightarrow{AM}(x_M - x_A\,;\,y_M - y_A) \quad \text{soit} \quad \overrightarrow{AM}(x_M - 1\,;\,y_M - 2). \]

\textbf{Étape 2.} Traduisons l'égalité vectorielle $\overrightarrow{AM} = \vec{u}$. Deux vecteurs sont égaux si et seulement si leurs coordonnées sont égales :
\[
\begin{cases}
x_M - 1 = 4 \\
y_M - 2 = -3
\end{cases}
\]

\textbf{Étape 3.} Résolvons chaque équation :
\[
\begin{cases}
x_M = 4 + 1 = 5 \\
y_M = -3 + 2 = -1
\end{cases}
\]

\textbf{Conclusion :} $M(5\,;\,-1)$.

\textbf{Vérification :} $\overrightarrow{AM}(5-1\,;\,-1-2) = (4\,;\,-3) = \vec{u}$. Le résultat est correct.

\textbf{Remarque :} on dit que $M$ est l'image de $A$ par la \emph{translation} de vecteur $\vec{u}$. Concrètement, on « colle » le déplacement $\vec{u}$ au bout du point $A$.
\end{exoresolu}

\begin{methode}[Trouver $M$ tel que $\overrightarrow{AM} = \vec{u}$]
\begin{enumerate}
\item Poser $M(x\,;\,y)$ et écrire les coordonnées de $\overrightarrow{AM}$ : $(x - x_A\,;\,y - y_A)$.
\item Écrire le système obtenu en égalant avec les coordonnées de $\vec{u}$.
\item Résoudre les deux équations du premier degré.
\item Conclure et vérifier en recalculant $\overrightarrow{AM}$.
\end{enumerate}
\end{methode}

\courssub{Point, vecteur position, vecteur libre : la distinction cruciale}

\begin{aretenir}[Trois objets, trois sens différents]
\begin{itemize}
\item \textbf{Coordonnées du point $A$} : une \emph{position} fixe dans le repère, par exemple $A(2\,;\,-3)$.
\item \textbf{Coordonnées du vecteur position $\overrightarrow{OA}$} : le déplacement qui mène de l'origine $O$ au point $A$. Ce sont \emph{les mêmes nombres} que ceux du point $A$ : $\overrightarrow{OA}(2\,;\,-3)$. C'est pourquoi on dit que « dans un repère d'origine $O$, point et vecteur position se confondent numériquement ».
\item \textbf{Coordonnées d'un vecteur libre $\overrightarrow{AB}$} : un déplacement entre deux points quelconques, calculé par $(x_B - x_A\,;\,y_B - y_A)$. Ces coordonnées ne dépendent \emph{pas} de l'origine du repère : si l'on dessine le même vecteur ailleurs (un \emph{représentant} de $\vec{u}$), il garde les mêmes coordonnées.
\end{itemize}
\end{aretenir}

\begin{exercice}[Pour s'entraîner]
Dans un repère orthonormé $(O\,;\vec{i},\vec{j})$, on donne $A(-2\,;\,1)$, $B(3\,;\,-4)$ et $C(0\,;\,5)$.
\begin{enumerate}
\item Calculer les coordonnées des vecteurs $\overrightarrow{AB}$, $\overrightarrow{BA}$ et $\overrightarrow{AC}$.
\item Que remarque-t-on pour $\overrightarrow{AB}$ et $\overrightarrow{BA}$ ?
\item Déterminer les coordonnées du point $D$ tel que $\overrightarrow{CD} = \overrightarrow{AB}$.
\end{enumerate}
\end{exercice}

\begin{corrige}[Exercice 1]
\begin{enumerate}
\item $\overrightarrow{AB}(3-(-2)\,;\,-4-1) = (5\,;\,-5)$ ; $\overrightarrow{BA}(-2-3\,;\,1-(-4)) = (-5\,;\,5)$ ; $\overrightarrow{AC}(0-(-2)\,;\,5-1) = (2\,;\,4)$.
\item Les coordonnées de $\overrightarrow{BA}$ sont les opposées de celles de $\overrightarrow{AB}$ : $\overrightarrow{BA} = -\overrightarrow{AB}$, ce qui est normal car les deux vecteurs ont même longueur, même direction mais des sens contraires.
\item Posons $D(x\,;\,y)$. Alors $\overrightarrow{CD}(x - 0\,;\,y - 5) = (x\,;\,y-5)$. L'égalité $\overrightarrow{CD} = \overrightarrow{AB}$ donne $x = 5$ et $y - 5 = -5$, soit $y = 0$. Donc $D(5\,;\,0)$. Vérification : $\overrightarrow{CD}(5-0\,;\,0-5) = (5\,;\,-5) = \overrightarrow{AB}$. 
\end{enumerate}
\end{corrige}

\begin{savaistu}[Les vecteurs dans les jeux vidéo et la CAO]
En informatique graphique (jeux vidéo, conception assistée par ordinateur, effets spéciaux du cinéma), chaque sommet d'un objet est mémorisé comme un \emph{vecteur position} : un couple $(x\,;\,y)$ en 2D, ou un triplet $(x\,;\,y\,;\,z)$ en 3D. Pour déplacer un personnage à l'écran, l'ordinateur ajoute simplement un \emph{vecteur translation} aux coordonnées de tous ses sommets : exactement le calcul $\overrightarrow{AM} = \vec{u}$ que tu viens d'apprendre ! Au Cameroun, les jeunes développeurs de studios comme ceux de Yaoundé ou de Douala qui créent des jeux mobiles utilisent chaque jour ces coordonnées vectorielles, souvent sans même s'en rendre compte.
\end{savaistu}

\begin{aretenir}[Bilan de la section]
\begin{itemize}
\item Tout vecteur $\vec{u}$ s'écrit de manière unique $\vec{u} = x\vec{i} + y\vec{j}$ : on note $\vec{u}(x\,;\,y)$.
\item Si $A(x_A\,;\,y_A)$ et $B(x_B\,;\,y_B)$, alors $\overrightarrow{AB}(x_B - x_A\,;\,y_B - y_A)$ : \textbf{arrivée moins départ}.
\item $\vec{u}(x\,;\,y) = \vec{v}(x'\,;\,y') \iff x = x'$ et $y = y'$.
\item $\vec{u} + \vec{v}(x+x'\,;\,y+y')$ et $k\vec{u}(kx\,;\,ky)$.
\item Pour trouver $M$ tel que $\overrightarrow{AM} = \vec{u}$, on écrit un système de deux équations et on le résout.
\end{itemize}
Dans la prochaine section, nous utiliserons ces coordonnées pour calculer des \textbf{longueurs} : la norme d'un vecteur et la distance entre deux points — exactement ce dont le géomètre du stade de Yaoundé a besoin pour mesurer les côtés du terrain.
\end{aretenir}

\coursec{Norme d'un vecteur et Distance entre deux points}

Dans la section précédente, nous avons appris à repérer un vecteur $\vec{u}$ par ses coordonnées $\vec{u}(x ; y)$ dans un repère orthonormé $(O ; \vec{\imath}, \vec{\jmath})$. Nous savons maintenant que $x$ est l'abscisse (projection sur l'axe des abscisses) et $y$ l'ordonnée (projection sur l'axe des ordonnées). Mais un vecteur, c'est aussi une \emph{longueur} (une norme) et une \emph{direction}. Comment retrouver cette longueur à partir des simples nombres $x$ et $y$ ? C'est tout l'objet de cette section : passer des coordonnées à la mesure, grâce au théorème de Pythagore, que tu as étudié en classe de 4ème.

\courssub{Norme d'un vecteur dans un repère orthonormé}

\textbf{Intuition géométrique.}\par
Reprenons le vecteur $\vec{u} = \overrightarrow{AB}$ de coordonnées $\vec{u}(x ; y)$. Dans le repère, ce vecteur est l'hypoténuse d'un triangle rectangle dont les côtés sont parallèles aux axes. Les longueurs de ces côtés sont exactement $|x|$ et $|y|$. Le théorème de Pythagore nous donne alors immédiatement la longueur de $\vec{u}$.

\begin{popfigure}
\begin{tikzpicture}[scale=1.2, >=Stealth]
    % Axes
    \draw[popline, thick, ->] (-0.5,0) -- (5,0) node[right] {$x$};
    \draw[popline, thick, ->] (0,-0.5) -- (0,4) node[above] {$y$};
    \node[below left] at (0,0) {$O$};
    
    % Vecteur u
    \coordinate (A) at (1,1);
    \coordinate (B) at (4,3);
    \draw[popblue, very thick, ->] (A) -- (B) node[midway, above left] {$\vec{u}(x;y)$};
    \node[below left] at (A) {$A$};
    \node[above right] at (B) {$B$};
    
    % Projections (Triangle rectangle)
    \coordinate (H) at (4,1);
    \draw[poporange, dashed] (A) -- (H) node[midway, below] {$x$};
    \draw[poporange, dashed] (H) -- (B) node[midway, right] {$y$};
    
    % Angle droit
    \draw[poporange] (3.85,1) -- (3.85,1.15) -- (4,1.15);
    
    % Vecteurs i, j du repère
    \draw[popmuted, ->] (0,0) -- (1,0) node[midway, below] {$\vec{\imath}$};
    \draw[popmuted, ->] (0,0) -- (0,1) node[midway, left] {$\vec{\jmath}$};
\end{tikzpicture}
\legende{Construction du triangle rectangle associé au vecteur $\vec{u}(x;y)$. Les côtés de l'angle droit mesurent $|x|$ et $|y|$.}
\end{popfigure}

\begin{definition}[Norme d'un vecteur]
Dans un repère orthonormé, la \textbf{norme} d'un vecteur $\vec{u}$, notée $\|\vec{u}\|$, est la longueur de n'importe lequel de ses représentants : si $\vec{u} = \overrightarrow{AB}$, alors $\|\vec{u}\| = AB$.
\end{definition}

\begin{propriete}[Norme d'un vecteur à partir de ses coordonnées]
Soit $\vec{u}(x ; y)$ un vecteur dans un repère orthonormé $(O ; \vec{\imath}, \vec{\jmath})$.
La \textbf{norme} de $\vec{u}$ est donnée par la formule :
\[
\|\vec{u}\| = \sqrt{x^2 + y^2}.
\]
\end{propriete}

\begin{experience}[Démonstration guidée : le théorème de Pythagore au service des vecteurs]
On considère un vecteur $\vec{u}$ de coordonnées $(x ; y)$.
\begin{enumerate}
    \item On représente $\vec{u}$ par le vecteur $\overrightarrow{AB}$, avec $A(x_A ; y_A)$ et $B(x_B ; y_B)$. On a alors $x = x_B - x_A$ et $y = y_B - y_A$.
    \item On construit le point $H$ de coordonnées $(x_B ; y_A)$ : $\overrightarrow{AH}$ est horizontal et $\overrightarrow{HB}$ est vertical. Le triangle $AHB$ est donc rectangle en $H$.
    \item Les longueurs des côtés de l'angle droit sont $AH = |x_B - x_A| = |x|$ et $HB = |y_B - y_A| = |y|$.
    \item Dans le triangle $AHB$ rectangle en $H$, d'hypoténuse $[AB]$, le théorème de Pythagore donne :
    \[
    AB^2 = AH^2 + HB^2 \quad \text{donc} \quad \|\vec{u}\|^2 = |x|^2 + |y|^2 = x^2 + y^2.
    \]
    \item Comme une longueur est un nombre positif, on prend la racine carrée : $\|\vec{u}\| = \sqrt{x^2 + y^2}$.
\end{enumerate}
\emph{Remarque :} cette démonstration est exigible au BEPC. Il faut savoir la rédiger clairement en nommant les points et en invoquant le théorème de Pythagore.
\end{experience}

\begin{exemplebox}[Calcul de normes]
Calculer la norme des vecteurs suivants :
\begin{enumerate}
    \item $\vec{u}(3 ; 4)$
    \item $\vec{v}(-5 ; 12)$
    \item $\vec{w}(0 ; -7)$
\end{enumerate}
\textbf{Solution.}
\begin{enumerate}
    \item $\|\vec{u}\| = \sqrt{3^2 + 4^2} = \sqrt{9 + 16} = \sqrt{25} = 5$.
    \item $\|\vec{v}\| = \sqrt{(-5)^2 + 12^2} = \sqrt{25 + 144} = \sqrt{169} = 13$.
    \item $\|\vec{w}\| = \sqrt{0^2 + (-7)^2} = \sqrt{49} = 7$. \emph{(Vecteur vertical : sa norme est la valeur absolue de son ordonnée.)}
\end{enumerate}
\end{exemplebox}

\begin{attention}[Pièges classiques sur la norme]
\begin{itemize}
    \item \textbf{Oublier la racine carrée :} écrire $\|\vec{u}\| = x^2 + y^2$ est une erreur grave (confusion entre la norme et le carré de la norme).
    \item \textbf{Confondre $\sqrt{x^2 + y^2}$ et $x + y$ :} $\sqrt{x^2+y^2} \neq |x| + |y|$ en général. Par exemple $\sqrt{3^2+4^2}=5$ mais $|3|+|4|=7$.
    \item \textbf{Le signe des coordonnées :} comme on élève au carré, le signe de $x$ et $y$ n'a aucune importance pour la norme : $\|(-3;4)\| = \|(3;4)\| = 5$.
\end{itemize}
\end{attention}

\courssub{Distance entre deux points}

La notion de distance entre deux points est l'application directe de la norme d'un vecteur : la distance $AB$ est par définition la longueur du vecteur $\overrightarrow{AB}$, c'est-à-dire $AB = \|\overrightarrow{AB}\|$.

\begin{propriete}[Distance entre deux points]
Dans un repère orthonormé, soient deux points $A(x_A ; y_A)$ et $B(x_B ; y_B)$.
La distance $AB$ est donnée par la formule :
\[
AB = \sqrt{(x_B - x_A)^2 + (y_B - y_A)^2}.
\]
\end{propriete}

\begin{popfigure}
\begin{tikzpicture}[scale=1.2, >=Stealth]
    % Grille légère pour repérage
    \draw[popline!20] (-1,-1) grid (5,4);
    % Axes
    \draw[popline, thick, ->] (-1,0) -- (5.5,0) node[right] {$x$};
    \draw[popline, thick, ->] (0,-1) -- (0,4.5) node[above] {$y$};
    
    % Points A et B
    \coordinate (A) at (1,1);
    \coordinate (B) at (4,3.5);
    \fill[popblue] (A) circle (2pt) node[below left] {$A(x_A;y_A)$};
    \fill[popblue] (B) circle (2pt) node[above right] {$B(x_B;y_B)$};
    
    % Segment AB
    \draw[popblue, thick] (A) -- (B) node[midway, above left] {$AB$};
    
    % Triangle rectangle
    \coordinate (H) at (4,1);
    \draw[poporange, dashed, thick] (A) -- (H) node[midway, below] {$x_B - x_A$};
    \draw[poporange, dashed, thick] (H) -- (B) node[midway, right] {$y_B - y_A$};
    
    % Angle droit
    \draw[poporange] (3.85,1) -- (3.85,1.15) -- (4,1.15);
\end{tikzpicture}
\legende{Distance $AB$ : hypoténuse du triangle rectangle de côtés $|x_B-x_A|$ et $|y_B-y_A|$.}
\end{popfigure}

\begin{aretenir}[Formule clé : distance entre deux points]
\[
\boxed{AB = \sqrt{(x_B - x_A)^2 + (y_B - y_A)^2}}
\]
Cette formule est à \textbf{connaître par cœur}. Elle généralise le théorème de Pythagore au plan muni d'un repère orthonormé.
\end{aretenir}

\begin{exemplebox}[Calcul de distance : le triangle 3-4-5]
Soient les points $A(1 ; 2)$ et $B(4 ; 6)$. Calculer la distance $AB$.

\textbf{Résolution.}
\begin{align*}
AB &= \sqrt{(x_B - x_A)^2 + (y_B - y_A)^2} \\
   &= \sqrt{(4 - 1)^2 + (6 - 2)^2} \\
   &= \sqrt{3^2 + 4^2} \\
   &= \sqrt{9 + 16} \\
   &= \sqrt{25} = 5.
\end{align*}
La distance $AB$ vaut $5$ unités de longueur. On retrouve le triplet pythagoricien classique $(3, 4, 5)$.
\end{exemplebox}

\begin{attention}[Erreurs fréquentes sur la formule de distance]
\begin{itemize}
    \item \textbf{Ordre des soustractions :} on calcule \og coordonnée du point d'arrivée $-$ coordonnée du point de départ \fg. Comme on élève au carré, $(x_B-x_A)^2 = (x_A-x_B)^2$ : l'ordre ne change pas le résultat final, mais la rigueur veut qu'on suive le sens du vecteur $\overrightarrow{AB}$.
    \item \textbf{Parenthèses obligatoires avec les nombres négatifs :} si $A(-2; 3)$ et $B(1; -1)$, alors $x_B - x_A = 1 - (-2) = 3$. Écrire $1 - -2$ sans parenthèses est une faute de rédaction ; écrire $1 - 2 = -1$ est une erreur de calcul.
    \item \textbf{Oublier les carrés :} $\sqrt{(x_B-x_A) + (y_B-y_A)}$ est faux : chaque différence doit être élevée au carré.
\end{itemize}
\end{attention}

\courssub{Propriétés fondamentales de la norme}

La norme possède des propriétés essentielles qu'il faut connaître et savoir utiliser.

\begin{propriete}[Propriétés de la norme]
Soient $\vec{u}$ et $\vec{v}$ deux vecteurs du plan et $k$ un nombre réel.
\begin{enumerate}
    \item \textbf{Positivité :} $\|\vec{u}\| \ge 0$, et $\|\vec{u}\| = 0$ si et seulement si $\vec{u} = \vec{0}$.
    \item \textbf{Homogénéité :} $\|k\vec{u}\| = |k| \times \|\vec{u}\|$.
    \item \textbf{Inégalité triangulaire :} $\|\vec{u} + \vec{v}\| \le \|\vec{u}\| + \|\vec{v}\|$.
\end{enumerate}
\end{propriete}

\begin{experience}[Démonstration de l'homogénéité]
Soit $\vec{u}(x ; y)$. Alors le vecteur $k\vec{u}$ a pour coordonnées $(kx ; ky)$. On calcule :
\[
\|k\vec{u}\| = \sqrt{(kx)^2 + (ky)^2} = \sqrt{k^2x^2 + k^2y^2} = \sqrt{k^2(x^2+y^2)} = \sqrt{k^2}\times\sqrt{x^2+y^2} = |k| \times \|\vec{u}\|.
\]
On a utilisé $\sqrt{k^2} = |k|$, car une racine carrée est toujours positive.
\end{experience}

\begin{exemplebox}[Application de l'homogénéité]
Soit $\vec{u}$ un vecteur de norme $5$. Quelle est la norme de $-3\vec{u}$ ?
\[
\|-3\vec{u}\| = |-3| \times \|\vec{u}\| = 3 \times 5 = 15.
\]
Le vecteur opposé $-\vec{u}$ a la même norme que $\vec{u}$ ; multiplier le vecteur par $3$ multiplie sa norme par $3$.
\end{exemplebox}

\begin{savaistu}[L'inégalité triangulaire : la géométrie dans le calcul]
L'inégalité $\|\vec{u}+\vec{v}\| \le \|\vec{u}\|+\|\vec{v}\|$ traduit un fait géométrique que tu connais bien : \emph{le plus court chemin entre deux points est la ligne droite}. Dans un triangle, la longueur d'un côté est toujours inférieure ou égale à la somme des longueurs des deux autres côtés. L'égalité a lieu exactement lorsque les deux vecteurs ont la même direction et le même sens : les trois points sont alors alignés.
\end{savaistu}

\courssub{Applications géométriques : périmètre, médiane, triangle rectangle}

Les coordonnées permettent de résoudre des problèmes de géométrie entièrement par le calcul : c'est ce qu'on appelle la \cle{géométrie analytique}.

\begin{methode}[Vérifier si un triangle est rectangle par le calcul]
Pour un triangle $ABC$ dont on connaît les coordonnées des sommets :
\begin{enumerate}
    \item Calculer les \textbf{carrés} des trois longueurs : $AB^2$, $AC^2$, $BC^2$ (inutile de calculer les racines carrées !).
    \item Repérer le plus grand des trois carrés : il correspond au plus long côté, l'hypoténuse potentielle.
    \item Comparer : \textbf{le triangle est rectangle si et seulement si la somme des deux plus petits carrés est égale au plus grand carré} (réciproque du théorème de Pythagore).
    \item Conclure en nommant l'angle droit : c'est le sommet opposé au plus long côté.
\end{enumerate}
\end{methode}

\begin{exoresolu}[Triangle rectangle et périmètre]
Soient les points $A(2 ; -1)$, $B(5 ; 3)$ et $C(0 ; 4)$.
\begin{enumerate}
    \item Calculer les longueurs $AB$, $AC$ et $BC$.
    \item En déduire le périmètre du triangle $ABC$.
    \item Le triangle $ABC$ est-il rectangle ? Justifier.
\end{enumerate}
\tcblower
\textbf{Solution détaillée.}
\begin{enumerate}
    \item \textbf{Calcul des carrés des longueurs (plus rapide) :}
    \begin{align*}
    AB^2 &= (5-2)^2 + (3-(-1))^2 = 3^2 + 4^2 = 9 + 16 = 25, \quad \text{donc } AB = 5. \\
    AC^2 &= (0-2)^2 + (4-(-1))^2 = (-2)^2 + 5^2 = 4 + 25 = 29, \quad \text{donc } AC = \sqrt{29}. \\
    BC^2 &= (0-5)^2 + (4-3)^2 = (-5)^2 + 1^2 = 25 + 1 = 26, \quad \text{donc } BC = \sqrt{26}.
    \end{align*}
    \item \textbf{Périmètre :} $\mathcal{P} = AB + AC + BC = 5 + \sqrt{29} + \sqrt{26}$.
    Valeur approchée : $\sqrt{29} \approx 5,39$ et $\sqrt{26} \approx 5,10$, donc $\mathcal{P} \approx 15,49$ unités de longueur.
    \item \textbf{Test de la réciproque de Pythagore :} les carrés sont $25$, $26$ et $29$. Le plus grand est $AC^2 = 29$.
    Somme des deux autres : $AB^2 + BC^2 = 25 + 26 = 51 \neq 29$.
    \textbf{Conclusion :} le triangle $ABC$ n'est \textbf{pas} rectangle.
\end{enumerate}
\end{exoresolu}

\begin{exemplebox}[Longueur d'une médiane]
Dans le triangle $ABC$ précédent, calculer la longueur de la médiane issue de $A$.

\textbf{Rappel :} la médiane issue de $A$ est le segment joignant $A$ au milieu $I$ du côté $[BC]$.

\textbf{Résolution.}
\begin{enumerate}
    \item Coordonnées du milieu $I$ de $[BC]$ :
    $x_I = \dfrac{x_B+x_C}{2} = \dfrac{5+0}{2} = 2,5$ et $y_I = \dfrac{y_B+y_C}{2} = \dfrac{3+4}{2} = 3,5$.
    Donc $I(2,5 ; 3,5)$.
    \item Coordonnées du vecteur $\overrightarrow{AI}$ : $(x_I-x_A ; y_I-y_A) = (2,5-2 ; 3,5-(-1)) = (0,5 ; 4,5)$.
    \item Norme : $AI = \|\overrightarrow{AI}\| = \sqrt{0,5^2 + 4,5^2} = \sqrt{0,25 + 20,25} = \sqrt{20,5} = \dfrac{\sqrt{82}}{2} \approx 4,53$.
\end{enumerate}
\end{exemplebox}

\courssub{Cas particulier : distance d'un point à l'origine}

\begin{propriete}[Distance d'un point à l'origine]
Dans un repère orthonormé $(O ; \vec{\imath}, \vec{\jmath})$, la distance d'un point $M(x ; y)$ à l'origine $O$ est la norme du vecteur $\overrightarrow{OM}$ :
\[
OM = \|\overrightarrow{OM}\| = \sqrt{x^2 + y^2}.
\]
\end{propriete}

C'est exactement la formule de la norme du vecteur de coordonnées $(x ; y)$ vue au début de la section. Conséquence géométrique remarquable : le point $M(x;y)$ appartient au cercle de centre $O$ et de rayon $R$ si et seulement si $x^2 + y^2 = R^2$.

\courssub{Application concrète : distance \og à vol d'oiseau \fg sur une carte}

Imaginons une carte simplifiée de la ville de \textbf{Yaoundé} munie d'un repère orthonormé où l'unité est le kilomètre.
\begin{itemize}
    \item Le lycée de Bastos est représenté par le point $L(2 ; 5)$.
    \item Le rond-point de la Poste (centre-ville) est représenté par le point $P(8 ; 1)$.
\end{itemize}
Quelle est la distance \og à vol d'oiseau \fg (en ligne droite) entre ces deux lieux ?

\begin{exoresolu}[Distance sur une carte]
\textbf{Données :} $L(2;5)$ et $P(8;1)$, l'unité de longueur étant le kilomètre.

\textbf{Calcul :}
\[
LP = \sqrt{(8-2)^2 + (1-5)^2} = \sqrt{6^2 + (-4)^2} = \sqrt{36 + 16} = \sqrt{52} = 2\sqrt{13} \approx 7,21 \text{ km}.
\]
\tcblower
\textbf{Interprétation :} la distance en ligne droite entre le lycée de Bastos et le rond-point de la Poste est d'environ $7,2$ km. En voiture, par le réseau routier (boulevards, carrefours), le trajet réel est plus long (environ $9$ à $10$ km) : la distance calculée est la borne inférieure théorique. C'est le principe utilisé par les applications de navigation des téléphones (Google Maps, etc.) pour estimer une distance : elles calculent $\sqrt{(\Delta x)^2 + (\Delta y)^2}$ sur une projection plane de la carte.
\end{exoresolu}

\begin{savaistu}[Et dans l'espace ?]
La formule de distance étudiée ici est la version \textbf{plane} (deux dimensions) du théorème de Pythagore. Dans l'espace, avec un repère à trois axes $(x, y, z)$, la distance entre deux points $A(x_A;y_A;z_A)$ et $B(x_B;y_B;z_B)$ devient :
\[
AB = \sqrt{(x_B-x_A)^2 + (y_B-y_A)^2 + (z_B-z_A)^2}.
\]
C'est la même logique : on ajoute le carré de la différence sur le troisième axe. Les satellites GPS qui permettent de localiser un téléphone au Cameroun calculent en permanence des distances en trois dimensions (latitude, longitude, altitude) !
\end{savaistu}

\begin{exercice}[Entraînement : distances et triangles]
\begin{enumerate}
    \item Calculer la distance entre les points $E(-3 ; 2)$ et $F(1 ; -3)$.
    \item Soit un triangle $RST$ avec $R(0;0)$, $S(4;0)$ et $T(4;3)$.
    \begin{enumerate}
        \item Calculer les longueurs des trois côtés.
        \item Quelle est la nature de ce triangle ? Justifier.
        \item Calculer son périmètre et son aire.
    \end{enumerate}
    \item Dans un repère orthonormé, on donne $A(-2;1)$ et $B(4;4)$. Déterminer le point $M$ de l'axe des abscisses tel que $MA = MB$.
    \item (Défi) Un drone se déplace au-dessus d'un terrain plat repéré en 2D. Il part du point $D(0;0)$, passe par $E(3;4)$ puis par $F(8;4)$, avant de revenir en $D$. Quelle est la distance totale parcourue ?
\end{enumerate}
\end{exercice}

\begin{corrige}[Exercice n°1]
\begin{enumerate}
    \item $EF = \sqrt{(1-(-3))^2 + (-3-2)^2} = \sqrt{4^2 + (-5)^2} = \sqrt{16+25} = \sqrt{41}$.
    \item 
    \begin{enumerate}
        \item $RS = \sqrt{(4-0)^2+(0-0)^2} = 4$ ; $ST = \sqrt{(4-4)^2+(3-0)^2} = 3$ ; $RT = \sqrt{(4-0)^2+(3-0)^2} = \sqrt{25} = 5$.
        \item $RS^2+ST^2 = 16+9 = 25 = RT^2$. D'après la réciproque du théorème de Pythagore, le triangle $RST$ est \textbf{rectangle en $S$}.
        \item Périmètre $= 3+4+5 = 12$ unités. Aire $= \dfrac{RS \times ST}{2} = \dfrac{4 \times 3}{2} = 6$ unités d'aire.
    \end{enumerate}
    \item $M$ est sur l'axe des abscisses, donc $M(x;0)$.
    $MA^2 = (x-(-2))^2 + (0-1)^2 = (x+2)^2 + 1$.
    $MB^2 = (x-4)^2 + (0-4)^2 = (x-4)^2 + 16$.
    Comme des distances sont positives, $MA=MB \iff MA^2=MB^2$ :
    \begin{align*}
    (x+2)^2 + 1 &= (x-4)^2 + 16 \\
    x^2+4x+4+1 &= x^2-8x+16+16 \\
    4x + 5 &= -8x + 32 \\
    12x &= 27 \\
    x &= \frac{27}{12} = \frac{9}{4} = 2,25.
    \end{align*}
    Donc $M(2,25 ; 0)$.
    \item $DE = \sqrt{3^2+4^2} = 5$ (triangle 3-4-5). $EF = \sqrt{(8-3)^2+(4-4)^2} = 5$ (déplacement horizontal). $FD = \sqrt{8^2+4^2} = \sqrt{80} = 4\sqrt{5} \approx 8,94$.
    Distance totale $= 5 + 5 + 4\sqrt{5} = 10 + 4\sqrt{5} \approx 18,94$ unités de longueur.
\end{enumerate}
\end{corrige}

\coursec{Condition de colinéarité de deux vecteurs}

Après avoir vu comment calculer la norme d'un vecteur et la distance entre deux points, nous abordons maintenant un outil fondamental : \textbf{la colinéarité}. Elle permet de décider si deux vecteurs ont la même direction (ou des directions opposées) et, par conséquent, si trois points sont alignés ou si deux droites sont parallèles. C'est la clé de nombreux problèmes de géométrie analytique au BEPC.

\courssub{Rappel : définition de la colinéarité}

\begin{definition}[Vecteurs colinéaires]
Deux vecteurs $\vec{u}$ et $\vec{v}$ du plan sont dits \textbf{colinéaires} s'il existe un nombre réel $k$ tel que
\[
\vec{u} = k\,\vec{v} \qquad\text{ou}\qquad \vec{v} = k\,\vec{u}.
\]
Le réel $k$ s'appelle le \textbf{coefficient de proportionnalité}.
\end{definition}

\emph{Intuition :} deux vecteurs colinéaires sont portés par la même droite (ou par des droites parallèles). Ils pointent dans le même sens si $k>0$, dans le sens opposé si $k<0$. Le vecteur nul $\vec{0}$ est colinéaire à \emph{tout} vecteur (on peut prendre $k=0$).

\courssub{Condition algébrique : le déterminant nul}

\begin{propriete}[Théorème : condition de colinéarité]
Soient $\vec{u}(x;y)$ et $\vec{v}(x';y')$ deux vecteurs exprimés dans un repère orthonormé. Alors
\[
\vec{u} \text{ et } \vec{v} \text{ sont colinéaires} \;\Longleftrightarrow\; x\,y' - y\,x' = 0.
\]
L'expression $x\,y' - y\,x'$ est appelée le \textbf{déterminant} des deux vecteurs (noté $\det(\vec{u},\vec{v})$).
\end{propriete}

\begin{experience}[Démonstration guidée (double implication)]
\emph{Sens $\Rightarrow$ (si colinéaires alors déterminant nul).}  
Supposons $\vec{u}=k\vec{v}$ avec $k\in\mathbb{R}$. Alors $x=kx'$ et $y=ky'$. Calculons :
\[
x\,y' - y\,x' = (kx')y' - (ky')x' = kx'y' - ky'x' = 0.
\]

\emph{Sens $\Leftarrow$ (si déterminant nul alors colinéaires).}  
Supposons $x\,y' - y\,x' = 0$, c'est-à-dire $x\,y' = y\,x'$.  
\begin{itemize}
\item Si $\vec{v}=\vec{0}$ ($x'=y'=0$), alors $\vec{u}=k\vec{v}$ pour tout $k$ (par exemple $k=0$) : colinéaires.
\item Si $\vec{v}\neq\vec{0}$, au moins l'une des coordonnées $x',y'$ est non nulle.
      \begin{itemize}
      \item Si $x'\neq 0$, posons $k = \dfrac{x}{x'}$. L'égalité $x\,y' = y\,x'$ donne $y = \dfrac{x}{x'}\,y' = k\,y'$. Donc $\vec{u}=k\vec{v}$.
      \item Si $x'=0$ alors $y'\neq 0$ (car $\vec{v}\neq\vec{0}$). L'égalité $x\,y' = y\,x'$ devient $x\,y' = 0$, donc $x=0$. Posons $k = \dfrac{y}{y'}$ ; on a $\vec{u}=k\vec{v}$.
      \end{itemize}
\end{itemize}
Dans tous les cas, il existe un réel $k$ tel que $\vec{u}=k\vec{v}$.
\end{experience}

\courssub{Cas particuliers et repères}

\begin{itemize}
\item \textbf{Vecteur nul :} $\vec{0}(0;0)$ est colinéaire à tout vecteur (déterminant toujours nul).
\item \textbf{Vecteurs des axes :} $\vec{\imath}(1;0)$ et $\vec{\jmath}(0;1)$ ne sont \textbf{pas} colinéaires car $1\cdot1 - 0\cdot0 = 1 \neq 0$.
\item \textbf{Vecteurs portés par le même axe :} $\vec{u}(a;0)$ et $\vec{v}(b;0)$ sont colinéaires ($a\cdot0 - 0\cdot b = 0$). De même pour $(0;a)$ et $(0;b)$.
\end{itemize}

\begin{popfigure}
\begin{tikzpicture}[scale=1.2,>=Stealth]
% Axes
\draw[->] (-0.5,0) -- (5,0) node[right] {$x$};
\draw[->] (0,-0.5) -- (0,4) node[above] {$y$};
% Vecteur v
\draw[popblue, thick, ->] (0,0) -- (2,1.5) node[midway, above left] {$\vec{v}(2;1,5)$};
% Vecteur u = 2 v
\draw[poporange, thick, ->] (0,0) -- (4,3) node[midway, above right] {$\vec{u}(4;3)$};
% Composantes de v
\draw[popblue, dashed] (2,0) -- (2,1.5);
\draw[popblue, dashed] (0,1.5) -- (2,1.5);
\node[popblue] at (2,-0.3) {$x'=2$};
\node[popblue] at (-0.3,1.5) {$y'=1,5$};
% Composantes de u
\draw[poporange, dashed] (4,0) -- (4,3);
\draw[poporange, dashed] (0,3) -- (4,3);
\node[poporange] at (4,-0.3) {$x=4$};
\node[poporange] at (-0.3,3) {$y=3$};
% Ratio
\node[align=center] at (5.5,2) {$\dfrac{x}{x'}=\dfrac{4}{2}=2$\\[2pt]$\dfrac{y}{y'}=\dfrac{3}{1,5}=2$};
\end{tikzpicture}
\legende{Deux vecteurs colinéaires $\vec{u}=2\vec{v}$ : les rapports des coordonnées sont égaux.}
\end{popfigure}

\courssub{Application géométrique : alignement de trois points}

\begin{propriete}[Alignement et colinéarité]
Trois points $A$, $B$, $C$ sont alignés \textbf{si et seulement si} les vecteurs $\overrightarrow{AB}$ et $\overrightarrow{AC}$ sont colinéaires, c'est-à-dire
\[
\det\bigl(\overrightarrow{AB},\overrightarrow{AC}\bigr)=0.
\]
\end{propriete}

\begin{popfigure}
\begin{tikzpicture}[scale=1.3,>=Stealth]
% Points
\coordinate (A) at (0,0);
\coordinate (B) at (2,1);
\coordinate (C) at (4,2);
% Droite
\draw[popmuted, thick] (-0.5,-0.25) -- (5,2.5);
% Points
\fill[popblue] (A) circle (2pt) node[below left] {$A$};
\fill[popblue] (B) circle (2pt) node[above left] {$B$};
\fill[popblue] (C) circle (2pt) node[above right] {$C$};
% Vecteurs
\draw[poporange, thick, ->] (A) -- (B) node[midway, above left] {$\overrightarrow{AB}$};
\draw[popgreen, thick, ->] (A) -- (C) node[midway, below right] {$\overrightarrow{AC}$};
% Composantes AB
\draw[poporange, dashed] (2,0) -- (2,1);
\draw[poporange, dashed] (0,1) -- (2,1);
% Composantes AC
\draw[popgreen, dashed] (4,0) -- (4,2);
\draw[popgreen, dashed] (0,2) -- (4,2);
% Légende déterminant
\node[align=center, popdark] at (5.5,1) {$\det(\overrightarrow{AB},\overrightarrow{AC}) = 2\times2 - 1\times4 = 0$};
\end{tikzpicture}
\legende{Trois points alignés : le déterminant des vecteurs $\overrightarrow{AB}$ et $\overrightarrow{AC}$ est nul (aire du parallélogramme = 0).}
\end{popfigure}

\courssub{Droites parallèles et coefficient directeur}

Deux droites sont parallèles \textbf{si et seulement si} leurs vecteurs directeurs sont colinéaires.  
Si une droite a pour équation cartésienne réduite $y = ax + b$, un vecteur directeur est $\vec{u}(1;a)$ (ou tout multiple). Le coefficient directeur $a$ est donc le rapport $\dfrac{y}{x}$ de tout vecteur directeur non vertical.

\begin{methode}[Prouver que trois points $A$, $B$, $C$ sont alignés]
\begin{enumerate}
\item Calculer les coordonnées des vecteurs $\overrightarrow{AB}(x_{AB};y_{AB})$ et $\overrightarrow{AC}(x_{AC};y_{AC})$.
\item Calculer le déterminant $D = x_{AB}\,y_{AC} - y_{AB}\,x_{AC}$.
\item \begin{itemize}
      \item Si $D = 0$ : les vecteurs sont colinéaires $\Rightarrow$ $A,B,C$ alignés.
      \item Si $D \neq 0$ : les vecteurs ne sont pas colinéaires $\Rightarrow$ $A,B,C$ non alignés (ils forment un triangle).
      \end{itemize}
\item Conclure par une phrase complète.
\end{enumerate}
\end{methode}

\courssub{Exemples résolus}

\begin{exoresolu}[Exemple 1 : tester la colinéarité]
Soient $\vec{u}(2;3)$ et $\vec{v}(4;6)$. Sont-ils colinéaires ?
\tcblower
Calculons le déterminant :
\[
\det(\vec{u},\vec{v}) = 2\times 6 - 3\times 4 = 12 - 12 = 0.
\]
Le déterminant est nul, donc $\vec{u}$ et $\vec{v}$ sont colinéaires. On vérifie d'ailleurs que $\vec{u} = \frac12\,\vec{v}$ (ou $\vec{v}=2\vec{u}$), le coefficient de proportionnalité est $k=2$.
\end{exoresolu}

\begin{exoresolu}[Exemple 2 : non-colinéarité]
Soient $\vec{u}(1;2)$ et $\vec{v}(2;5)$. Sont-ils colinéaires ?
\tcblower
\[
\det(\vec{u},\vec{v}) = 1\times 5 - 2\times 2 = 5 - 4 = 1 \neq 0.
\]
Le déterminant n'est pas nul, donc $\vec{u}$ et $\vec{v}$ ne sont \textbf{pas} colinéaires. Il n'existe aucun réel $k$ tel que $\vec{u}=k\vec{v}$.
\end{exoresolu}

\begin{exoresolu}[Exemple 3 : alignement de trois points]
Donnés $A(1;2)$, $B(3;5)$, $C(5;8)$. Les points $A,B,C$ sont-ils alignés ?
\tcblower
\begin{align*}
\overrightarrow{AB} &= (3-1;\;5-2) = (2;3),\\
\overrightarrow{AC} &= (5-1;\;8-2) = (4;6).
\end{align*}
Déterminant :
\[
\det(\overrightarrow{AB},\overrightarrow{AC}) = 2\times 6 - 3\times 4 = 12 - 12 = 0.
\]
Les vecteurs sont colinéaires, donc $A$, $B$, $C$ sont alignés. (On remarque d'ailleurs que $\overrightarrow{AC}=2\,\overrightarrow{AB}$.)
\end{exoresolu}

\begin{exoresolu}[Exemple 4 : trouver une coordonnée pour l'alignement]
Soient $A(2;1)$, $B(4;4)$ et $C(x;7)$. Déterminer $x$ pour que $A,B,C$ soient alignés.
\tcblower
$\overrightarrow{AB} = (2;3)$, $\overrightarrow{AC} = (x-2;\;6)$.  
Condition d'alignement : $\det(\overrightarrow{AB},\overrightarrow{AC}) = 0$.
\[
2\times 6 - 3\times(x-2) = 0 \;\Longrightarrow\; 12 - 3x + 6 = 0 \;\Longrightarrow\; 18 = 3x \;\Longrightarrow\; x = 6.
\]
Donc $C(6;7)$.
\end{exoresolu}

\courssub{Attention : pièges à éviter}

\begin{attention}[Division par zéro et produit en croix]
\begin{itemize}
\item \textbf{Ne jamais écrire} $k = \dfrac{x}{x'}$ ou $k = \dfrac{y}{y'}$ sans vérifier que le dénominateur est non nul. Si $x'=0$ ou $y'=0$, la division est interdite.
\item \textbf{Utilisez toujours le produit en croix} (le déterminant) : $x\,y' = y\,x'$. Cette écriture est valide même quand une coordonnée est nulle.
\item \textbf{Vecteur nul :} si l'un des deux vecteurs est $\vec{0}$, le déterminant est automatiquement nul et les vecteurs sont colinéaires (par définition). Ne cherchez pas de coefficient $k$ unique.
\item \textbf{Ordre des vecteurs :} $\det(\vec{u},\vec{v}) = -\det(\vec{v},\vec{u})$. Le signe change mais la nullité ne change pas.
\end{itemize}
\end{attention}

\courssub{Le savais-tu ? — Aire du parallélogramme}

\begin{savaistu}[Interprétation géométrique du déterminant]
Le déterminant $x\,y' - y\,x'$ de deux vecteurs $\vec{u}(x;y)$ et $\vec{v}(x';y')$ représente l'\textbf{aire algébrique} du parallélogramme construit sur ces deux vecteurs.  
\begin{itemize}
\item Si l'aire est nulle, le parallélogramme est « aplati » : les deux vecteurs sont sur la même droite $\Rightarrow$ colinéaires.
\item Le signe du déterminant indique l'orientation (directe ou indirecte) du couple $(\vec{u},\vec{v})$.
\end{itemize}
C'est pour cela que la condition de colinéarité s'exprime par « déterminant nul ».
\end{savaistu}

\courssub{Exercices d'entraînement}

\begin{exercice}[Exercice 1]
Déterminer si les couples de vecteurs suivants sont colinéaires :
\begin{enumerate}
\item $\vec{u}(3;-2)$ et $\vec{v}(-6;4)$
\item $\vec{u}(0;5)$ et $\vec{v}(0;-3)$
\item $\vec{u}(1;1)$ et $\vec{v}(2;3)$
\end{enumerate}
\end{exercice}

\begin{exercice}[Exercice 2]
Soient $A(-1;2)$, $B(2;5)$ et $C(5;8)$. Montrer que $A$, $B$, $C$ sont alignés.
\end{exercice}

\begin{exercice}[Exercice 3]
Trouver la valeur de $k$ pour que les vecteurs $\vec{u}(k;4)$ et $\vec{v}(3;6)$ soient colinéaires.
\end{exercice}

\begin{exercice}[Exercice 4]
Une droite $\mathcal{D}$ passe par $A(1;3)$ et a pour vecteur directeur $\vec{u}(2;-1)$. Une droite $\mathcal{D}'$ passe par $B(0;0)$ et a pour vecteur directeur $\vec{v}(-4;2)$. Les droites $\mathcal{D}$ et $\mathcal{D}'$ sont-elles parallèles ? Justifier.
\end{exercice}

\begin{corrige}[Exercice 1]
\begin{enumerate}
\item $\det = 3\times4 - (-2)\times(-6) = 12 - 12 = 0$ $\Rightarrow$ colinéaires ($k=-2$).
\item $\det = 0\times(-3) - 5\times0 = 0$ $\Rightarrow$ colinéaires (vecteurs verticaux).
\item $\det = 1\times3 - 1\times2 = 3 - 2 = 1 \neq 0$ $\Rightarrow$ non colinéaires.
\end{enumerate}
\end{corrige}

\begin{corrige}[Exercice 2]
$\overrightarrow{AB} = (3;3)$, $\overrightarrow{AC} = (6;6)$.  
$\det = 3\times6 - 3\times6 = 18 - 18 = 0$.  
Les vecteurs sont colinéaires, donc $A$, $B$, $C$ sont alignés. (On a $\overrightarrow{AC}=2\,\overrightarrow{AB}$.)
\end{corrige}

\begin{corrige}[Exercice 3]
Condition : $k\times6 - 4\times3 = 0 \;\Rightarrow\; 6k - 12 = 0 \;\Rightarrow\; k = 2$.
\end{corrige}

\begin{corrige}[Exercice 4]
$\det(\vec{u},\vec{v}) = 2\times2 - (-1)\times(-4) = 4 - 4 = 0$.  
Les vecteurs directeurs sont colinéaires, donc les droites sont parallèles (elles sont distinctes car $B(0;0)$ n'appartient pas à $\mathcal{D}$ : $0 \neq -1+3$).
\end{corrige}

\coursec{Condition d'orthogonalité de deux vecteurs (Produit scalaire)}

\noindent Dans la section précédente, nous avons vu comment détecter si deux vecteurs ont \emph{la même direction} (colinéarité) grâce au déterminant. Nous allons maintenant aborder le cas tout aussi fondamental : comment savoir si deux vecteurs sont \emph{perpendiculaires} ? La réponse réside dans un outil puissant de la géométrie analytique : le \cle{produit scalaire}. Ce n'est pas une simple formule à mémoriser, c'est la clé qui relie l'algèbre (les coordonnées) à la géométrie (les angles, les distances, les projections). Maîtriser le produit scalaire, c'est s'ouvrir les portes de la trigonométrie vectorielle et de la géométrie dans l'espace que vous découvrirez au lycée.

\courssub{Rappel : Produit scalaire géométrique et orthogonalité}

\noindent Vous connaissez déjà la notion d'angle entre deux vecteurs. Le produit scalaire permet de \textbf{mesurer} cet angle par un calcul algébrique.

\begin{definition}[Produit scalaire de deux vecteurs -- Définition géométrique]
Soient $\vec{u}$ et $\vec{v}$ deux vecteurs non nuls, et $(\vec{u}, \vec{v})$ l'angle orienté entre eux (pris dans $[0; \pi]$). Le \textbf{produit scalaire} de $\vec{u}$ par $\vec{v}$, noté $\vec{u} \cdot \vec{v}$ (ou parfois $\langle \vec{u}, \vec{v} \rangle$), est le nombre réel défini par :
\[
\vec{u} \cdot \vec{v} = \|\vec{u}\| \times \|\vec{v}\| \times \cos(\vec{u}, \vec{v})
\]
Par convention, si l'un des vecteurs est nul, le produit scalaire vaut $0$.
\end{definition}

\begin{aretenir}[Interprétation géométrique immédiate]
Le produit scalaire est le produit de la norme d'un vecteur par la \textbf{projection orthogonale} de l'autre sur sa direction. C'est l'``ombre'' d'un vecteur sur l'autre, multipliée par la longueur du support.
\end{aretenir}

\noindent Cette définition géométrique donne tout de suite la condition fondamentale d'orthogonalité.

\begin{propriete}[Caractérisation de l'orthogonalité par le produit scalaire]
Soient $\vec{u}$ et $\vec{v}$ deux vecteurs.
\[
\vec{u} \perp \vec{v} \quad \Longleftrightarrow \quad \vec{u} \cdot \vec{v} = 0
\]
\emph{(À condition que les deux vecteurs soient non nuls).}
\end{propriete}

\begin{experience}[Découverte : Pourquoi le produit scalaire nul signifie perpendicularité ?]
Prenons deux vecteurs $\vec{u}$ et $\vec{v}$.
\begin{enumerate}
    \item Si $\vec{u} \perp \vec{v}$, l'angle $(\vec{u}, \vec{v}) = 90^\circ$. Or $\cos 90^\circ = 0$. Donc $\vec{u} \cdot \vec{v} = \|\vec{u}\| \|\vec{v}\| \times 0 = 0$.
    \item Réciproquement, si $\vec{u} \cdot \vec{v} = 0$ et que $\vec{u} \neq \vec{0}, \vec{v} \neq \vec{0}$, alors $\|\vec{u}\| \|\vec{v}\| \neq 0$. Il faut donc que $\cos(\vec{u}, \vec{v}) = 0$, ce qui impose $(\vec{u}, \vec{v}) = 90^\circ$.
\end{enumerate}
C'est la traduction algébrique exacte de la perpendicularité.
\end{experience}

\courssub{Expression algébrique du produit scalaire dans un repère orthonormé}

\noindent La définition géométrique est belle, mais pour calculer concrètement avec des coordonnées, nous avons besoin d'une formule algébrique. C'est ici que le choix d'un \textbf{repère orthonormé} $(O; \vec{i}, \vec{j})$ devient crucial.

\begin{propriete}[Produit scalaire en coordonnées -- Théorème fondamental]
Dans un repère \textbf{orthonormé} $(O; \vec{i}, \vec{j})$, soient $\vec{u}(x; y)$ et $\vec{v}(x'; y')$ deux vecteurs. Leur produit scalaire s'écrit :
\[
\boxed{\vec{u} \cdot \vec{v} = x x' + y y'}
\]

\noindent \textbf{Démonstration (exigible au BEPC -- à savoir rédiger) :}
On développe $\vec{u}$ et $\vec{v}$ dans la base $(\vec{i}, \vec{j})$ :
$\vec{u} = x\vec{i} + y\vec{j}$ et $\vec{v} = x'\vec{i} + y'\vec{j}$.
Par bilinéarité du produit scalaire (distributivité par rapport à l'addition et factorisation des scalaires) :
\begin{align*}
\vec{u} \cdot \vec{v} &= (x\vec{i} + y\vec{j}) \cdot (x'\vec{i} + y'\vec{j}) \\
&= x x' (\vec{i} \cdot \vec{i}) + x y' (\vec{i} \cdot \vec{j}) + y x' (\vec{j} \cdot \vec{i}) + y y' (\vec{j} \cdot \vec{j})
\end{align*}
Dans un repère \textbf{orthonormé}, les vecteurs de base sont unitaires ($\|\vec{i}\| = \|\vec{j}\| = 1$) et orthogonaux ($\vec{i} \perp \vec{j}$).
Donc :
$\vec{i} \cdot \vec{i} = \|\vec{i}\|^2 = 1$,
$\vec{j} \cdot \vec{j} = \|\vec{j}\|^2 = 1$,
$\vec{i} \cdot \vec{j} = \vec{j} \cdot \vec{i} = 0$.
En substituant :
$\vec{u} \cdot \vec{v} = x x' (1) + x y' (0) + y x' (0) + y y' (1) = x x' + y y'$.
\end{propriete}

\begin{attention}[Piège mortel : Repère NON orthonormé]
Cette formule $x x' + y y'$ \textbf{ne fonctionne QUE dans un repère orthonormé}.
Si le repère n'est pas orthonormé (vecteurs de base non perpendiculaires ou non unitaires), les produits croisés $\vec{i} \cdot \vec{j}$ ne sont pas nuls et la formule est fausse. \textbf{Vérifiez toujours : ``Dans un repère orthonormé...''} avant d'appliquer $xx' + yy'$.
\end{attention}

\begin{popfigure}
\begin{tikzpicture}[scale=0.9, >=Stealth]
    % Axes
    \draw[->, thick] (-0.5,0) -- (5,0) node[right] {$x$};
    \draw[->, thick] (0,-0.5) -- (0,4) node[above] {$y$};
    \foreach \x in {1,...,4} \draw (\x, -0.1) -- (\x, 0.1) node[below, font=\tiny] {\x};
    \foreach \y in {1,...,3} \draw (-0.1, \y) -- (0.1, \y) node[left, font=\tiny] {\y};
    % Vecteur u = (3, 1)
    \draw[popblue, very thick, ->] (0,0) -- (3,1) node[midway, above left, popblue] {$\vec{u}(3;1)$};
    \draw[popblue, dashed, thin] (3,0) -- (3,1) node[below, font=\small, popblue] {$x=3$};
    \draw[popblue, dashed, thin] (0,1) -- (3,1) node[left, font=\small, popblue] {$y=1$};
    % Vecteur v = (-1, 3)
    \draw[poporange, very thick, ->] (0,0) -- (-1,3) node[midway, above left, poporange] {$\vec{v}(-1;3)$};
    \draw[poporange, dashed, thin] (-1,0) -- (-1,3) node[below, font=\small, poporange] {$x'=-1$};
    \draw[poporange, dashed, thin] (0,3) -- (-1,3) node[left, font=\small, poporange] {$y'=3$};
    % Angle droit
    \draw[popgreen, thick] (0.2, 0.06) -- (0.26, 0.26) -- (0.06, 0.32);
    % Formule
    \node[align=center, fill=popcream, draw=popline, rounded corners, inner sep=4pt] at (3.5, 2.5) {
        $\vec{u} \cdot \vec{v} = 3 \times (-1) + 1 \times 3 = -3 + 3 = 0$ \\
        $\Rightarrow \vec{u} \perp \vec{v}$
    };
\end{tikzpicture}
\legende{Visualisation du produit scalaire nul : $\vec{u}(3;1)$ et $\vec{v}(-1;3)$. La somme des produits des coordonnées ($xx' + yy'$) annule les composantes.}
\end{popfigure}

\courssub{Condition d'orthogonalité de deux vecteurs : Théorème et applications}

\noindent En combinant la propriété géométrique ($\vec{u} \perp \vec{v} \iff \vec{u}\cdot\vec{v}=0$) et la formule algébrique ($\vec{u}\cdot\vec{v} = xx' + yy'$), nous obtenons le critère de calcul ultime.

\begin{aretenir}[Théorème : Condition d'orthogonalité en coordonnées]
Dans un repère \textbf{orthonormé}, soient $\vec{u}(x; y)$ et $\vec{v}(x'; y')$ deux vecteurs \textbf{non nuls}.
\[
\vec{u} \perp \vec{v} \quad \Longleftrightarrow \quad x x' + y y' = 0
\]

\emph{Preuve :} $\vec{u} \perp \vec{v} \iff \vec{u}\cdot\vec{v}=0 \iff xx' + yy' = 0$. \qed
\end{aretenir}

\begin{exemplebox}[Vérification d'un angle droit]
Soient les points $A(1;2)$, $B(4;3)$ et $C(3;6)$ dans un repère orthonormé. Montrer que le triangle $ABC$ est rectangle en $B$.

\textbf{Résolution :}
\begin{enumerate}
    \item Calculons les vecteurs $\overrightarrow{BA}$ et $\overrightarrow{BC}$ :
    $\overrightarrow{BA} = A - B = (1-4; 2-3) = (-3; -1)$.
    $\overrightarrow{BC} = C - B = (3-4; 6-3) = (-1; 3)$.
    \item Calculons leur produit scalaire :
    $\overrightarrow{BA} \cdot \overrightarrow{BC} = (-3) \times (-1) + (-1) \times 3 = 3 - 3 = 0$.
    \item Comme les deux vecteurs sont non nuls et leur produit scalaire est nul, $\overrightarrow{BA} \perp \overrightarrow{BC}$.
    \item \textbf{Conclusion :} L'angle $\widehat{ABC}$ est un angle droit. Le triangle $ABC$ est rectangle en $B$.
\end{enumerate}
\end{exemplebox}

\begin{exemplebox}[Recherche d'une inconnue pour l'orthogonalité]
Déterminer $k$ pour que les vecteurs $\vec{u}(2; k)$ et $\vec{v}(-3; 4)$ soient orthogonaux.

\textbf{Résolution :}
$\vec{u} \perp \vec{v} \iff 2 \times (-3) + k \times 4 = 0 \iff -6 + 4k = 0 \iff 4k = 6 \iff k = 1,5$.
Pour $k = 1,5$, $\vec{u}(2; 1,5)$ et $\vec{v}(-3; 4)$ sont orthogonaux.
\end{exemplebox}

\courssub{Lien avec les pentes des droites : Produit des pentes = -1}

\noindent C'est une application majeure pour l'étude des droites. Dans un repère orthonormé, une droite $d$ non parallèle à l'axe des ordonnées a une \cle{pente} $m$. Un vecteur directeur de $d$ est $\vec{u}(1; m)$ (ou tout multiple).

\begin{propriete}[Condition de perpendicularité de deux droites par les pentes]
Soient $d$ et $d'$ deux droites non parallèles aux axes, de pentes respectives $m$ et $m'$.
\[
d \perp d' \quad \Longleftrightarrow \quad m \times m' = -1
\]

\noindent \textbf{Démonstration :}
Un vecteur directeur de $d$ est $\vec{u}(1; m)$. Un vecteur directeur de $d'$ est $\vec{v}(1; m')$.
$d \perp d' \iff \vec{u} \perp \vec{v} \iff 1 \times 1 + m \times m' = 0 \iff m m' = -1$.
\end{propriete}

\begin{attention}[Cas des droites parallèles aux axes]
La formule $m m' = -1$ ne s'applique \textbf{pas} si l'une des droites est verticale (pente infinie) ou horizontale (pente 0).
\begin{itemize}
    \item Une droite horizontale ($m=0$) est perpendiculaire à une droite verticale (pas de pente).
    \item Vérifiez toujours l'orthogonalité par les \textbf{vecteurs directeurs} (produit scalaire nul), cela marche dans \textbf{tous les cas} sans exception.
\end{itemize}
\end{attention}

\begin{exemplebox}[Droites perpendiculaires et pentes]
La droite $d$ a pour équation $y = \frac{1}{3}x + 2$. La droite $d'$ a pour équation $y = -3x - 1$.
Pente de $d$ : $m = \frac{1}{3}$. Pente de $d'$ : $m' = -3$.
$m \times m' = \frac{1}{3} \times (-3) = -1$.
Donc $d \perp d'$.
(Rappel : vecteurs directeurs $\vec{u}(3;1)$ et $\vec{v}(1;-3)$ ou $(-1;3)$ : $3\times(-1) + 1\times 3 = 0$).
\end{exemplebox}

\courssub{Équation cartésienne d'une droite et vecteur normal}

\noindent Le produit scalaire permet de passer instantanément d'un vecteur directeur à une équation cartésienne (réduite ou générale).

\begin{propriete}[Vecteur normal et équation cartésienne]
Soit $d$ une droite de vecteur directeur $\vec{u}(a; b)$. Un \textbf{vecteur normal} à $d$ est un vecteur $\vec{n}$ orthogonal à $\vec{u}$.
On peut choisir $\vec{n}(-b; a)$ ou $\vec{n}(b; -a)$ car $a(-b) + b(a) = 0$.
L'équation cartésienne de $d$ s'écrit alors : $\vec{n} \cdot \overrightarrow{M_0M} = 0$ pour tout point $M(x;y)$ de la droite, $M_0(x_0;y_0)$ étant un point connu de $d$.
Soit : $-b(x - x_0) + a(y - y_0) = 0 \iff ax - by + c = 0$ (forme générale).
\emph{Plus généralement, une équation de la forme $ax + by + c = 0$ a pour vecteur normal $\vec{n}(a; b)$.}
\end{propriete}

\begin{methode}[Prouver que deux droites sont perpendiculaires]
\begin{enumerate}
    \item Identifier un vecteur directeur $\vec{u}$ pour la première droite et $\vec{v}$ pour la seconde.
    \item Calculer le produit scalaire $\vec{u} \cdot \vec{v} = x_u x_v + y_u y_v$.
    \item Si le résultat est $0$ (et vecteurs non nuls), conclure : ``Les vecteurs directeurs sont orthogonaux, donc les droites sont perpendiculaires.''
    \item Si les droites sont données par des équations cartésiennes $ax+by+c=0$ et $a'x+b'y+c'=0$, leurs vecteurs normaux sont $\vec{n}(a;b)$ et $\vec{n'}(a';b')$. Les droites sont perpendiculaires $\iff \vec{n} \perp \vec{n'} \iff aa' + bb' = 0$.
\end{enumerate}
\end{methode}

\courssub{Problème concret : Projection orthogonale d'un point sur une droite}

\noindent C'est un classique du BEPC et une application synthétique : colinéarité + orthogonalité.

\begin{exoresolu}[Projection orthogonale -- Résolution par système]
Dans un repère orthonormé, on donne la droite $d$ d'équation $y = 2x - 1$ et le point $M(3; 4)$.
Déterminer les coordonnées du point $H$, projection orthogonale de $M$ sur $d$.

\textbf{Analyse :}
$H$ appartient à $d$ $\implies$ $\overrightarrow{OH}$ vérifie l'équation de $d$ (Colinéarité avec un vecteur directeur).
$MH \perp d$ $\implies$ $\overrightarrow{MH} \cdot \vec{u} = 0$ où $\vec{u}$ est un vecteur directeur de $d$ (Orthogonalité).

\textbf{Résolution :}
\begin{enumerate}
    \item \textbf{Vecteur directeur de $d$ :} Pente $m=2 \Rightarrow \vec{u}(1; 2)$.
    \item \textbf{Inconnues :} $H(x_H; y_H)$.
    \item \textbf{Condition 1 : $H \in d$ (Colinéarité).}
    $y_H = 2x_H - 1$. \hfill (Équation 1)
    \item \textbf{Condition 2 : $MH \perp d$ (Orthogonalité).}
    $\overrightarrow{MH} = (x_H - 3; y_H - 4)$.
    $\overrightarrow{MH} \cdot \vec{u} = 0 \iff (x_H - 3) \times 1 + (y_H - 4) \times 2 = 0$.
    $\iff x_H - 3 + 2y_H - 8 = 0 \iff x_H + 2y_H - 11 = 0$. \hfill (Équation 2)
    \item \textbf{Résolution du système :}
    \[\begin{cases}
    y_H = 2x_H - 1 \\
    x_H + 2y_H = 11
    \end{cases}\]
    $\implies x_H + 2(2x_H - 1) = 11 \iff x_H + 4x_H - 2 = 11 \iff 5x_H = 13 \iff x_H = 2{,}6$.
    $y_H = 2(2,6) - 1 = 5,2 - 1 = 4,2$.
    \item \textbf{Conclusion :} $H(2,6; 4,2)$.
\end{enumerate}

\textbf{Vérification rapide :} $\overrightarrow{MH} = (-0,4; 0,2)$. $\overrightarrow{MH} \cdot \vec{u} = -0,4 \times 1 + 0,2 \times 2 = -0,4 + 0,4 = 0$. \checkmark
\tcblower
\textbf{Autre méthode (vecteur normal) :}
Équation cartésienne de $d$ : $2x - y - 1 = 0$. Vecteur normal $\vec{n}(2; -1)$.
Droite $(MH)$ : passe par $M(3;4)$, directeur $\vec{n}(2;-1)$.
Paramétrique : $x = 3 + 2t$, $y = 4 - t$.
Intersection avec $d$ : $2(3+2t) - (4-t) - 1 = 0 \iff 6+4t-4+t-1=0 \iff 5t+1=0 \iff t = -0,2$.
$H(3+2(-0,2); 4-(-0,2)) = H(2,6; 4,2)$. Même résultat, souvent plus rapide.
\end{exoresolu}

\begin{popfigure}
\begin{tikzpicture}[scale=0.8, >=Stealth]
    % Axes
    \draw[->, thin, gray] (-1,0) -- (5,0) node[right] {$x$};
    \draw[->, thin, gray] (0,-1) -- (0,6) node[above] {$y$};
    \foreach \x in {-1,...,4} \draw (\x, -0.05) -- (\x, 0.05) node[below, font=\tiny] {\x};
    \foreach \y in {-1,...,5} \draw (-0.05, \y) -- (0.05, \y) node[left, font=\tiny] {\y};
    % Droite d : y = 2x - 1
    \draw[popblue, thick, domain=-0.5:3] plot (\x, {2*\x - 1}) node[right, popblue] {$d: y=2x-1$};
    % Point M
    \fill[poporange] (3,4) circle (2.5pt) node[above right] {$M(3;4)$};
    % Point H
    \fill[popgreen] (2.6, 4.2) circle (2.5pt) node[above left] {$H(2,6;4,2)$};
    % Segment MH
    \draw[popgreen, thick, dashed] (3,4) -- (2.6, 4.2) node[midway, right] {$\overrightarrow{MH}$};
    % Vecteur directeur u(1,2) à partir de H
    \draw[popblue, very thick, ->] (2.6, 4.2) -- (3.6, 6.2) node[above right] {$\vec{u}(1;2)$};
    % Angle droit
    \draw[popgreen, thick] (2.8, 4.5) -- (2.9, 4.3) -- (3.1, 4.4);
    % Vecteur normal n(2,-1) à partir de M
    \draw[poppurple, very thick, ->] (3,4) -- (5,3) node[below right] {$\vec{n}(2;-1)$};
\end{tikzpicture}
\legende{Projection orthogonale $H$ de $M$ sur $d$. $\overrightarrow{MH}$ est orthogonal à $\vec{u}$ (directeur de $d$) et colinéaire à $\vec{n}$ (normal de $d$).}
\end{popfigure}

\courssub{Exercices d'entraînement progressif}

\begin{exercice}[Vérification d'orthogonalité]
Dans un repère orthonormé, on donne $\vec{u}(-2; 5)$ et $\vec{v}(10; 4)$.
\begin{enumerate}
    \item Calculer $\vec{u} \cdot \vec{v}$.
    \item En déduire la nature de l'angle $(\vec{u}, \vec{v})$.
    \item Trouver une valeur de $k$ telle que $\vec{u}$ et $\vec{w}(k; 2)$ soient orthogonaux.
\end{enumerate}
\end{exercice}

\begin{exercice}[Droite et perpendicularité]
La droite $d$ passe par $A(1; -2)$ et $B(3; 2)$. La droite $d'$ a pour équation $x + 2y - 5 = 0$.
\begin{enumerate}
    \item Déterminer un vecteur directeur de $d$ et un vecteur normal de $d'$.
    \item Montrer que $d \perp d'$.
\end{enumerate}
\end{exercice}

\begin{exercice}[Projection orthogonale -- Situation concrète]
Un pylône de télécommunication \textbf{MTN} est implanté au point $P(2; 3)$ sur une carte (repère orthonormé, unité : 100 m). Une route suit la droite $d$ d'équation $y = -\frac{1}{2}x + 4$. On veut raccorder le pylône à la route par le chemin le plus court (perpendiculaire à la route).
\begin{enumerate}
    \item Justifier que le point de raccordement $R$ est la projection orthogonale de $P$ sur $d$.
    \item Déterminer les coordonnées de $R$.
    \item Calculer la longueur du câble nécessaire (distance $PR$) en mètres.
\end{enumerate}
\end{exercice}

\begin{corrige}[Exercice 1]
\begin{enumerate}
    \item $\vec{u} \cdot \vec{v} = (-2)(10) + (5)(4) = -20 + 20 = 0$.
    \item Le produit scalaire est nul, les vecteurs sont non nuls $\Rightarrow \vec{u} \perp \vec{v}$. L'angle est droit ($90^\circ$).
    \item $\vec{u} \perp \vec{w} \iff -2k + 5(2) = 0 \iff -2k + 10 = 0 \iff k = 5$.
\end{enumerate}
\end{corrige}

\begin{corrige}[Exercice 2]
\begin{enumerate}
    \item $\overrightarrow{AB} = (2; 4)$ ou $\vec{u}(1; 2)$. Équation $d' : x + 2y - 5 = 0 \Rightarrow \vec{n}(1; 2)$.
    \item Le vecteur directeur de $d$ est $\vec{u}(1;2)$. Le vecteur normal de $d'$ est $\vec{n}(1;2)$.
    $\vec{u}$ et $\vec{n}$ sont colinéaires (identiques). Donc $\vec{u}$ est orthogonal à tout vecteur directeur de $d'$.
    \textbf{Conclusion :} $d \perp d'$.
    \emph{(Vérification : vecteur directeur de $d'$ : $\vec{v}(-2; 1)$ car $1(-2)+2(1)=0$. $\vec{u} \cdot \vec{v} = 1(-2) + 2(1) = 0$.)}
\end{enumerate}
\end{corrige}

\begin{corrige}[Exercice 3]
\begin{enumerate}
    \item Le plus court chemin d'un point à une droite est la perpendiculaire. Donc $PR \perp d$, $R \in d \Rightarrow R$ est la projection orthogonale.
    \item $d : y = -0,5x + 4$. $\vec{u}(2; -1)$ (ou $(1; -0,5)$). $P(2;3)$. $R(x_R; y_R)$.
    Système : $\begin{cases} y_R = -0,5 x_R + 4 \\ \overrightarrow{PR} \cdot \vec{u} = 0 \end{cases}$
    $\overrightarrow{PR} = (x_R-2; y_R-3)$.
    $(x_R-2)(2) + (y_R-3)(-1) = 0 \iff 2x_R - 4 - y_R + 3 = 0 \iff 2x_R - y_R - 1 = 0$.
    Substitution : $2x_R - (-0,5x_R + 4) - 1 = 0 \iff 2,5x_R - 5 = 0 \iff x_R = 2$.
    $y_R = -0,5(2) + 4 = 3$.
    $R(2; 3)$.
    \item $R = P$ ! Le pylône est \textbf{sur} la route. Distance $PR = 0$ m. (C'était un piège : vérifier si le point appartient à la droite au départ ! $3 = -0,5(2)+4 = 3$, oui).
\end{enumerate}
\end{corrige}

\begin{savaistu}[Le produit scalaire : l'outil universel de la géométrie]
Le produit scalaire n'est pas qu'une formule de 3ème. C'est la \textbf{colonne vertébrale} de la géométrie analytique supérieure.
\begin{itemize}
    \item \textbf{Angles quelconques :} $\cos \theta = \frac{\vec{u}\cdot\vec{v}}{\|\vec{u}\|\|\vec{v}\|}$. Permet de calculer n'importe quel angle, pas seulement le droit.
    \item \textbf{Projection orthogonale :} La projection de $\vec{u}$ sur $\vec{v}$ est le vecteur $\frac{\vec{u}\cdot\vec{v}}{\|\vec{v}\|^2}\vec{v}$. C'est l'``ombre'' de $\vec{u}$ sur $\vec{v}$.
    \item \textbf{Applications réelles (Cameroun et monde) :}
    \begin{itemize}
        \item \textbf{GPS / Navigation :} Calcul des distances orthodromiques (sur la sphère terrestre) via produits scalaires de vecteurs 3D.
        \item \textbf{Infographie / Jeux vidéo :} Détection des collisions, éclairage (produit scalaire normale/lumière), ombres portées.
        \item \textbf{Physique :} Travail d'une force $W = \vec{F} \cdot \vec{d}$. Si force $\perp$ déplacement, travail nul (ex: force centripète).
        \item \textbf{Télécoms (MTN/Orange) :} Alignement d'antennes, calcul de zones de couverture (faisceaux), interférences.
    \end{itemize}
\end{itemize}
Maîtriser $xx' + yy' = 0$ aujourd'hui, c'est poser la première pierre pour comprendre le monde numérique et physique de demain.
\end{savaistu}

\coursec{Coordonnées d'un vecteur dans un repère orthonormé}

\courssub{Repère orthonormé et coordonnées d'un point}

\begin{definition}[Repère orthonormé]
Un \cle{repère orthonormé} du plan est un triplet $(O\,;\vec{\imath},\vec{\jmath})$ où :
\begin{itemize}
\item $O$ est l'origine,
\item $\vec{\imath}$ et $\vec{\jmath}$ sont deux vecteurs unitaires ($||\vec{\imath}||=||\vec{\jmath}||=1$) et orthogonaux ($\vec{\imath}\perp\vec{\jmath}$).
\end{itemize}
Les axes $(O\vec{\imath})$ et $(O\vec{\jmath})$ sont donc deux droites perpendiculaires, graduées avec la \textbf{même unité de longueur}.
\end{definition}

\begin{propriete}[Coordonnées d'un point]
Dans un repère orthonormé $(O\,;\vec{\imath},\vec{\jmath})$, tout point $M$ du plan est repéré par un couple unique $(x\,;y)$ de nombres réels tels que :
\[ \overrightarrow{OM} = x\,\vec{\imath} + y\,\vec{\jmath}. \]
$x$ est l'\cle{abscisse} de $M$, $y$ son \cle{ordonnée}. On note $M(x\,;y)$.
\end{propriete}

\begin{attention}[Ordre des coordonnées]
On écrit toujours \textbf{abscisse ; ordonnée} (séparées par un point-virgule en France et au Cameroun). $M(3\,;-2)$ signifie $x=3$ et $y=-2$. Ne confonds pas avec la notation vectorielle $\vec{u}(3\,;-2)$ qui utilise les mêmes symboles mais désigne un vecteur, pas un point.
\end{attention}

\courssub{Coordonnées d'un vecteur $\overrightarrow{AB}$}

\begin{propriete}[Coordonnées du vecteur $\overrightarrow{AB}$]
Soient $A(x_A\,;y_A)$ et $B(x_B\,;y_B)$ deux points dans un repère orthonormé. Les coordonnées du vecteur $\overrightarrow{AB}$ sont :
\[ \overrightarrow{AB}\bigl(x_B - x_A \ ; \ y_B - y_A\bigr). \]
\emph{Moyen mnémotechnique :} « Coordonnées de l'arrivée moins coordonnées du départ ».
\end{propriete}

\begin{exemplebox}[Calcul de coordonnées]
Dans un repère orthonormé, $A(2\,;-1)$ et $B(-3\,;4)$.
\[ \overrightarrow{AB}(-3-2 \ ; \ 4-(-1)) = (-5 \ ; \ 5). \]
Vérification graphique : on part de $A$, on va 5 cases à gauche, 5 cases en haut, on arrive bien en $B$.
\end{exemplebox}

\courssub{Relation de Chasles et Milieu d'un segment}

\begin{propriete}[Relation de Chasles (version coordonnées)]
Pour trois points $A$, $B$, $C$ :
\[ \overrightarrow{AB} + \overrightarrow{BC} = \overrightarrow{AC}. \]
En coordonnées : $(x_B-x_A+x_C-x_B \ ; \ y_B-y_A+y_C-y_A) = (x_C-x_A \ ; \ y_C-y_A)$.
\end{propriete}

\begin{propriete}[Coordonnées du milieu]
Le milieu $I$ du segment $[AB]$ a pour coordonnées :
\[ I\left(\frac{x_A+x_B}{2} \ ; \ \frac{y_A+y_B}{2}\right). \]
C'est la moyenne des abscisses et la moyenne des ordonnées.
\end{propriete}

\begin{methode}[Calculer les coordonnées d'un vecteur ou d'un milieu]
\begin{enumerate}
\item Identifier les coordonnées des points donnés : $A(x_A\,;y_A)$, $B(x_B\,;y_B)$.
\item Appliquer la formule : $\overrightarrow{AB}(x_B-x_A \ ; \ y_B-y_A)$.
\item Pour un milieu $I$ de $[AB]$ : $I\bigl(\frac{x_A+x_B}{2} \ ; \ \frac{y_A+y_B}{2}\bigr)$.
\item Vérifier le signe des différences (sens du vecteur).
\end{enumerate}
\end{methode}

\begin{exoresolu}[Application : Chasles et milieu]
Donnés $A(1\,;3)$, $B(4\,;-1)$, $C(-2\,;5)$.
\begin{enumerate}
\item Calculer $\overrightarrow{AB}$, $\overrightarrow{BC}$, $\overrightarrow{AC}$ et vérifier Chasles.
\item Déterminer les coordonnées du milieu $I$ de $[AC]$.
\end{enumerate}
\tcblower
\textbf{1.}
$\overrightarrow{AB}(4-1 \ ; \ -1-3) = (3 \ ; \ -4)$.
$\overrightarrow{BC}(-2-4 \ ; \ 5-(-1)) = (-6 \ ; \ 6)$.
$\overrightarrow{AC}(-2-1 \ ; \ 5-3) = (-3 \ ; \ 2)$.
Somme : $(3+(-6) \ ; \ -4+6) = (-3 \ ; \ 2) = \overrightarrow{AC}$. \checkmark

\textbf{2.}
$I\left(\frac{1+(-2)}{2} \ ; \ \frac{3+5}{2}\right) = \left(-\frac{1}{2} \ ; \ 4\right)$.
\end{exoresolu}

\begin{exercice}[Entraînement : coordonnées et milieu]
Dans un repère orthonormé, $A(-2\,;5)$, $B(4\,;-3)$, $C(1\,;2)$.
\begin{enumerate}
\item Calculer $\overrightarrow{AB}$, $\overrightarrow{BA}$, $\overrightarrow{AC}$.
\item Vérifier que $\overrightarrow{AB} + \overrightarrow{BC} = \overrightarrow{AC}$.
\item Déterminer les coordonnées du milieu $I$ de $[AB]$ et du milieu $J$ de $[BC]$.
\item Que représente le vecteur $\overrightarrow{IJ}$ par rapport à $\overrightarrow{AC}$ ? (Conjecture)
\end{enumerate}
\end{exercice}

\begin{corrige}[Exercice : coordonnées et milieu]
\textbf{1.} $\overrightarrow{AB}(6\,;-8)$, $\overrightarrow{BA}(-6\,;8)$, $\overrightarrow{AC}(3\,;-3)$.
\textbf{2.} $\overrightarrow{BC}(-3\,;5)$. $\overrightarrow{AB}+\overrightarrow{BC} = (6-3\,;-8+5) = (3\,;-3) = \overrightarrow{AC}$.
\textbf{3.} $I(1\,;1)$, $J(2,5\,;-0,5)$.
\textbf{4.} $\overrightarrow{IJ}(1,5\,;-1,5) = \frac{1}{2}\overrightarrow{AC}$. C'est la propriété de la droite des milieux (programme de 4ème, rappel utile).
\end{corrige}

\coursec{Norme d'un vecteur et Distance entre deux points}

\courssub{Norme d'un vecteur}

\begin{definition}[Norme (longueur) d'un vecteur]
Dans un repère \textbf{orthonormé}, la \cle{norme} (ou longueur) d'un vecteur $\vec{u}(x\,;y)$ est le nombre positif noté $||\vec{u}||$ et défini par :
\[ ||\vec{u}|| = \sqrt{x^2 + y^2}. \]
C'est une application directe du théorème de Pythagore dans le triangle rectangle dont $\vec{u}$ est l'hypoténuse.
\end{definition}

\begin{attention}[Hypothèse cruciale]
La formule $||\vec{u}|| = \sqrt{x^2+y^2}$ n'est valable \textbf{que dans un repère orthonormé}. Dans un repère quelconque, elle est fausse. Vérifie toujours l'énoncé : « repère orthonormé » ou « axes perpendiculaires et même unité ».
\end{attention}

\courssub{Distance entre deux points}

\begin{propriete}[Distance $AB$]
La distance entre deux points $A(x_A\,;y_A)$ et $B(x_B\,;y_B)$ dans un repère orthonormé est la norme du vecteur $\overrightarrow{AB}$ :
\[ AB = ||\overrightarrow{AB}|| = \sqrt{(x_B-x_A)^2 + (y_B-y_A)^2}. \]
\end{propriete}

\begin{exemplebox}[Calcul de distances]
$A(1\,;-2)$, $B(4\,;2)$.
$\overrightarrow{AB}(3\,;4)$. $AB = \sqrt{3^2+4^2} = \sqrt{25} = 5$.
$C(-1\,;3)$, $D(2\,;-1)$.
$\overrightarrow{CD}(3\,;-4)$. $CD = \sqrt{3^2+(-4)^2} = 5$.
Les segments $[AB]$ et $[CD]$ ont même longueur.
\end{exemplebox}

\courssub{Propriétés de la distance}

\begin{aretenir}[Propriétés fondamentales]
Pour tous points $A, B, C$ dans un repère orthonormé :
\begin{itemize}
\item $AB \ge 0$ (positivité) ; $AB = 0 \iff A = B$.
\item $AB = BA$ (symétrie).
\item $AC \le AB + BC$ (inégalité triangulaire), avec égalité si et seulement si $A, B, C$ sont alignés \textbf{et} $B$ est entre $A$ et $C$.
\end{itemize}
\end{aretenir}

\begin{methode}[Calculer une distance ou un périmètre]
\begin{enumerate}
\item Noter les coordonnées des points.
\item Calculer les différences d'abscisses et d'ordonnées pour chaque segment.
\item Appliquer $d = \sqrt{(\Delta x)^2 + (\Delta y)^2}$.
\item Simplifier la racine si possible ($\sqrt{20}=2\sqrt{5}$).
\item Pour un périmètre : sommer les longueurs des côtés.
\end{enumerate}
\end{methode}

\begin{exoresolu}[Périmètre d'un triangle]
Triangle $MNP$ avec $M(0\,;0)$, $N(3\,;0)$, $P(3\,;4)$.
\tcblower
$MN = \sqrt{(3-0)^2+(0-0)^2} = 3$.
$NP = \sqrt{(3-3)^2+(4-0)^2} = 4$.
$MP = \sqrt{(3-0)^2+(4-0)^2} = 5$.
Périmètre $\mathcal{P} = 3+4+5 = 12$.
Remarque : $3^2+4^2=5^2$, le triangle est rectangle en $N$ (converse de Pythagore).
\end{exoresolu}

\begin{popfigure}
\begin{center}
\begin{tikzpicture}[scale=0.8]
\draw[popline, dashed, step=1] (-0.5,-0.5) grid (4.5,5.5);
\draw[-{Stealth}, thick] (-0.5,0) -- (4.5,0) node[below]{$x$};
\draw[-{Stealth}, thick] (0,-0.5) -- (0,5.5) node[left]{$y$};
\coordinate (M) at (0,0); \coordinate (N) at (3,0); \coordinate (P) at (3,4);
\draw[very thick, popblue] (M) -- (N) -- (P) -- cycle;
\draw[popdark] (2.8,0) -- (2.8,0.2) -- (3,0.2); % angle droit
\foreach \p/\n/\pos in {M/M/below left, N/N/below right, P/P/above right}
  \fill[popink] (\p) circle (2pt) node[\pos, font=\small]{\n};
\node[popblue, font=\small] at (1.5,-0.4) {3};
\node[popblue, font=\small] at (3.4,2) {4};
\node[popblue, font=\small] at (1.2,2.2) {5};
\end{tikzpicture}
\end{center}
\legende{Triangle rectangle $MNP$ : côtés 3, 4, 5.}
\end{popfigure}

\begin{exercice}[Entraînement : distances]
Dans un repère orthonormé, $E(-1\,;2)$, $F(3\,;5)$, $G(6\,;-1)$.
\begin{enumerate}
\item Calculer les longueurs $EF$, $FG$, $EG$.
\item Le triangle $EFG$ est-il rectangle ? Justifier.
\item Calculer son périmètre (valeur exacte puis approchée au dixième).
\end{enumerate}
\end{exercice}

\begin{corrige}[Exercice : distances]
\textbf{1.} $\overrightarrow{EF}(4\,;3) \Rightarrow EF=5$. $\overrightarrow{FG}(3\,;-6) \Rightarrow FG=\sqrt{45}=3\sqrt{5}$. $\overrightarrow{EG}(7\,;-3) \Rightarrow EG=\sqrt{58}$.
\textbf{2.} $EF^2+FG^2 = 25+45=70 \neq EG^2=58$. $EF^2+EG^2 \neq FG^2$. $FG^2+EG^2 \neq EF^2$. Non, pas rectangle.
\textbf{3.} $\mathcal{P} = 5 + 3\sqrt{5} + \sqrt{58} \approx 5 + 6,7 + 7,6 = 19,3$.
\end{corrige}

\coursec{Condition de colinéarité de deux vecteurs}

\courssub{Définition et caractérisation algébrique}

\begin{definition}[Vecteurs colinéaires]
Deux vecteurs $\vec{u}$ et $\vec{v}$ sont \cle{colinéaires} s'ils ont la même direction (ils sont portés par des droites parallèles ou confondues). On note $\vec{u} \parallel \vec{v}$.
Le vecteur nul $\vec{0}$ est colinéaire à tout vecteur.
\end{definition}

\begin{propriete}[Condition de colinéarité par les coordonnées]
Dans un repère (orthonormé ou non), deux vecteurs $\vec{u}(x\,;y)$ et $\vec{v}(x'\,;y')$ sont colinéaires \textbf{si et seulement si} :
\[ x\,y' - x'\,y = 0. \]
Cette quantité s'appelle le \cle{déterminant} de $\vec{u}$ et $\vec{v}$, noté $\det(\vec{u},\vec{v})$.
\end{propriete}

\begin{experience}[Découverte de la formule du déterminant]
Si $\vec{u}$ et $\vec{v}$ sont colinéaires et non nuls, il existe $k$ tel que $\vec{u}=k\vec{v}$, c'est-à-dire $x=kx'$ et $y=ky'$.
Alors $xy' - x'y = (kx')y' - x'(ky') = kx'y' - kx'y' = 0$.
Réciproquement, si $xy'-x'y=0$, alors $xy'=x'y$.
\begin{itemize}
\item Si $x'\neq 0$ et $y'\neq 0$, on a $\frac{x}{x'}=\frac{y}{y'}=k$, donc $\vec{u}=k\vec{v}$.
\item Si $x'=0$, alors $xy'=0$. Comme $\vec{v}\neq\vec{0}$, $y'\neq 0$, donc $x=0$. Les deux vecteurs sont verticaux, donc colinéaires.
\item Cas $y'=0$ symétrique.
\end{itemize}
\end{experience}

\begin{attention}[Pièges du déterminant]
\begin{itemize}
\item \textbf{Ordre :} $\det(\vec{u},\vec{v}) = -\det(\vec{v},\vec{u})$. Le signe change, mais la nullité reste la même.
\item \textbf{Vecteur nul :} $\det(\vec{0},\vec{v})=0$ toujours vrai. $\vec{0}$ est colinéaire à tout vecteur.
\item \textbf{Repère :} Cette formule fonctionne dans \textbf{tout} repère (pas besoin d'orthonormé), car la colinéarité est une notion affine.
\end{itemize}
\end{attention}

\courssub{Applications : alignement et parallélisme}

\begin{propriete}[Alignement de trois points]
Trois points $A, B, C$ sont alignés \textbf{si et seulement si} les vecteurs $\overrightarrow{AB}$ et $\overrightarrow{AC}$ sont colinéaires, c'est-à-dire :
\[ \det(\overrightarrow{AB},\overrightarrow{AC}) = 0. \]
\end{propriete}

\begin{propriete}[Droites parallèles]
Deux droites $(AB)$ et $(CD)$ sont parallèles \textbf{si et seulement si} $\overrightarrow{AB}$ et $\overrightarrow{CD}$ sont colinéaires.
\end{propriete}

\begin{methode}[Montrer que des points sont alignés ou des droites parallèles]
\begin{enumerate}
\item Calculer les coordonnées des deux vecteurs directeurs (ex: $\overrightarrow{AB}$ et $\overrightarrow{AC}$).
\item Calculer leur déterminant $x_1y_2 - x_2y_1$.
\item Si le résultat est $0$, conclure : « Les vecteurs sont colinéaires, donc les points sont alignés / les droites sont parallèles. »
\item Si le résultat n'est pas $0$, conclure : « Les vecteurs ne sont pas colinéaires, donc les points ne sont pas alignés / les droites ne sont pas parallèles. »
\end{enumerate}
\end{methode}

\begin{exoresolu}[Alignement et parallélisme]
$A(1\,;2)$, $B(3\,;5)$, $C(5\,;8)$, $D(0\,;0)$, $E(2\,;3)$.
\begin{enumerate}
\item Montrer que $A, B, C$ sont alignés.
\item Les droites $(AB)$ et $(DE)$ sont-elles parallèles ?
\end{enumerate}
\tcblower
\textbf{1.} $\overrightarrow{AB}(2\,;3)$, $\overrightarrow{AC}(4\,;6)$.
$\det = 2\times 6 - 4\times 3 = 12 - 12 = 0$.
$\overrightarrow{AB} \parallel \overrightarrow{AC} \implies A, B, C$ alignés.

\textbf{2.} $\overrightarrow{DE}(2\,;3)$. $\overrightarrow{AB}(2\,;3)$.
$\det(\overrightarrow{AB},\overrightarrow{DE}) = 2\times 3 - 2\times 3 = 0$.
Les vecteurs sont colinéaires, donc $(AB) \parallel (DE)$. (Ici même direction, droites distinctes car $A\notin (DE)$).
\end{exoresolu}

\begin{exercice}[Entraînement : colinéarité]
Dans un repère, $P(-2\,;1)$, $Q(2\,;3)$, $R(6\,;5)$, $S(1\,;-1)$, $T(3\,;2)$.
\begin{enumerate}
\item Montrer que $P, Q, R$ sont alignés.
\item Déterminer $a$ pour que les vecteurs $\vec{u}(2\,;a)$ et $\vec{v}(4\,;6)$ soient colinéaires.
\item Les droites $(PQ)$ et $(ST)$ sont-elles parallèles ?
\end{enumerate}
\end{exercice}

\begin{corrige}[Exercice : colinéarité]
\textbf{1.} $\overrightarrow{PQ}(4\,;2)$, $\overrightarrow{PR}(8\,;4)$. $\det = 4\times 4 - 8\times 2 = 0$. Alignés.
\textbf{2.} $2\times 6 - 4\times a = 0 \iff 12 - 4a = 0 \iff a = 3$.
\textbf{3.} $\overrightarrow{ST}(2\,;3)$. $\overrightarrow{PQ}(4\,;2)$. $\det = 4\times 3 - 2\times 2 = 12 - 4 = 8 \neq 0$. Non parallèles.
\end{corrige}

\coursec{Condition d'orthogonalité de deux vecteurs (Produit scalaire)}

\courssub{Produit scalaire et orthogonalité}

\begin{definition}[Produit scalaire]
Dans un repère \textbf{orthonormé}, le \cle{produit scalaire} de deux vecteurs $\vec{u}(x\,;y)$ et $\vec{v}(x'\,;y')$ est le nombre réel noté $\vec{u}\cdot\vec{v}$ (ou $\langle\vec{u},\vec{v}\rangle$) défini par :
\[ \vec{u}\cdot\vec{v} = x\,x' + y\,y'. \]
\end{definition}

\begin{propriete}[Condition d'orthogonalité]
Deux vecteurs $\vec{u}$ et $\vec{v}$ sont \cle{orthogonaux} (perpendiculaires) \textbf{si et seulement si} leur produit scalaire est nul :
\[ \vec{u}\cdot\vec{v} = 0 \quad\iff\quad x\,x' + y\,y' = 0. \]
\end{propriete}

\begin{attention}[Hypothèse orthonormée indispensable]
La condition $xx'+yy'=0$ n'est valable \textbf{que dans un repère orthonormé}. Dans un repère non orthonormé, l'angle entre les axes n'est pas $90^\circ$ et la formule ne donne pas l'orthogonalité géométrique.
\end{attention}

\begin{aretenir}[Propriétés du produit scalaire (rappels utiles)]
Pour tous vecteurs $\vec{u}, \vec{v}, \vec{w}$ et tout réel $k$ :
\begin{itemize}
\item \textbf{Symétrie :} $\vec{u}\cdot\vec{v} = \vec{v}\cdot\vec{u}$.
\item \textbf{Bilinéarité :} $(k\vec{u})\cdot\vec{v} = k(\vec{u}\cdot\vec{v})$ ; $(\vec{u}+\vec{v})\cdot\vec{w} = \vec{u}\cdot\vec{w} + \vec{v}\cdot\vec{w}$.
\item \textbf{Norme :} $\vec{u}\cdot\vec{u} = ||\vec{u}||^2 = x^2+y^2$.
\end{itemize}
\end{aretenir}

\courssub{Applications : angle droit, hauteur, médiatrice}

\begin{propriete}[Triangle rectangle]
Un triangle $ABC$ est rectangle en $A$ \textbf{si et seulement si} $\overrightarrow{AB}\cdot\overrightarrow{AC} = 0$.
\end{propriete}

\begin{propriete}[Droite perpendiculaire]
Une droite $(AB)$ est perpendiculaire à une droite $(CD)$ \textbf{si et seulement si} $\overrightarrow{AB}\cdot\overrightarrow{CD} = 0$.
\end{propriete}

\begin{methode}[Montrer un angle droit ou des droites perpendiculaires]
\begin{enumerate}
\item Vérifier que le repère est orthonormé.
\item Calculer les coordonnées des deux vecteurs concernés (côtés de l'angle ou directeurs des droites).
\item Calculer le produit scalaire $xx'+yy'$.
\item Si le résultat est $0$, conclure : « Le produit scalaire est nul, donc les vecteurs sont orthogonaux, donc l'angle est droit / les droites sont perpendiculaires. »
\end{enumerate}
\end{methode}

\begin{exoresolu}[Triangle rectangle et hauteur]
$A(0\,;0)$, $B(4\,;0)$, $C(1\,;3)$.
\begin{enumerate}
\item Montrer que $ABC$ est rectangle en $A$.
\item Déterminer une équation de la hauteur issue de $C$.
\end{enumerate}
\tcblower
\textbf{1.} $\overrightarrow{AB}(4\,;0)$, $\overrightarrow{AC}(1\,;3)$.
$\overrightarrow{AB}\cdot\overrightarrow{AC} = 4\times 1 + 0\times 3 = 4 \neq 0$.
Attention ! Pas rectangle en $A$.
Testons en $B$ : $\overrightarrow{BA}(-4\,;0)$, $\overrightarrow{BC}(-3\,;3)$. Produit scalaire $= (-4)(-3)+0\times 3 = 12 \neq 0$.
Testons en $C$ : $\overrightarrow{CA}(-1\,;-3)$, $\overrightarrow{CB}(3\,;-3)$. Produit scalaire $= (-1)\times 3 + (-3)\times(-3) = -3+9=6 \neq 0$.
Ce triangle n'est \textbf{pas rectangle}. (C'est un triangle quelconque).

\textbf{Modification pour l'exemple :} Prenons $C(0\,;3)$. Alors $\overrightarrow{AB}(4\,;0)$, $\overrightarrow{AC}(0\,;3)$. Produit scalaire $= 0$. Rectangle en $A$.
Hauteur issue de $C$ : droite passant par $C$ et $\perp$ à $(AB)$. $(AB)$ est horizontale (vecteur $(1\,;0)$), la hauteur est verticale : équation $x=0$.
\end{exoresolu}

\begin{exemplebox}[Médiatrice par produit scalaire]
$A(-2\,;1)$, $B(4\,;3)$. Milieu $I(1\,;2)$. Vecteur directeur de $(AB)$ : $\overrightarrow{AB}(6\,;2)$.
La médiatrice est $\perp$ à $(AB)$ et passe par $I$.
Soit $M(x\,;y)$ sur la médiatrice. $\overrightarrow{IM}(x-1\,;y-2)$.
Condition : $\overrightarrow{IM}\cdot\overrightarrow{AB} = 0 \iff 6(x-1) + 2(y-2) = 0 \iff 6x+2y-10=0 \iff 3x+y=5$.
Équation de la médiatrice : $3x+y=5$.
\end{exemplebox}

\begin{exercice}[Entraînement : orthogonalité]
Dans un repère orthonormé, $A(1\,;1)$, $B(4\,;5)$, $C(6\,;2)$.
\begin{enumerate}
\item Calculer $\overrightarrow{AB}\cdot\overrightarrow{AC}$. En déduire la nature de l'angle $\widehat{BAC}$.
\item Déterminer une équation de la droite passant par $A$ et perpendiculaire à $(BC)$.
\item Le triangle $ABC$ est-il rectangle ? Si oui, en quel sommet ?
\end{enumerate}
\end{exercice}

\begin{corrige}[Exercice : orthogonalité]
\textbf{1.} $\overrightarrow{AB}(3\,;4)$, $\overrightarrow{AC}(5\,;1)$. PS $= 3\times5 + 4\times1 = 19 \neq 0$. Angle non droit.
\textbf{2.} $\overrightarrow{BC}(2\,;-3)$. Droite $\perp (BC)$ par $A$ : $\overrightarrow{AM}(x-1\,;y-1)\cdot(2\,;-3)=0 \iff 2x-2 -3y+3=0 \iff 2x-3y+1=0$.
\textbf{3.} $\overrightarrow{BA}(-3\,;-4)$, $\overrightarrow{BC}(2\,;-3)$. PS $= -6+12=6\neq0$. $\overrightarrow{CA}(-5\,;-1)$, $\overrightarrow{CB}(-2\,;3)$. PS $= 10-3=7\neq0$. Triangle non rectangle.
\end{corrige}

\coursec{Synthèse et Résolution de problèmes complexes}

Dans les sections précédentes, nous avons découvert la \cle{condition d'orthogonalité} : deux vecteurs $\vec{u}(x\,;y)$ et $\vec{v}(x'\,;y')$ sont orthogonaux si et seulement si $xx' + yy' = 0$. Nous disposons désormais de \textbf{quatre outils algébriques} pour étudier la géométrie : les coordonnées d'un vecteur, la distance, la colinéarité et l'orthogonalité. Cette dernière section a un objectif ambitieux et passionnant : apprendre à \textbf{résoudre des problèmes complets} en combinant ces outils, comme on te le demandera au BEPC. C'est toute la puissance de la \cle{géométrie analytique} : transformer une figure en calculs, et les calculs en conclusions géométriques.

\courssub{La boîte à outils du géomètre analytique}

Avant de résoudre des problèmes, rassemblons dans un tableau tout ce que nous avons appris dans cette leçon. Ce tableau est ton meilleur ami : apprends-le par c\oe ur.

\begin{aretenir}[Boîte à outils : les 4 formules clés]
Dans un repère \textbf{orthonormé} $(O\,;\vec{\imath},\vec{\jmath})$, soient $A(x_A\,;y_A)$, $B(x_B\,;y_B)$, $\vec{u}(x\,;y)$ et $\vec{v}(x'\,;y')$.
\begin{center}
\begin{tabularx}{\linewidth}{|l|X|}
\hline
\rowcolor{popblueL} \textbf{Outil} & \textbf{Formule} \\
\hline
Coordonnées de $\overrightarrow{AB}$ & $\overrightarrow{AB}(x_B - x_A \ ; \ y_B - y_A)$ \\
\hline
Distance $AB$ (norme de $\overrightarrow{AB}$) & $AB = \sqrt{(x_B - x_A)^2 + (y_B - y_A)^2}$ \\
\hline
Colinéarité de $\vec{u}$ et $\vec{v}$ & $xy' - x'y = 0$ \quad (déterminant nul) \\
\hline
Orthogonalité de $\vec{u}$ et $\vec{v}$ & $xx' + yy' = 0$ \quad (produit scalaire nul) \\
\hline
\end{tabularx}
\end{center}
On retiendra aussi : le milieu $I$ de $[AB]$ a pour coordonnées $I\left(\dfrac{x_A+x_B}{2}\,;\dfrac{y_A+y_B}{2}\right)$.
\end{aretenir}

\begin{attention}[Pièges classiques à éviter absolument]
\begin{itemize}
\item \textbf{Repère non orthonormé} : les formules de distance et d'orthogonalité sont \emph{fausses} si le repère n'est pas orthonormé. Vérifie toujours cette hypothèse dans l'énoncé avant de calculer $AB$ ou un produit scalaire.
\item \textbf{Le vecteur nul} : $\vec{0}$ est colinéaire à \emph{tout} vecteur (son déterminant avec n'importe quel vecteur vaut $0$), mais on ne parle pas d'orthogonalité pour le vecteur nul en géométrie (un « angle droit » avec un vecteur nul n'a pas de sens).
\item \textbf{Division par zéro} : quand tu cherches un réel $k$ tel que $\vec{u} = k\vec{v}$, tu peux diviser par une coordonnée de $\vec{v}$ \emph{uniquement si elle est non nulle}. Si une coordonnée est nulle, utilise plutôt le déterminant $xy' - x'y = 0$, qui fonctionne dans tous les cas.
\item \textbf{Confusion des formules} : colinéarité $\to$ produit en croix $xy' - x'y$ ; orthogonalité $\to$ somme de produits $xx' + yy'$. Ne les intervertis pas !
\item \textbf{Sens du vecteur} : $\overrightarrow{AB}(x_B-x_A\,;y_B-y_A)$. Ne fais pas l'inverse ($x_A-x_B$) sauf si tu as besoin de $\overrightarrow{BA}$.
\end{itemize}
\end{attention}

\courssub{Stratégie générale de résolution : les 6 étapes}

Face à un problème de géométrie repérée, le plus grand danger est de se lancer dans des calculs au hasard. Voici la démarche méthodique qui mène au succès.

\begin{methode}[Résoudre un problème géométrique par les coordonnées : les 6 étapes]
\begin{enumerate}
\item \textbf{Choisir ou identifier le repère.} Si l'énoncé n'impose pas de repère, en choisir un adapté (souvent un sommet de la figure comme origine). Vérifier qu'il est orthonormé.
\item \textbf{Noter les coordonnées connues.} Lire ou calculer les coordonnées de tous les points donnés.
\item \textbf{Exprimer les vecteurs utiles.} Calculer les coordonnées des vecteurs qui interviennent dans la propriété à démontrer (côtés, diagonales, hauteurs...).
\item \textbf{Traduire la propriété géométrique en équation algébrique.}
      \begin{itemize}
      \item Alignement ou parallélisme $\to$ déterminant nul ($xy'-x'y=0$).
      \item Angle droit, perpendiculaires $\to$ produit scalaire nul ($xx'+yy'=0$).
      \item Longueurs égales $\to$ distances égales (ou carrés des distances égaux).
      \item Milieu $\to$ moyennes des coordonnées.
      \end{itemize}
\item \textbf{Résoudre} l'équation ou vérifier l'égalité obtenue.
\item \textbf{Interpréter et conclure} en revenant à la géométrie : « donc le triangle est rectangle en... », « donc les points sont alignés », « donc $ABCD$ est un parallélogramme », etc.
\end{enumerate}
\end{methode}

L'arbre de décision suivant t'aide à choisir le bon outil en fonction de la question posée.

\begin{popfigure}
\begin{center}
\begin{tikzpicture}[
node distance=1.1cm and 1.8cm,
qbox/.style={rectangle, rounded corners, draw=popblue, fill=popblueL, text width=3.8cm, align=center, font=\small\bfseries},
tool/.style={rectangle, rounded corners, draw=poporange, fill=poporangeL, text width=3.6cm, align=center, font=\small},
arr/.style={-{Stealth}, thick, popdark}]
\node[qbox] (q) {Que veut-on montrer ou calculer ?};
\node[qbox, below left=1.2cm and 2.5cm of q] (d) {Une longueur, un périmètre, un trajet ?};
\node[qbox, below=1.2cm of q] (c) {Alignement de points ou droites parallèles ?};
\node[qbox, below right=1.2cm and 2.5cm of q] (o) {Angle droit, droites perpendiculaires, hauteur ?};
\node[tool, below=1cm of d] (td) {Distance : $AB=\sqrt{(x_B-x_A)^2+(y_B-y_A)^2}$};
\node[tool, below=1cm of c] (tc) {Déterminant : $xy'-x'y=0$};
\node[tool, below=1cm of o] (to) {Produit scalaire : $xx'+yy'=0$};
\draw[arr] (q) -- (d);
\draw[arr] (q) -- (c);
\draw[arr] (q) -- (o);
\draw[arr] (d) -- (td);
\draw[arr] (c) -- (tc);
\draw[arr] (o) -- (to);
\end{tikzpicture}
\end{center}
\legende{Arbre de décision : quel outil utiliser selon la question posée ?}
\end{popfigure}

\courssub{Problème type 1 : étude d'un quadrilatère}

C'est le grand classique du BEPC : on donne les coordonnées de quatre points et on demande la nature du quadrilatère. Rappelons la stratégie :
\begin{itemize}
\item \textbf{Parallélogramme :} $\overrightarrow{AB} = \overrightarrow{DC}$ (côtés opposés parallèles et égaux) \textbf{ou} les diagonales ont même milieu.
\item \textbf{Rectangle :} Parallélogramme + un angle droit (produit scalaire nul) \textbf{ou} diagonales égales.
\item \textbf{Losange :} Parallélogramme + deux côtés consécutifs égaux \textbf{ou} diagonales perpendiculaires.
\item \textbf{Carré :} Rectangle + Losange (côtés égaux et angle droit).
\end{itemize}

\begin{exoresolu}[Type BEPC : triangle rectangle puis rectangle]
On se place dans un repère orthonormé $(O\,;\vec{\imath},\vec{\jmath})$. Soient $A(1\,;2)$, $B(5\,;4)$ et $C(3\,;8)$.
\begin{enumerate}
\item Démontrer que le triangle $ABC$ est rectangle en $B$.
\item Calculer l'aire du triangle $ABC$.
\item Déterminer les coordonnées du point $D$ tel que $ABCD$ soit un rectangle.
\end{enumerate}
\tcblower
\textbf{1.} On calcule les vecteurs des deux côtés issus de $B$ :
\[ \overrightarrow{BA}(1-5\,;2-4) = (-4\,;-2) \qquad \overrightarrow{BC}(3-5\,;8-4) = (-2\,;4). \]
Le repère est orthonormé, donc on peut utiliser la condition d'orthogonalité :
\[ x_{\overrightarrow{BA}}\,x_{\overrightarrow{BC}} + y_{\overrightarrow{BA}}\,y_{\overrightarrow{BC}} = (-4)\times(-2) + (-2)\times 4 = 8 - 8 = 0. \]
Or le produit scalaire est nul, donc $\overrightarrow{BA}$ et $\overrightarrow{BC}$ sont orthogonaux : \textbf{le triangle $ABC$ est rectangle en $B$}.

\textbf{2.} Calculons les longueurs des côtés de l'angle droit :
\[ BA = \sqrt{(-4)^2+(-2)^2} = \sqrt{20} = 2\sqrt{5}, \qquad BC = \sqrt{(-2)^2+4^2} = \sqrt{20} = 2\sqrt{5}. \]
L'aire d'un triangle rectangle est la moitié du produit des côtés de l'angle droit :
\[ \mathcal{A} = \frac{BA \times BC}{2} = \frac{2\sqrt{5}\times 2\sqrt{5}}{2} = \frac{20}{2} = 10. \]
\textbf{L'aire de $ABC$ vaut $10$ unités d'aire.} (On remarque au passage que $BA = BC$ : le triangle est rectangle \emph{et isocèle} en $B$.)

\textbf{3.} $ABCD$ est un parallélogramme si et seulement si $\overrightarrow{AB} = \overrightarrow{DC}$. Posons $D(x\,;y)$. On a $\overrightarrow{AB}(4\,;2)$ et $\overrightarrow{DC}(3-x\,;8-y)$. Donc :
\[ \begin{cases} 3 - x = 4 \\ 8 - y = 2 \end{cases} \iff \begin{cases} x = -1 \\ y = 6 \end{cases} \]
Donc $D(-1\,;6)$. Comme $ABCD$ est un parallélogramme possédant un angle droit en $B$ (question 1), \textbf{$ABCD$ est bien un rectangle} (et même un carré, car $BA = BC$).
\end{exoresolu}

\begin{popfigure}
\begin{center}
\begin{tikzpicture}[scale=0.62]
\draw[popline, dashed, step=1] (-2.5,-0.5) grid (6.5,9.5);
\draw[-{Stealth}, thick] (-2.5,0) -- (6.5,0) node[below left]{$x$};
\draw[-{Stealth}, thick] (0,-0.5) -- (0,9.5) node[below left]{$y$};
\node[below left] at (0,0) {\small $O$};
\coordinate (A) at (1,2); \coordinate (B) at (5,4);
\coordinate (C) at (3,8); \coordinate (D) at (-1,6);
\draw[very thick, popblue] (A) -- (B) -- (C) -- (D) -- cycle;
\draw[poporange, thick, dashed] (A) -- (C);
\draw[poporange, thick, dashed] (B) -- (D);
\draw[popgreen, thick, -{Stealth}] (A) -- (B) node[midway, below right, font=\small]{$\overrightarrow{AB}(4;2)$};
\draw[popgreen, thick, -{Stealth}] (B) -- (C) node[midway, right, font=\small]{$\overrightarrow{BC}(-2;4)$};
\draw[popdark] (4.4,4) -- (4.4,4.8) -- (3.6,4.8);
\foreach \p/\n/\pos in {A/A/below, B/B/right, C/C/above, D/D/left}
  \fill[popink] (\p) circle (2.2pt) node[\pos, font=\small]{\n};
\node[font=\small, poporange] at (2,5.6) {diagonale $[AC]$};
\node[font=\small, poporange, rotate=-25] at (2.9,3.4) {diagonale $[BD]$};
\end{tikzpicture}
\end{center}
\legende{Le quadrilatère $ABCD$ : côtés (vecteurs), angle droit en $B$ et diagonales. Toutes les vérifications se font par le calcul.}
\end{popfigure}

\courssub{Problème type 2 : médiatrice, hauteur et médiane d'un triangle}

Rappelons les définitions : la \cle{médiatrice} de $[AB]$ est l'ensemble des points équidistants de $A$ et $B$ ; la \cle{hauteur} issue de $C$ est la droite passant par $C$ et perpendiculaire à $(AB)$ ; la \cle{médiane} issue de $C$ passe par $C$ et par le milieu de $[AB]$. Chacune se traduit algébriquement.

\begin{exoresolu}[Médiatrice par le calcul]
Dans un repère orthonormé, soient $A(1\,;2)$ et $B(5\,;4)$. Déterminer une relation vérifiée par les coordonnées $(x\,;y)$ de tout point $M$ de la médiatrice de $[AB]$.
\tcblower
$M(x\,;y)$ appartient à la médiatrice de $[AB]$ si et seulement si $MA = MB$, c'est-à-dire $MA^2 = MB^2$ (des longueurs positives égales ont des carrés égaux). Or :
\begin{align*}
MA^2 &= (x-1)^2 + (y-2)^2 = x^2 - 2x + 1 + y^2 - 4y + 4,\\
MB^2 &= (x-5)^2 + (y-4)^2 = x^2 - 10x + 25 + y^2 - 8y + 16.
\end{align*}
L'égalité $MA^2 = MB^2$ donne, après simplification des $x^2$ et $y^2$ :
\[ -2x - 4y + 5 = -10x - 8y + 41 \iff 8x + 4y - 36 = 0 \iff 2x + y = 9. \]
\textbf{Conclusion} : la médiatrice de $[AB]$ est l'ensemble des points $M(x\,;y)$ tels que $2x + y = 9$ : c'est une droite, ce que l'on savait déjà géométriquement. Vérification : le milieu $I(3\,;3)$ de $[AB]$ satisfait $2\times 3 + 3 = 9$, il est donc bien sur la médiatrice.
\end{exoresolu}

\begin{methode}[Trouver la hauteur ou la médiane]
\begin{itemize}
\item \textbf{Hauteur issue de $C$ dans $\triangle ABC$ :} $M(x\,;y)$ est sur la hauteur si $\overrightarrow{CM}$ et $\overrightarrow{AB}$ sont orthogonaux, c'est-à-dire $x_{\overrightarrow{CM}}\,x_{\overrightarrow{AB}} + y_{\overrightarrow{CM}}\,y_{\overrightarrow{AB}} = 0$.
\item \textbf{Médiane issue de $C$ :} calculer le milieu $I$ de $[AB]$, puis écrire que $\overrightarrow{CM}$ et $\overrightarrow{CI}$ sont colinéaires (déterminant nul : $x_{\overrightarrow{CM}}y_{\overrightarrow{CI}} - x_{\overrightarrow{CI}}y_{\overrightarrow{CM}} = 0$).
\item \textbf{Médiatrice de $[AB]$ :} écrire $MA^2 = MB^2$ (ou utiliser le produit scalaire $\overrightarrow{IM}\cdot\overrightarrow{AB}=0$ avec $I$ milieu).
\end{itemize}
\end{methode}

\courssub{Problème type 3 : lieu de points}

Un \cle{lieu de points} est l'ensemble de tous les points $M$ vérifiant une condition donnée. La méthode est toujours la même : poser $M(x\,;y)$, traduire la condition par une équation, puis reconnaître la nature de l'ensemble (droite, cercle, demi-droite, parabole...). Au programme de 3ème, on rencontre surtout des droites et des cercles.

\begin{exoresolu}[Lieu de points : $MA^2 - MB^2 = k$]
Dans un repère orthonormé, soient $A(0\,;0)$ et $B(4\,;0)$. Déterminer l'ensemble des points $M(x\,;y)$ tels que $MA^2 - MB^2 = 8$.
\tcblower
On calcule : $MA^2 = x^2 + y^2$ et $MB^2 = (x-4)^2 + y^2 = x^2 - 8x + 16 + y^2$. Donc :
\[ MA^2 - MB^2 = (x^2+y^2) - (x^2 - 8x + 16 + y^2) = 8x - 16. \]
La condition $MA^2 - MB^2 = 8$ s'écrit $8x - 16 = 8$, soit $8x = 24$, soit $x = 3$.
\textbf{Conclusion} : l'ensemble cherché est la droite d'équation $x = 3$, c'est-à-dire la droite verticale passant par le point $(3\,;0)$. Les carrés $x^2$ et $y^2$ se sont simplifiés : c'est pourquoi le lieu est toujours une \textbf{droite} pour une condition du type $MA^2 - MB^2 = k$.
\end{exoresolu}

\begin{exemplebox}[Lieu de points : $MA = k \cdot MB$ (Cercle d'Apollonius)]
Soient $A(0\,;0)$, $B(4\,;0)$. Trouver le lieu des points $M$ tels que $MA = 2\,MB$.
$MA^2 = 4\,MB^2 \iff x^2+y^2 = 4[(x-4)^2+y^2] = 4(x^2-8x+16+y^2)$.
$x^2+y^2 = 4x^2 - 32x + 64 + 4y^2 \iff 0 = 3x^2 + 3y^2 - 32x + 64$.
Divisons par 3 : $x^2+y^2 - \frac{32}{3}x + \frac{64}{3} = 0$.
C'est l'équation d'un \textbf{cercle} (compléter le carré pour trouver centre et rayon). Au BEPC, on s'arrête souvent à l'équation cartésienne simplifiée.
\end{exemplebox}

\courssub{Modélisation de situations de la vie courante}

La géométrie analytique n'est pas qu'un exercice scolaire : elle sert chaque jour au Cameroun.

\begin{itemize}
\item \textbf{Alignement de poteaux électriques (ENEO).} Une équipe plante des poteaux le long d'une route. Sur le plan (repère orthonormé, unité : 100 m), trois poteaux sont en $P_1(0\,;1)$, $P_2(2\,;3)$ et $P_3(5\,;7)$. Sont-ils alignés ? On a $\overrightarrow{P_1P_2}(2\,;2)$ et $\overrightarrow{P_1P_3}(5\,;6)$. Déterminant : $2\times 6 - 5\times 2 = 12 - 10 = 2 \neq 0$. Donc les poteaux \textbf{ne sont pas alignés} : il faut repositionner $P_3$.
\item \textbf{Optimisation de trajet (Moto-taxi).} Un moto-taxi part du marché $M(1\,;1)$ pour rejoindre l'école $E(4\,;5)$ (unité : km). La distance à vol d'oiseau est $ME = \sqrt{3^2+4^2} = \sqrt{25} = 5$ km : c'est le trajet minimal possible, celui que la ligne droite réaliserait.
\item \textbf{Panneaux solaires (Énergie renouvelable).} Pour un rendement maximal, un panneau solaire doit être perpendiculaire aux rayons du soleil. Si la direction des rayons est $\vec{s}(3\,;-2)$ et celle du panneau $\vec{p}(2\,;a)$, la condition d'orthogonalité $3\times 2 + (-2)\times a = 0$ donne $6 - 2a = 0$, donc $a = 3$ : on règle l'inclinaison du panneau en conséquence.
\item \textbf{Cartographie et GPS.} Ton téléphone calcule sa position par trilatération : il mesure sa distance à au moins 3 satellites (coordonnées connues) et résout un système d'équations du type $(x-x_i)^2+(y-y_i)^2 = d_i^2$. C'est de la géométrie analytique en temps réel !
\end{itemize}

\courssub{Rédaction soignée et auto-évaluation}

Au BEPC, la qualité de la rédaction compte autant que le résultat. Utilise des \textbf{connecteurs logiques} : \emph{« or »} pour introduire une hypothèse ou un calcul connu, \emph{« donc »} pour une conclusion, \emph{« car »} pour une justification, \emph{« si... alors »} pour une implication. Encadre ou souligne tes résultats finaux. Exemple de phrase modèle : « \emph{Or} le repère est orthonormé ; \emph{de plus} $xx' + yy' = 0$, \emph{donc} les vecteurs sont orthogonaux, \emph{c'est-à-dire} que les droites sont perpendiculaires. »

\begin{attention}[Check-list avant de rendre ta copie]
\begin{itemize}
\item[$\square$] Ai-je vérifié que le repère est \textbf{orthonormé} avant d'utiliser distances et produit scalaire ?
\item[$\square$] Ai-je calculé $\overrightarrow{AB}$ dans le bon sens : $x_B - x_A$ et non $x_A - x_B$ ?
\item[$\square$] Ai-je distingué colinéarité ($xy' - x'y = 0$) et orthogonalité ($xx' + yy' = 0$) ?
\item[$\square$] Ai-je pensé à la \textbf{racine carrée} dans la formule de la distance ?
\item[$\square$] Ai-je \textbf{conclu en français} (alignés, rectangle, perpendiculaires...) et pas seulement par un calcul ?
\item[$\square$] Ai-je vérifié mon résultat sur un point connu ou par un ordre de grandeur ?
\end{itemize}
\end{attention}

\begin{savaistu}[Descartes, Fermat... et ton téléphone]
La géométrie analytique est née au XVII\textsuperscript{e} siècle avec René \textbf{Descartes} et Pierre de \textbf{Fermat}, qui ont eu l'idée révolutionnaire de réunir l'algèbre et la géométrie grâce aux coordonnées. Aujourd'hui, elle est partout : l'\textbf{infographie} et les jeux vidéo déplacent des personnages par des calculs de vecteurs ; la \textbf{robotique} pilote les bras articulés par coordonnées ; le \textbf{GPS} localise ton téléphone par trilatération, c'est-à-dire par des calculs de distances à plusieurs satellites. Dans les classes supérieures, tu utiliseras les coordonnées pour étudier des courbes et des transformations géométriques encore plus riches.
\end{savaistu}

\begin{exercice}[Pour t'entraîner : Synthèse BEPC]
Dans un repère orthonormé, soient $A(-2\,;1)$, $B(2\,;3)$ et $C(4\,;-1)$.
\begin{enumerate}
\item Montrer que $ABC$ est un triangle isocèle en $B$.
\item Le triangle $ABC$ est-il rectangle ? Justifier.
\item Déterminer les coordonnées de $D$ tel que $ABDC$ soit un parallélogramme.
\item Déterminer l'ensemble des points $M$ tels que $MA = MB$.
\end{enumerate}
\end{exercice}

\begin{corrige}[Exercice : triangle isocèle et parallélogramme]
\textbf{1.} Calculons les trois longueurs :
$AB = \sqrt{(2-(-2))^2 + (3-1)^2} = \sqrt{4^2+2^2} = \sqrt{20}$.
$BC = \sqrt{(4-2)^2 + (-1-3)^2} = \sqrt{2^2+(-4)^2} = \sqrt{20}$.
$AC = \sqrt{(4-(-2))^2 + (-1-1)^2} = \sqrt{6^2+(-2)^2} = \sqrt{40}$.
On a $AB = BC = \sqrt{20} \neq AC$. Le triangle est \textbf{isocèle en $B$} (sommet commun aux deux côtés égaux).

\textbf{2.} Vérifions l'angle en $B$ (sommet de l'angle isocèle) :
$\overrightarrow{BA}(-4\,;-2)$, $\overrightarrow{BC}(2\,;-4)$.
Produit scalaire : $(-4)\times 2 + (-2)\times(-4) = -8 + 8 = 0$.
Donc $\overrightarrow{BA} \perp \overrightarrow{BC}$. $ABC$ est \textbf{rectangle en $B$}.
C'est un triangle \textbf{rectangle isocèle en $B$}.

\textbf{3.} $ABDC$ parallélogramme $\iff \overrightarrow{AB} = \overrightarrow{CD}$.
$\overrightarrow{AB}(4\,;2)$. Soit $D(x\,;y)$, $\overrightarrow{CD}(x-4\,;y+1)$.
$\begin{cases} x-4 = 4 \\ y+1 = 2 \end{cases} \iff \begin{cases} x = 8 \\ y = 1 \end{cases}$.
Donc $D(8\,;1)$.

\textbf{4.} $MA = MB \iff M$ appartient à la médiatrice de $[AB]$.
$MA^2 = MB^2 \iff (x+2)^2+(y-1)^2 = (x-2)^2+(y-3)^2$.
$x^2+4x+4+y^2-2y+1 = x^2-4x+4+y^2-6y+9$.
$4x - 2y + 5 = -4x - 6y + 13 \iff 8x + 4y = 8 \iff 2x + y = 2$.
L'ensemble est la \textbf{droite d'équation $2x+y=2$} (médiatrice de $[AB]$).
\end{corrige}

\newpage

\coursec{Activité d'intégration}

\courssub{Objectifs de l'activité}
Cette activité d'intégration a pour but de mobiliser l'ensemble des savoirs et savoir-faire de la leçon « Coordonnées d'un vecteur » dans une situation unique, contextualisée et complète. Elle vise à :
\begin{itemize}
    \item Vérifier la maîtrise du calcul des coordonnées d'un vecteur, de sa norme et de la distance entre deux points.
    \item Évaluer la capacité à utiliser les conditions de colinéarité (déterminant nul) et d'orthogonalité (produit scalaire nul) pour résoudre des problèmes géométriques.
    \item Mettre en œuvre la résolution de systèmes d'équations pour déterminer un projeté orthogonal.
    \item Appliquer la caractérisation vectorielle du parallélogramme et la propriété des diagonales.
    \item Développer la compétence « Rédiger et justifier une démarche » en produisant un rapport structuré, argumenté et mathématiquement rigoureux, comme attendu au BEPC.
\end{itemize}

\courssub{Situation}
Dans le cadre du plan d'urbanisme de la commune de \textbf{Nouveau Départ}, une zone d'extension située en périphérie de Yaoundé, le conseil municipal doit valider le tracé des voiries et l'implantation des équipements publics avant le lancement des travaux. Le plan de masse est fourni sur un quadrillage orthonormé où \textbf{1 cm représente 10 m sur le terrain}. Le repère orthonormé $(O;\vec{\imath},\vec{\jmath})$ est matérialisé par les lignes du quadrillage.

Trois points stratégiques sont déjà identifiés :
\begin{itemize}
    \item La \textbf{Mairie annexe} $A(2\,;\,1)$,
    \item L'\textbf{École publique} $B(8\,;\,1)$,
    \item Le \textbf{Puits communautaire} $C(5\,;\,6)$.
\end{itemize}

Le maire et son bureau d'études doivent prendre cinq décisions techniques majeures qui conditionnent la viabilité du quartier. Votre mission, en tant qu'expert géomètre-topographe junior, est de fournir les justifications mathématiques exactes pour chacune de ces décisions.

\begin{popfigure}
\begin{tikzpicture}[scale=0.7, >=Stealth]
    % Grille
    \draw[popmuted, step=1cm] (0,0) grid (12,8);
    % Axes
    \draw[popdark, thick, ->] (0,0) -- (12.5,0) node[right] {$x$};
    \draw[popdark, thick, ->] (0,0) -- (0,8.5) node[above] {$y$};
    \foreach \x in {1,...,12} \draw (\x,0.1) -- (\x,-0.1) node[below, font=\small] {\x};
    \foreach \y in {1,...,8} \draw (0.1,\y) -- (-0.1,\y) node[left, font=\small] {\y};
    % Points principaux
    \coordinate (A) at (2,1);
    \coordinate (B) at (8,1);
    \coordinate (C) at (5,6);
    \fill[popblue] (A) circle (3pt) node[below left] {$A(2;1)$ Mairie};
    \fill[popgreen] (B) circle (3pt) node[below right] {$B(8;1)$ École};
    \fill[poporange] (C) circle (3pt) node[above] {$C(5;6)$ Puits};
    % Triangle ABC
    \draw[popblue, thick] (A) -- (B) -- (C) -- cycle;
    % Droite (BC) complète : intersections avec le cadre [0,12]x[0,8]
    % 5x+3y-43=0 -> x=0 => y=14.3 (out) ; y=8 => x=3.8 ; y=0 => x=8.6 ; x=12 => y=-5.6 (out)
    \draw[popgreen!50!popdark, thick] (3.8,8) -- (8.6,0);
    % Route R par A // BC : 5x+3y-13=0 -> x=0 => y=4.33 ; y=0 => x=2.6
    \draw[poppurple, thick, dashed] (0,4.33) -- (2.6,0) node[below right, font=\small, poppurple] {Route $\mathcal{R}$};
    % Point L (projeté orthogonal) L(109/17, 62/17) approx (6.41, 3.65)
    \coordinate (L) at (6.41, 3.65);
    \fill[poppink] (L) circle (3pt) node[right] {$L$ Lampadaire};
    \draw[poppink, thick] (A) -- (L);
    % Point D (11,6)
    \coordinate (D) at (11,6);
    \fill[popgold] (D) circle (3pt) node[above right] {$D(11;6)$ Lotissement};
    % Parallélogramme ABDC
    \draw[popgold, thick] (A) -- (B) -- (D) -- (C) -- cycle;
    % Diagonales
    \draw[popmuted, thin, densely dotted] (A) -- (D);
    \draw[popmuted, thin, densely dotted] (B) -- (C);
    % Milieu I
    \coordinate (I) at (6.5, 3.5);
    \fill[popdark] (I) circle (2pt) node[below left, font=\tiny] {$I$};
    % Échelle
    \draw[popdark, thick] (0.5,-0.8) -- (1.5,-0.8) node[midway, below] {1 cm = 10 m};
\end{tikzpicture}
\legende{Plan de masse du quartier « Nouveau Départ » (repère orthonormé, 1 cm = 10 m).}
\end{popfigure}

\courssub{Consignes}
Sur la base du plan ci-dessus et des coordonnées données, traiter les cinq points suivants en justifiant chaque étape par un calcul de coordonnées ou une propriété vectorielle.

\begin{enumerate}
    \item \textbf{Étude du triangle $ABC$ :} Calculer les distances réelles (en mètres) $AB$, $AC$ et $BC$. Déterminer la nature du triangle $ABC$ (isocèle ? rectangle ? équilatéral ?). Justifier par les calculs.
    \item \textbf{Tracé de la route principale :} La mairie demande une route droite $\mathcal{R}$ passant par $A$ et \textbf{parallèle} à la droite $(BC)$. Déterminer une équation cartésienne de $\mathcal{R}$.
    \item \textbf{Implantation d'un lampadaire :} On veut placer un lampadaire $L$ sur la droite $(BC)$ de sorte que le segment $[AL]$ soit \textbf{perpendiculaire} à $(BC)$. 
    \begin{enumerate}
        \item Déterminer les coordonnées exactes de $L$ (projeté orthogonal de $A$ sur $(BC)$).
        \item Calculer la distance $AL$ en mètres. Interpréter ce résultat pour l'éclairage public.
    \end{enumerate}
    \item \textbf{Nouveau lotissement :} Un lotissement doit former un parallélogramme $ABDC$ (dans cet ordre). 
    \begin{enumerate}
        \item Déterminer les coordonnées du point $D$.
        \item Vérifier par le calcul que les diagonales $[AD]$ et $[BC]$ ont le même milieu.
    \end{enumerate}
    \item \textbf{Rapport final :} Rédiger le rapport technique adressé au maire, synthétisant les résultats des questions 1 à 4. Chaque décision géométrique doit être justifiée par un calcul de coordonnées ou une propriété vectorielle clairement énoncée.
\end{enumerate}

\courssub{Ressources à mobiliser}
\begin{aretenir}[Formulaire de synthèse — Vecteurs dans un repère orthonormé]
\textbf{1. Coordonnées d'un vecteur :} $\vec{AB} \begin{pmatrix} x_B - x_A \\ y_B - y_A \end{pmatrix}$.
\textbf{2. Norme et Distance :} $\|\vec{AB}\| = AB = \sqrt{(x_B-x_A)^2 + (y_B-y_A)^2}$.
\textbf{3. Colinéarité :} $\vec{u}\begin{pmatrix} x_1 \\ y_1 \end{pmatrix}$ et $\vec{v}\begin{pmatrix} x_2 \\ y_2 \end{pmatrix}$ sont colinéaires $\iff$ $\det(\vec{u},\vec{v}) = x_1y_2 - x_2y_1 = 0$.
\textbf{4. Orthogonalité (Produit scalaire) :} $\vec{u} \cdot \vec{v} = x_1x_2 + y_1y_2$. $\vec{u} \perp \vec{v} \iff \vec{u} \cdot \vec{v} = 0$.
\textbf{5. Équation cartésienne d'une droite :} Droite passant par $A(x_A,y_A)$ de vecteur directeur $\vec{u}\begin{pmatrix} a \\ b \end{pmatrix}$ ($a \neq 0$) : $y - y_A = \frac{b}{a}(x - x_A)$ ou forme $bx - ay + (ay_A - bx_A) = 0$.
\textbf{6. Parallélogramme :} $ABDC$ parallélogramme $\iff \vec{AB} = \vec{CD}$ $\iff$ $\vec{AC} = \vec{BD}$ $\iff$ Milieu de $[AD]$ = Milieu de $[BC]$.
\textbf{7. Milieu :} $I$ milieu de $[AB] \iff I\left(\frac{x_A+x_B}{2};\frac{y_A+y_B}{2}\right)$.
\end{aretenir}

\courssub{Démarche et corrigé commenté}
\begin{exoresolu}[Corrigé détaillé de l'activité d'intégration]
\textbf{Étape 1 : Calcul des vecteurs de base et distances.}
On détermine d'abord les vecteurs reliant les points remarquables :
\[
\vec{AB} = \begin{pmatrix} 8-2 \\ 1-1 \end{pmatrix} = \begin{pmatrix} 6 \\ 0 \end{pmatrix},\quad
\vec{AC} = \begin{pmatrix} 5-2 \\ 6-1 \end{pmatrix} = \begin{pmatrix} 3 \\ 5 \end{pmatrix},\quad
\vec{BC} = \begin{pmatrix} 5-8 \\ 6-1 \end{pmatrix} = \begin{pmatrix} -3 \\ 5 \end{pmatrix}.
\]
Les distances sur le plan (en cm) puis sur le terrain (en m, échelle 1/1000 soit 1 cm = 10 m) :
\begin{align*}
AB &= \|\vec{AB}\| = \sqrt{6^2 + 0^2} = 6 \text{ cm} \quad\Rightarrow\quad \textbf{60 m}.\\
AC &= \|\vec{AC}\| = \sqrt{3^2 + 5^2} = \sqrt{34} \text{ cm} \quad\Rightarrow\quad \textbf{10$\sqrt{34}$ m} \approx 58,3 \text{ m}.\\
BC &= \|\vec{BC}\| = \sqrt{(-3)^2 + 5^2} = \sqrt{34} \text{ cm} \quad\Rightarrow\quad \textbf{10$\sqrt{34}$ m} \approx 58,3 \text{ m}.
\end{align*}
\textbf{Nature du triangle :} On a $AC = BC = 10\sqrt{34}$ m. Le triangle $ABC$ est \textbf{isocèle en $C$} (sommet au puits).
Vérification rectangle : $\vec{AC} \cdot \vec{BC} = 3\times(-3) + 5\times 5 = -9 + 25 = 16 \neq 0$. $\vec{AB} \cdot \vec{AC} = 18 \neq 0$. $\vec{AB} \cdot \vec{BC} = -18 \neq 0$. Aucune orthogonalité. Le triangle n'est \textbf{pas rectangle}.

\textbf{Étape 2 : Équation de la route $\mathcal{R}$ (par $A$, parallèle à $(BC)$).}
La droite $(BC)$ a pour vecteur directeur $\vec{BC} = \begin{pmatrix} -3 \\ 5 \end{pmatrix}$.
La route $\mathcal{R}$ passant par $A(2;1)$ et parallèle à $(BC)$ a le même vecteur directeur.
Son coefficient directeur (pente) est $m = \frac{5}{-3} = -\frac{5}{3}$.
Équation réduite : $y - 1 = -\frac{5}{3}(x - 2)$.
On met sous forme cartésienne $ax+by+c=0$ :
\[
3(y-1) = -5(x-2) \iff 3y - 3 = -5x + 10 \iff \boxed{5x + 3y - 13 = 0}.
\]

\textbf{Étape 3 : Projeté orthogonal $L$ de $A$ sur $(BC)$ et distance $AL$.}
\emph{Méthode :} $L$ est l'intersection de $(BC)$ et de la perpendiculaire à $(BC)$ passant par $A$.
\begin{itemize}
    \item \textbf{Équation de $(BC)$ :} Passe par $B(8;1)$, directeur $\vec{BC}=\begin{pmatrix}-3\\5\end{pmatrix}$.
    $y - 1 = -\frac{5}{3}(x - 8) \iff 3y - 3 = -5x + 40 \iff \boxed{5x + 3y - 43 = 0}$.
    \item \textbf{Équation de $(AL)$ :} $(AL) \perp (BC) \iff \vec{AL} \cdot \vec{BC} = 0$.
    Soit $L(x_L;y_L)$. $\vec{AL} = \begin{pmatrix} x_L-2 \\ y_L-1 \end{pmatrix}$.
    Condition d'orthogonalité : $(x_L-2)(-3) + (y_L-1)(5) = 0 \iff -3x_L + 6 + 5y_L - 5 = 0 \iff \boxed{3x_L - 5y_L = 1}$.
\end{itemize}
$L$ est solution du système :
\[
\begin{cases}
5x_L + 3y_L = 43 \quad (1)\\
3x_L - 5y_L = 1 \quad (2)
\end{cases}
\]
Résolution par combinaison linéaire : $(1)\times 5 + (2)\times 3$ :
$25x_L + 15y_L + 9x_L - 15y_L = 215 + 3 \iff 34x_L = 218 \iff x_L = \frac{218}{34} = \frac{109}{17}$.
En remplaçant dans (2) : $3\times\frac{109}{17} - 5y_L = 1 \iff \frac{327}{17} - 5y_L = \frac{17}{17} \iff 5y_L = \frac{310}{17} \iff y_L = \frac{62}{17}$.
Donc \textbf{$L\left(\frac{109}{17}\,;\,\frac{62}{17}\right)$} (soit environ $L(6,41\,;\,3,65)$ sur le plan).

\emph{Distance $AL$ :}
$\vec{AL} = \begin{pmatrix} \frac{109}{17}-2 \\ \frac{62}{17}-1 \end{pmatrix} = \begin{pmatrix} \frac{75}{17} \\ \frac{45}{17} \end{pmatrix}$.
$AL = \|\vec{AL}\| = \sqrt{\left(\frac{75}{17}\right)^2 + \left(\frac{45}{17}\right)^2} = \frac{1}{17}\sqrt{5625 + 2025} = \frac{\sqrt{7650}}{17} = \frac{\sqrt{225\times 34}}{17} = \frac{15\sqrt{34}}{17} \text{ cm}$.
Distance réelle : $\boxed{AL = \frac{150\sqrt{34}}{17} \text{ m} \approx 51,4 \text{ m}}$.
\textit{Interprétation :} C'est la distance minimale entre la mairie et la voirie $(BC)$ ; le lampadaire $L$ est le point de $(BC)$ le plus proche de la mairie.

\textbf{Étape 4 : Construction du parallélogramme $ABDC$.}
L'ordre des sommets est $A \to B \to D \to C$. On a $\vec{AB} = \vec{CD}$.
$\vec{AB} = \begin{pmatrix} 6 \\ 0 \end{pmatrix}$. Soit $D(x_D;y_D)$. $\vec{CD} = \begin{pmatrix} x_D-5 \\ y_D-6 \end{pmatrix}$.
Égalité vectorielle :
\[
\begin{cases}
x_D - 5 = 6 \Rightarrow x_D = 11\\
y_D - 6 = 0 \Rightarrow y_D = 6
\end{cases}
\]
Donc \textbf{$D(11\,;\,6)$}.
\textit{Vérification alternative :} $\vec{AC} = \begin{pmatrix} 3 \\ 5 \end{pmatrix}$ et $\vec{BD} = \begin{pmatrix} 11-8 \\ 6-1 \end{pmatrix} = \begin{pmatrix} 3 \\ 5 \end{pmatrix}$. $\vec{AC} = \vec{BD}$, bien un parallélogramme.

\textbf{Milieux des diagonales :}
Diagonale $[AD]$ : $I_{AD}\left(\frac{2+11}{2}\,;\,\frac{1+6}{2}\right) = \left(\frac{13}{2}\,;\,\frac{7}{2}\right) = (6,5\,;\,3,5)$.
Diagonale $[BC]$ : $I_{BC}\left(\frac{8+5}{2}\,;\,\frac{1+6}{2}\right) = \left(\frac{13}{2}\,;\,\frac{7}{2}\right) = (6,5\,;\,3,5)$.
Les milieux coïncident : \textbf{les diagonales se coupent en leur milieu commun $I\left(\frac{13}{2}\,;\,\frac{7}{2}\right)$}.

\textbf{Étape 5 : Structure du rapport final (modèle).}
\begin{quote}
\textsc{Rapport technique n° 2026/03 — Aménagement « Nouveau Départ »}\\
\textbf{Objet :} Validation géométrique du plan de masse.\\
\textbf{1. Implantation existante :} Le triangle Mairie-École-Puits ($ABC$) est isocèle en $C$ ($AC=BC\approx 58,3$ m), non rectangle. Distance Mairie-École : 60 m.\\
\textbf{2. Voirie principale :} La route $\mathcal{R}$ (équation $5x+3y-13=0$) relie la mairie et est parallèle à l'axe Puits-École $(BC)$, assurant la continuité du réseau.\\
\textbf{3. Éclairage public :} Le lampadaire $L$ est implanté au projeté orthogonal de la mairie sur $(BC)$ : $L(109/17;62/17)$. Distance mairie-lampadaire $AL \approx 51,4$ m (optimale pour la couverture).\\
\textbf{4. Extension lotissement :} Le parallélogramme $ABDC$ place le nouveau lot $D$ en $(11;6)$. La symétrie centrale (milieu commun $I(6,5;3,5)$ des diagonales) garantit l'équilibre du cadastre.\\
\textbf{Conclusion :} Tous les paramètres sont validés mathématiquement. Le dossier est transmis pour exécution.
\end{quote}
\tcblower
\textbf{Commentaires pédagogiques :}
\begin{itemize}
    \item \textbf{Échelle :} Ne pas oublier la conversion cm $\to$ m ($\times 10$) pour les distances finales. C'est une source fréquente d'erreur au BEPC.
    \item \textbf{Projeté orthogonal :} La méthode par système d'équations (droite support + orthogonalité) est la plus robuste en 3ème. La méthode paramétrique ($L = B + t\vec{BC}$) est élégante mais demande de maîtriser le paramètre $t$ (souvent vu en Seconde).
    \item \textbf{Parallélogramme :} L'ordre des lettres $ABDC$ impose $\vec{AB}=\vec{CD}$. Attention à ne pas confondre avec $ABCD$ qui donnerait $\vec{AB}=\vec{DC}$ (point $D$ différent).
    \item \textbf{Rédaction :} Au BEPC, « Justifier » exige d'écrire la propriété utilisée (ex: « $\vec{u}\cdot\vec{v}=0$ donc $(AL)\perp(BC)$ ») et le calcul numérique.
\end{itemize}
\end{exoresolu}

\courssub{Bilan de l'activité}
Cette activité d'intégration a permis de valider l'acquisition des quatre piliers de la leçon dans une situation unique et cohérente :
\begin{enumerate}
    \item \textbf{Coordonnées et distances :} Le calcul des vecteurs $\vec{AB}, \vec{AC}, \vec{BC}$ et de leurs normes a servi de base à toute la résolution. La conversion d'échelle (cm $\to$ m) ancre les mathématiques dans le réel.
    \item \textbf{Colinéarité :} Utilisée implicitement pour écrire l'équation de la route $\mathcal{R}$ parallèle à $(BC)$ (même vecteur directeur) et pour vérifier le parallélogramme ($\vec{AB}=\vec{CD}$).
    \item \textbf{Orthogonalité :} Clé de voûte de la question 3 (recherche du projeté orthogonal $L$ via $\vec{AL}\cdot\vec{BC}=0$) et de la vérification de la nature du triangle (produits scalaires nuls ?).
    \item \textbf{Synthèse vectorielle :} La question 4 (parallélogramme $ABDC$) a mobilisé l'égalité vectorielle $\vec{AB}=\vec{CD}$ et la propriété des diagonales (même milieu), reliant géométrie analytique et géométrie classique.
\end{enumerate}
Sur le plan des \textbf{savoir-faire}, l'élève a dû :
\begin{itemize}
    \item Organiser ses calculs de manière lisible (vecteurs, distances, systèmes).
    \item Passer du géométrique (figure) à l'analytique (coordonnées, équations) et inversement.
    \item Rédiger une argumentation complète (rapport final), compétence transversale essentielle pour l'épreuve écrite du BEPC.
\end{itemize}
\begin{attention}[Pièges évités dans cette activité]
\begin{itemize}
    \item Confondre « isocèle en $C$ » (côtés $AC$ et $BC$ égaux) et « isocèle en $A$ ».
    \item Oublier le changement d'échelle pour donner les distances en mètres.
    \item Écrire l'équation de la parallèle sans passer par le vecteur directeur (risque d'erreur sur le coefficient directeur).
    \item Chercher $L$ comme milieu de $[BC]$ au lieu du projeté orthogonal (erreur classique : confondre médiane et hauteur dans un triangle isocèle ; ici la hauteur issue de $A$ n'est pas médiane car $AB \neq AC$).
    \item Inverser l'ordre des points dans le parallélogramme ($ABCD$ au lieu de $ABDC$), ce qui place $D$ en $(-1;6)$ au lieu de $(11;6)$.
\end{itemize}
\end{attention}

\begin{savaistu}[Le géomètre-topographe au Cameroun]
Le métier de géomètre-topographe est au cœur de l'aménagement du territoire au Cameroun. De la délimitation des parcelles pour le titre foncier (réforme foncière en cours) au tracé des routes (projet d'autoroute Yaoundé-Douala), en passant par l'implantation des réseaux Eneo, Camwater ou Camtel, la maîtrise des coordonnées, distances et angles est une compétence technique quotidienne. Les outils modernes (GPS différentiel, stations totales, drones photogrammétriques) reposent tous sur les mêmes fondements mathématiques : le repérage dans un repère orthonormé (souvent Lambert ou UTM au Cameroun), le calcul vectoriel et la trigonométrie. Cette activité est une simulation simplifiée de ce que fait un jeune technicien géomètre au début de sa carrière au MINHDU (Ministère de l'Habitat et du Développement Urbain) ou dans un cabinet privé.
\end{savaistu}

\newpage

\coursec{Méthodes et exercices résolus}

\coursec{Coordonnées d'un vecteur}

\courssub{Coordonnées d'un vecteur dans un repère orthonormé}

\begin{definition}[Repère orthonormé et coordonnées d'un point]
    Un \textbf{repère orthonormé} $(O;\vec{\imath},\vec{\jmath})$ est formé d'un point origine $O$ et de deux vecteurs unitaires $\vec{\imath}$ et $\vec{\jmath}$ perpendiculaires ($\vec{\imath}\perp\vec{\jmath}$, $\|\vec{\imath}\|=\|\vec{\jmath}\|=1$).
    Tout point $M$ du plan se repère par un couple unique de nombres $(x_M;y_M)$ appelés \textbf{coordonnées} de $M$ dans ce repère, tels que $\overrightarrow{OM}=x_M\vec{\imath}+y_M\vec{\jmath}$.
    On note aussi $M(x_M;y_M)$.
\end{definition}

\begin{propriete}[Coordonnées d'un vecteur]
    Soient $A(x_A;y_A)$ et $B(x_B;y_B)$ deux points dans un repère orthonormé $(O;\vec{\imath},\vec{\jmath})$.
    Les coordonnées du vecteur $\overrightarrow{AB}$ sont :
    \[
    \overrightarrow{AB}\begin{pmatrix} x_B-x_A \\ y_B-y_A \end{pmatrix}
    \]
    \textbf{Démonstration :} $\overrightarrow{AB}=\overrightarrow{OB}-\overrightarrow{OA}=(x_B\vec{\imath}+y_B\vec{\jmath})-(x_A\vec{\imath}+y_A\vec{\jmath})=(x_B-x_A)\vec{\imath}+(y_B-y_A)\vec{\jmath}$.
\end{propriete}

\begin{aretenir}[Règle pratique : « Arrivée moins Départ »]
    Pour obtenir les coordonnées de $\overrightarrow{AB}$, on soustrait les coordonnées du point de départ $A$ aux coordonnées du point d'arrivée $B$ :
    \begin{center}
    $\overrightarrow{AB}\Bigl(\; \underbrace{x_B-x_A}_{\text{abscisse}} \;;\; \underbrace{y_B-y_A}_{\text{ordonnée}} \;\Bigr)$
    \end{center}
    Le vecteur nul $\vec{0}$ a pour coordonnées $\begin{pmatrix} 0 \\ 0 \end{pmatrix}$.
\end{aretenir}

\begin{exemplebox}[Repérage sur une carte du Cameroun]
    Dans un repère orthonormé $(O;\vec{\imath},\vec{\jmath})$ (unité : 10 km), on place :
    \begin{itemize}
        \item \textbf{Yaoundé} (capitale) en $Y(4;2)$,
        \item \textbf{Bafoussam} en $B(1;6)$,
        \item \textbf{Douala} (port) en $D(8;-1)$.
    \end{itemize}
    \begin{enumerate}
        \item Déterminer les coordonnées du vecteur $\overrightarrow{YB}$ (trajet Yaoundé $\to$ Bafoussam).
        \item Déterminer les coordonnées du vecteur $\overrightarrow{BD}$ (trajet Bafoussam $\to$ Douala).
        \item Que représente le vecteur $\overrightarrow{YD}$ ? Calculer ses coordonnées.
    \end{enumerate}
    \tcblower
    \textbf{Solution :}
    \begin{enumerate}
        \item $\overrightarrow{YB}\begin{pmatrix} 1-4 \\ 6-2 \end{pmatrix} = \begin{pmatrix} -3 \\ 4 \end{pmatrix}$. Déplacement de 30 km vers l'Ouest et 40 km vers le Nord.
        \item $\overrightarrow{BD}\begin{pmatrix} 8-1 \\ -1-6 \end{pmatrix} = \begin{pmatrix} 7 \\ -7 \end{pmatrix}$. Déplacement de 70 km vers l'Est et 70 km vers le Sud.
        \item $\overrightarrow{YD}$ représente le trajet direct Yaoundé $\to$ Douala. $\overrightarrow{YD}\begin{pmatrix} 8-4 \\ -1-2 \end{pmatrix} = \begin{pmatrix} 4 \\ -3 \end{pmatrix}$.
        \textit{Remarque :} $\overrightarrow{YD} = \overrightarrow{YB} + \overrightarrow{BD} = \begin{pmatrix} -3 \\ 4 \end{pmatrix} + \begin{pmatrix} 7 \\ -7 \end{pmatrix} = \begin{pmatrix} 4 \\ -3 \end{pmatrix}$ (Relation de Chasles).
    \end{enumerate}
\end{exemplebox}

\begin{popfigure}
\begin{tikzpicture}[scale=0.6, >=Stealth]
    % Grid
    \draw[popmuted, very thin] (-1,-3) grid (10,8);
    % Axes
    \draw[popdark, thick, ->] (-1,0) -- (10.5,0) node[right] {$x$ (Est/Ouest)};
    \draw[popdark, thick, ->] (0,-3) -- (0,8.5) node[above] {$y$ (Nord/Sud)};
    % Origin
    \node[below left] at (0,0) {$O$};
    % Points
    \node[popblue, circle, fill, inner sep=2pt, label=above right:{$Y(4;2)$}] (Y) at (4,2) {};
    \node[poporange, circle, fill, inner sep=2pt, label=above left:{$B(1;6)$}] (B) at (1,6) {};
    \node[popgreen, circle, fill, inner sep=2pt, label=below right:{$D(8;-1)$}] (D) at (8,-1) {};
    % Vectors
    \draw[popblue, thick, ->] (Y) -- (B) node[midway, above left, sloped] {$\overrightarrow{YB}(-3;4)$};
    \draw[poporange, thick, ->] (B) -- (D) node[midway, below right, sloped] {$\overrightarrow{BD}(7;-7)$};
    \draw[poppurple, very thick, ->, dashed] (Y) -- (D) node[midway, below right, sloped] {$\overrightarrow{YD}(4;-3)$};
    % Unit vectors
    \draw[popdark, thick, ->] (0,0) -- (1,0) node[midway, below] {$\vec{\imath}$};
    \draw[popdark, thick, ->] (0,0) -- (0,1) node[midway, left] {$\vec{\jmath}$};
\end{tikzpicture}
\legende{Repérage de villes et vecteurs déplacements (unité : 10 km).}
\end{popfigure}

\begin{attention}[Piège classique : L'ordre de la soustraction]
    $\overrightarrow{AB} \neq \overrightarrow{BA}$. On a $\overrightarrow{BA} = -\overrightarrow{AB}$.
    \textbf{Toujours faire :} Coordonnées de l'arrivée $-$ Coordonnées du départ.
    \textit{Astuce :} Le vecteur « tire » la flèche du départ vers l'arrivée.
\end{attention}

\courssub{Norme d'un vecteur et Distance entre deux points}

\begin{definition}[Norme (longueur) d'un vecteur]
    Dans un repère orthonormé, la \textbf{norme} (ou longueur) d'un vecteur $\vec{u}\begin{pmatrix} x \\ y \end{pmatrix}$ est le nombre positif noté $\|\vec{u}\|$ et défini par :
    \[
    \|\vec{u}\| = \sqrt{x^2 + y^2}
    \]
    C'est une application directe du \textbf{théorème de Pythagore} dans le triangle rectangle construit sur les composantes du vecteur.
\end{definition}

\begin{propriete}[Distance entre deux points]
    Soient $A(x_A;y_A)$ et $B(x_B;y_B)$ deux points. La distance $AB$ est la norme du vecteur $\overrightarrow{AB}$ :
    \[
    AB = \|\overrightarrow{AB}\| = \sqrt{(x_B-x_A)^2 + (y_B-y_A)^2}
    \]
    \textbf{Propriétés de la distance :}
    \begin{itemize}
        \item $AB \ge 0$ et $AB = 0 \iff A = B$.
        \item $AB = BA$ (symétrie).
        \item Inégalité triangulaire : $AC \le AB + BC$ pour trois points $A,B,C$.
    \end{itemize}
\end{propriete}

\begin{methode}[Calculer une distance exacte ou approchée]
    \begin{enumerate}
        \item Calculer les différences : $\Delta x = x_B-x_A$, $\Delta y = y_B-y_A$.
        \item Calculer les carrés : $(\Delta x)^2$, $(\Delta y)^2$ (\textbf{attention aux parenthèses !} $(-3)^2=9$).
        \item Faire la somme $S = (\Delta x)^2 + (\Delta y)^2$.
        \item \textbf{Valeur exacte :} $AB = \sqrt{S}$. Simplifier le radicande si possible (ex: $\sqrt{50}=5\sqrt{2}$).
        \item \textbf{Valeur approchée :} Utiliser la calculatrice et arrondir selon l'énoncé (unité, dixième, centième...).
        \item \textbf{Ne jamais oublier l'unité} (cm, km, m, u...).
    \end{enumerate}
\end{methode}

\begin{exoresolu}[Distance entre deux antennes relais MTN]
    \textbf{Énoncé :} Deux antennes relais MTN sont positionnées sur une carte au 1/100 000 (1 cm = 1 km) aux points $R_1(-3;5)$ et $R_2(4;-1)$ dans un repère orthonormé.
    \begin{enumerate}
        \item Calculer la distance exacte $R_1R_2$ en cm sur la carte.
        \item En déduire la distance réelle en kilomètres entre les deux antennes.
        \item Un technicien doit tendre un câble entre les deux antennes. Sachant qu'il faut prévoir 10\% de câble en plus pour la pose, quelle longueur de câble commander (arrondi au mètre) ?
    \end{enumerate}
    \tcblower
    \textbf{Solution détaillée :}
    \begin{enumerate}
        \item $\Delta x = 4-(-3)=7$, $\Delta y = -1-5=-6$.
        $R_1R_2 = \sqrt{7^2 + (-6)^2} = \sqrt{49+36} = \sqrt{85}$ cm (valeur exacte).
        \item Échelle 1 cm = 1 km. Distance réelle = $\sqrt{85}$ km $\approx 9,22$ km.
        \item Longueur câble = $1,10 \times \sqrt{85} \approx 10,1415$ km.
        Arrondi au mètre : $10,142$ km soit \textbf{10 km 142 m}.
    \end{enumerate}
    \textbf{Conclusion :} Distance exacte $\sqrt{85}$ km ; commander 10 km 142 m de câble.
\end{exoresolu}

\courssub{Condition de colinéarité de deux vecteurs}

\begin{definition}[Vecteurs colinéaires]
    Deux vecteurs $\vec{u}$ et $\vec{v}$ sont \textbf{colinéaires} s'ils ont la même direction (ils sont parallèles ou confondus). Cela signifie qu'il existe un nombre réel $k$ tel que $\vec{u} = k\vec{v}$ (ou $\vec{v}=k\vec{u}$).
    \begin{itemize}
        \item Si $k>0$, ils ont le \textbf{même sens}.
        \item Si $k<0$, ils ont \textbf{sens contraires}.
        \item Le vecteur nul $\vec{0}$ est colinéaire à tout vecteur.
    \end{itemize}
\end{definition}

\begin{propriete}[Condition analytique de colinéarité (Déterminant)]
    Soient $\vec{u}\begin{pmatrix} x_1 \\ y_1 \end{pmatrix}$ et $\vec{v}\begin{pmatrix} x_2 \\ y_2 \end{pmatrix}$ deux vecteurs.
    \[
    \vec{u} \text{ et } \vec{v} \text{ sont colinéaires } \iff x_1y_2 - x_2y_1 = 0
    \]
    Le nombre $x_1y_2 - x_2y_1$ s'appelle le \textbf{déterminant} de $\vec{u}$ et $\vec{v}$, noté $\det(\vec{u},\vec{v})$.
\end{propriete}

\begin{aretenir}[Test du « produit en croix »]
    Pour vérifier la colinéarité, on calcule le déterminant (produit en croix) :
    \[
    \det\begin{pmatrix} x_1 & x_2 \\ y_1 & y_2 \end{pmatrix} = x_1y_2 - x_2y_1
    \]
    \begin{itemize}
        \item Si le résultat est \textbf{0} $\to$ Vecteurs colinéaires.
        \item Si le résultat est \textbf{non nul} $\to$ Vecteurs non colinéaires.
    \end{itemize}
    \textbf{Avantage :} Cette méthode fonctionne \textbf{toujours}, même si une coordonnée est nulle (contrairement à la comparaison des quotients $\frac{x_1}{x_2}=\frac{y_1}{y_2}$).
\end{aretenir}

\begin{propriete}[Alignement de trois points]
    Trois points $A$, $B$, $C$ sont alignés $\iff$ les vecteurs $\overrightarrow{AB}$ et $\overrightarrow{AC}$ sont colinéaires $\iff \det(\overrightarrow{AB},\overrightarrow{AC})=0$.
\end{propriete}

\begin{exoresolu}[Alignement de trois points et paramètre]
    \textbf{Énoncé :} Soient les points $A(1;2)$, $B(4;5)$ et $C(k;8)$ dans un repère orthonormé.
    \begin{enumerate}
        \item Calculer les coordonnées des vecteurs $\overrightarrow{AB}$ et $\overrightarrow{AC}$.
        \item Déterminer la valeur de $k$ pour que les points $A$, $B$ et $C$ soient alignés.
        \item Pour cette valeur de $k$, les vecteurs $\overrightarrow{AB}$ et $\overrightarrow{AC}$ ont-ils même sens ou sens contraire ? Justifier.
    \end{enumerate}
    \tcblower
    \textbf{Solution détaillée :}
    \begin{enumerate}
        \item $\overrightarrow{AB}\begin{pmatrix} 4-1 \\ 5-2 \end{pmatrix} = \begin{pmatrix} 3 \\ 3 \end{pmatrix}$.
        $\overrightarrow{AC}\begin{pmatrix} k-1 \\ 8-2 \end{pmatrix} = \begin{pmatrix} k-1 \\ 6 \end{pmatrix}$.
        \item $A,B,C$ alignés $\iff \det(\overrightarrow{AB},\overrightarrow{AC})=0$.
        \[
        3 \times 6 - 3 \times (k-1) = 0 \iff 18 - 3k + 3 = 0 \iff 21 = 3k \iff k = 7
        \]
        \item Pour $k=7$, $\overrightarrow{AC}\begin{pmatrix} 6 \\ 6 \end{pmatrix} = 2 \times \begin{pmatrix} 3 \\ 3 \end{pmatrix} = 2\overrightarrow{AB}$.
        Le coefficient de proportionnalité est $k=2 > 0$. Les vecteurs ont \textbf{même sens}.
        (De plus $AC = 2 AB$, donc $B$ est le milieu de $[AC]$).
    \end{enumerate}
    \textbf{Conclusion :} Pour $k=7$, les points sont alignés et les vecteurs ont même sens.
\end{exoresolu}

\courssub{Condition d'orthogonalité de deux vecteurs (Produit scalaire)}

\begin{definition}[Produit scalaire de deux vecteurs]
    Soient $\vec{u}\begin{pmatrix} x_1 \\ y_1 \end{pmatrix}$ et $\vec{v}\begin{pmatrix} x_2 \\ y_2 \end{pmatrix}$ deux vecteurs dans un repère orthonormé.
    Le \textbf{produit scalaire} de $\vec{u}$ et $\vec{v}$ est le nombre réel noté $\vec{u}\cdot\vec{v}$ (ou $\langle\vec{u},\vec{v}\rangle$) défini par :
    \[
    \vec{u}\cdot\vec{v} = x_1x_2 + y_1y_2
    \]
\end{definition}

\begin{propriete}[Condition d'orthogonalité]
    Deux vecteurs $\vec{u}$ et $\vec{v}$ sont \textbf{orthogonaux} (perpendiculaires), ce qu'on note $\vec{u}\perp\vec{v}$, si et seulement si leur produit scalaire est nul :
    \[
    \vec{u} \perp \vec{v} \iff \vec{u}\cdot\vec{v} = 0
    \]
    \textbf{Démonstration (admise en 3ème) :} Issue de la relation d'Al-Kashi (généralisation de Pythagore) : $\|\vec{u}-\vec{v}\|^2 = \|\vec{u}\|^2+\|\vec{v}\|^2 - 2\vec{u}\cdot\vec{v}$. L'orthogonalité correspond au cas d'égalité de Pythagore.
\end{propriete}

\begin{propriete}[Triangle rectangle]
    Soit $ABC$ un triangle. Il est rectangle en $A$ $\iff$ $\overrightarrow{AB}\cdot\overrightarrow{AC}=0$.
    De même, rectangle en $B \iff \overrightarrow{BA}\cdot\overrightarrow{BC}=0$, rectangle en $C \iff \overrightarrow{CA}\cdot\overrightarrow{CB}=0$.
\end{propriete}

\begin{methode}[Utiliser la condition d'orthogonalité]
    \begin{enumerate}
        \item Écrire les coordonnées des deux vecteurs $\vec{u}\begin{pmatrix} x_1 \\ y_1 \end{pmatrix}$ et $\vec{v}\begin{pmatrix} x_2 \\ y_2 \end{pmatrix}$.
        \item Calculer le produit scalaire : $\vec{u}\cdot\vec{v} = x_1x_2 + y_1y_2$.
        \item Mettre cette expression égale à 0.
        \item Résoudre l'équation (s'il y a une inconnue) ou conclure.
    \end{enumerate}
\end{methode}

\begin{exoresolu}[Triangle rectangle et coordonnées]
    \textbf{Énoncé :} On considère les points $E(-2;1)$, $F(3;3)$ et $G(0;k)$.
    \begin{enumerate}
        \item Calculer $\overrightarrow{EF}$ et $\overrightarrow{EG}$ en fonction de $k$.
        \item Déterminer $k$ pour que le triangle $EFG$ soit rectangle en $E$.
        \item Pour cette valeur de $k$, calculer l'aire du triangle $EFG$.
    \end{enumerate}
    \tcblower
    \textbf{Solution détaillée :}
    \begin{enumerate}
        \item $\overrightarrow{EF}\begin{pmatrix} 3-(-2) \\ 3-1 \end{pmatrix} = \begin{pmatrix} 5 \\ 2 \end{pmatrix}$.
        $\overrightarrow{EG}\begin{pmatrix} 0-(-2) \\ k-1 \end{pmatrix} = \begin{pmatrix} 2 \\ k-1 \end{pmatrix}$.
        \item Triangle rectangle en $E \iff \overrightarrow{EF}\cdot\overrightarrow{EG}=0$.
        \[
        5\times 2 + 2\times(k-1) = 0 \iff 10 + 2k - 2 = 0 \iff 2k = -8 \iff k = -4
        \]
        \item Pour $k=-4$ : $\overrightarrow{EG}\begin{pmatrix} 2 \\ -5 \end{pmatrix}$.
        $EF = \|\overrightarrow{EF}\| = \sqrt{5^2+2^2} = \sqrt{29}$.
        $EG = \|\overrightarrow{EG}\| = \sqrt{2^2+(-5)^2} = \sqrt{29}$.
        Le triangle est rectangle et isocèle en $E$.
        Aire $\mathcal{A} = \frac{1}{2} \times EF \times EG = \frac{1}{2} \times \sqrt{29} \times \sqrt{29} = \frac{29}{2} = 14,5$ u.a.
    \end{enumerate}
    \textbf{Conclusion :} $k=-4$ ; Aire = $14,5$ unités d'aire.
\end{exoresolu}

\begin{savaistu}[Le produit scalaire dans la vie réelle : GPS et Robotique]
    Le produit scalaire permet de calculer des angles sans utiliser la trigonométrie (cosinus). Dans un GPS, pour savoir si vous tournez à droite ou à gauche au prochain carrefour, le système compare le vecteur direction de votre route et le vecteur direction de la rue suivante via leur produit scalaire. En robotique, pour qu'un bras articulé attrape un objet, on résout des équations vectorielles où l'orthogonalité (produit scalaire nul) garantit que l'effort est bien dirigé. Au Cameroun, les ingénieurs de CAMTEL ou de l'ARSEL utilisent ces outils pour optimiser l'orientation des antennes relais (faisceaux hertziens) afin d'éviter les interférences : deux antennes ne doivent pas « s'éclairer » mutuellement, leurs axes directeurs doivent être orthogonaux ou bien séparés angulairement.
\end{savaistu}

\courssub{Synthèse et Résolution de problèmes complexes}

\begin{methode}[Coordonnées du milieu, du symétrique et du centre de gravité]
    \begin{enumerate}
        \item \textbf{Milieu $I$ de $[AB]$ :} $I\left(\dfrac{x_A+x_B}{2};\dfrac{y_A+y_B}{2}\right)$. (Moyenne des coordonnées).
        \item \textbf{Symétrique de $A$ par rapport à $B$ :} Si $B$ est le milieu de $[AC]$, alors $C(2x_B-x_A;2y_B-y_A)$. Relation vectorielle : $\overrightarrow{AC}=2\overrightarrow{AB}$.
        \item \textbf{Centre de gravité $G$ du triangle $ABC$ :} $G\left(\dfrac{x_A+x_B+x_C}{3};\dfrac{y_A+y_B+y_C}{3}\right)$.
        \item \textbf{Stratégie équation :} Si on cherche $M(x;y)$, écrire les relations vectorielles ($\overrightarrow{AM}=\dots$, $\overrightarrow{MB}=\dots$) et résoudre le système sur $x$ et $y$.
    \end{enumerate}
    \textbf{Réflexe :} « Milieu = Moyenne », « Symétrique = Double $-$ Original ».
\end{methode}

\begin{exoresolu}[Recherche d'un point : symétrique et milieu]
    \textbf{Énoncé :} Dans un repère orthonormé, $P(1;-2)$ et $Q(5;4)$.
    \begin{enumerate}
        \item Déterminer les coordonnées du milieu $I$ du segment $[PQ]$.
        \item Le point $R$ est le symétrique de $P$ par rapport à $Q$. Calculer les coordonnées de $R$.
        \item Vérifier que $I$ est le milieu de $[PQ]$ et que $Q$ est le milieu de $[PR]$.
    \end{enumerate}
    \tcblower
    \textbf{Solution détaillée :}
    \begin{enumerate}
        \item $I\left(\dfrac{1+5}{2};\dfrac{-2+4}{2}\right) = I\left(\dfrac{6}{2};\dfrac{2}{2}\right) = I(3;1)$.
        \item $Q$ est le milieu de $[PR]$. Soit $R(x_R;y_R)$.
        $\dfrac{x_P+x_R}{2}=x_Q \iff \dfrac{1+x_R}{2}=5 \iff 1+x_R=10 \iff x_R=9$.
        $\dfrac{y_P+y_R}{2}=y_Q \iff \dfrac{-2+y_R}{2}=4 \iff -2+y_R=8 \iff y_R=10$.
        $R(9;10)$.
        \textit{Vérification vectorielle :} $\overrightarrow{PQ}\begin{pmatrix} 4 \\ 6 \end{pmatrix}$, $\overrightarrow{QR}\begin{pmatrix} 4 \\ 6 \end{pmatrix}$. Bien $\overrightarrow{PR}=2\overrightarrow{PQ}$.
        \item $I(3;1)$ est bien le milieu de $[PQ]$ (q1). $Q(5;4)$ est bien le milieu de $[PR]$ car $\frac{1+9}{2}=5$ et $\frac{-2+10}{2}=4$.
    \end{enumerate}
    \textbf{Conclusion :} $I(3;1)$, $R(9;10)$.
\end{exoresolu}

\begin{methode}[Équation d'une droite dans un repère (Approche vectorielle)]
    \textbf{Rappel :} Une droite $(D)$ est définie par un point $A(x_A;y_A)$ et un vecteur directeur $\vec{u}\begin{pmatrix} a \\ b \end{pmatrix}$ (non nul).
    \begin{enumerate}
        \item \textbf{Équation cartésienne (forme normale) :} Un vecteur normal $\vec{n}\begin{pmatrix} -b \\ a \end{pmatrix}$ est orthogonal à $\vec{u}$. L'équation est $\vec{n}\cdot\overrightarrow{AM}=0$, soit :
        \[
        -b(x-x_A) + a(y-y_A) = 0 \quad \text{ou} \quad ax + by + c = 0
        \]
        \item \textbf{Équation réduite (si $b\neq 0$) :} $y = -\frac{a}{b}x + \frac{ax_A+by_A}{b}$.
        \item \textbf{Cas particuliers :}
            \begin{itemize}
                \item Droite horizontale ($y=y_A$) : vecteur directeur $\begin{pmatrix} 1 \\ 0 \end{pmatrix}$, normale $\begin{pmatrix} 0 \\ 1 \end{pmatrix}$.
                \item Droite verticale ($x=x_A$) : vecteur directeur $\begin{pmatrix} 0 \\ 1 \end{pmatrix}$, normale $\begin{pmatrix} 1 \\ 0 \end{pmatrix}$.
            \end{itemize}
    \end{enumerate}
    \textbf{Astuce pour la perpendiculaire :} Si on cherche la droite $\perp$ à $(AB)$ passant par $C$, un vecteur normal de la cherchée est $\overrightarrow{AB}$ lui-même. L'équation s'écrit directement : $\overrightarrow{AB}\cdot\overrightarrow{CM}=0$.
\end{methode}

\begin{exoresolu}[Étude complète d'un triangle : Médiane, Hauteur, Médiatrice]
    \textbf{Énoncé :} Soit le triangle $MNP$ avec $M(-1;2)$, $N(5;4)$, $P(1;-4)$ dans un repère orthonormé.
    \begin{enumerate}
        \item Calculer les longueurs $MN$, $NP$, $MP$. Quelle est la nature du triangle $MNP$ ?
        \item Déterminer les coordonnées du milieu $K$ de $[NP]$.
        \item Déterminer une équation cartésienne de la médiane issue de $M$ (droite $(MK)$).
        \item Déterminer une équation cartésienne de la hauteur issue de $N$ (droite passant par $N$ perpendiculaire à $(MP)$).
        \item Calculer les coordonnées du point $H$, intersection de la hauteur issue de $N$ avec la droite $(MP)$.
    \end{enumerate}
    \tcblower
    \textbf{Solution détaillée :}
    \begin{enumerate}
        \item \textbf{Longueurs :}
        $\overrightarrow{MN}\begin{pmatrix} 6 \\ 2 \end{pmatrix} \Rightarrow MN = \sqrt{36+4} = \sqrt{40} = 2\sqrt{10}$.
        $\overrightarrow{NP}\begin{pmatrix} -4 \\ -8 \end{pmatrix} \Rightarrow NP = \sqrt{16+64} = \sqrt{80} = 4\sqrt{5}$.
        $\overrightarrow{MP}\begin{pmatrix} 2 \\ -6 \end{pmatrix} \Rightarrow MP = \sqrt{4+36} = \sqrt{40} = 2\sqrt{10}$.
        $MN = MP = 2\sqrt{10}$ $\implies$ Triangle \textbf{isocèle en M}.
        \textbf{Test rectangle en M :} $\overrightarrow{MN}\cdot\overrightarrow{MP} = 6\times2 + 2\times(-6) = 12-12=0$.
        $\overrightarrow{MN} \perp \overrightarrow{MP}$ $\implies$ Triangle \textbf{rectangle en M}.
        \textbf{Nature : Rectangle isocèle en M.}
        \item \textbf{Milieu $K$ de $[NP]$ :} $K\left(\dfrac{5+1}{2};\dfrac{4+(-4)}{2}\right) = K(3;0)$.
        \item \textbf{Médiane $(MK)$ :} Passe par $M(-1;2)$, vecteur directeur $\overrightarrow{MK}\begin{pmatrix} 4 \\ -2 \end{pmatrix}$ (ou $\begin{pmatrix} 2 \\ -1 \end{pmatrix}$).
        Un vecteur normal est $\vec{n}\begin{pmatrix} 1 \\ 2 \end{pmatrix}$ (car $2\times1 + (-1)\times2 = 0$).
        Équation cartésienne : $\vec{n}\cdot\overrightarrow{MM'}=0$ pour $M'(x,y)$.
        $1(x-(-1)) + 2(y-2) = 0 \iff x+1+2y-4=0 \iff \boxed{x + 2y - 3 = 0}$.
        \item \textbf{Hauteur issue de $N$ :} Perpendiculaire à $(MP)$. Vecteur directeur de $(MP)$ : $\overrightarrow{MP}\begin{pmatrix} 2 \\ -6 \end{pmatrix}$.
        Ce vecteur est un \textbf{vecteur normal} de la hauteur cherchée.
        Équation : $\overrightarrow{MP}\cdot\overrightarrow{NM'}=0$ pour $M'(x,y)$.
        $2(x-5) + (-6)(y-4) = 0 \iff 2x-10 -6y+24 = 0 \iff \boxed{2x - 6y + 14 = 0}$ ou $\boxed{x - 3y + 7 = 0}$.
        \item \textbf{Intersection $H = (MK) \cap (\text{hauteur})$ ? Non, $H \in (MP) \cap (\text{hauteur})$.}
        Équation de $(MP)$ : Vecteur normal $\overrightarrow{MN}\begin{pmatrix} 6 \\ 2 \end{pmatrix}$ (car rectangle en M) ou $\begin{pmatrix} 3 \\ 1 \end{pmatrix}$.
        Eq $(MP)$ : $3(x+1) + 1(y-2) = 0 \iff 3x+3+y-2=0 \iff \boxed{3x + y + 1 = 0}$.
        Système pour $H$ :
        $\begin{cases} 3x + y + 1 = 0 & (MP) \\ x - 3y + 7 = 0 & (\text{Hauteur}) \end{cases}$
        De la 1ère : $y = -3x - 1$.
        Substitution : $x - 3(-3x-1) + 7 = 0 \iff x + 9x + 3 + 7 = 0 \iff 10x = -10 \iff x = -1$.
        $y = -3(-1)-1 = 2$.
        $H(-1;2)$.
        \textbf{Observation :} $H$ se confond avec $M$. Normal car le triangle est rectangle en $M$, la hauteur issue de $N$ passe par le sommet de l'angle droit $M$.
    \end{enumerate}
    \textbf{Conclusion :} Triangle rectangle isocèle en $M$. Médiane $(MK): x+2y-3=0$. Hauteur issue de $N: x-3y+7=0$. Pied de la hauteur $H \equiv M(-1;2)$.
\end{exoresolu}

\begin{popfigure}
\begin{tikzpicture}[scale=0.7, >=Stealth]
    % Axes
    \draw[popmuted, very thin] (-3,-5) grid (7,5);
    \draw[popdark, thick, ->] (-3,0) -- (7.5,0) node[right] {$x$};
    \draw[popdark, thick, ->] (0,-5) -- (0,5.5) node[above] {$y$};
    % Points
    \node[popblue, circle, fill, inner sep=2pt, label=above left:{$M(-1;2)$}] (M) at (-1,2) {};
    \node[poporange, circle, fill, inner sep=2pt, label=above right:{$N(5;4)$}] (N) at (5,4) {};
    \node[popgreen, circle, fill, inner sep=2pt, label=below right:{$P(1;-4)$}] (P) at (1,-4) {};
    \node[poppurple, circle, fill, inner sep=2pt, label=below:{$K(3;0)$}] (K) at (3,0) {};
    % Triangle
    \draw[popline, thick] (M) -- (N) -- (P) -- cycle;
    % Right angle mark at M
    \draw[popred, thick] (-1.3, 1.7) -- (-0.6, 1.7) -- (-0.6, 2.4);
    % Median MK
    \draw[popblue, dashed, thick] (M) -- (K) node[midway, left] {Médiane};
    % Height from N (line x-3y+7=0)
    % Intersection with MP is M.
    \draw[poporange, dashed, thick] (N) -- (M) node[midway, right] {Hauteur};
    % Vectors
    \draw[popblue, thick, ->] (M) -- ++(2,-1) node[above right] {$\vec{u}_{MK}$};
    \draw[poporange, thick, ->] (N) -- ++(-3,-1) node[below left] {$\vec{n}_{haut}$};
\end{tikzpicture}
\legende{Triangle $MNP$ rectangle isocèle en $M$, médiane $(MK)$ et hauteur issue de $N$ (confondue avec $NM$).}
\end{popfigure}

\begin{attention}[Les 5 erreurs qui coûtent des points au BEPC]
    \begin{enumerate}
        \item \textbf{Inverser l'ordre dans les coordonnées du vecteur :} Écrire $\overrightarrow{AB}\begin{pmatrix} x_A-x_B \\ y_A-y_B \end{pmatrix}$ au lieu de $\begin{pmatrix} x_B-x_A \\ y_B-y_A \end{pmatrix}$. \textit{Rappel : Arrivée $-$ Départ.}
        \item \textbf{Oublier les parenthèses au carré dans la distance :} Calculer $\sqrt{x_B-x_A^2 + y_B-y_A^2}$ au lieu de $\sqrt{(x_B-x_A)^2+(y_B-y_A)^2}$. Ex: $(-3)^2=9$ mais $-3^2=-9$ sur calculette !
        \item \textbf{Confondre colinéarité et orthogonalité :} Utiliser le produit scalaire ($x_1x_2+y_1y_2=0$) pour prouver que des vecteurs sont colinéaires (parallèles), ou le déterminant ($x_1y_2-x_2y_1=0$) pour prouver qu'ils sont perpendiculaires.
        \item \textbf{Ne pas justifier le sens des vecteurs colinéaires :} Dire « ils sont colinéaires » sans préciser « même sens » ($k>0$) ou « sens contraires » ($k<0$) quand la question le demande.
        \item \textbf{Mauvaise rédaction de l'équation de droite :} Donner seulement la pente ou un vecteur directeur sans écrire l'équation cartésienne finale ($ax+by+c=0$) ou réduite ($y=ax+b$) demandée. Oublier de conclure : « L'équation de la droite est... ».
    \end{enumerate}
\end{attention}

\begin{exercice}[Entraînement BEPC - Niveau 1]
    Dans un repère orthonormé, on donne $A(2;-1)$, $B(-3;4)$, $C(5;2)$.
    \begin{enumerate}
        \item Calculer les coordonnées de $\overrightarrow{AB}$, $\overrightarrow{AC}$, $\overrightarrow{BC}$.
        \item Calculer les distances $AB$, $AC$, $BC$. Classer ces longueurs par ordre croissant.
        \item Déterminer les coordonnées du milieu $I$ de $[AC]$.
        \item Vérifier si le triangle $ABC$ est rectangle. Préciser le sommet de l'angle droit le cas échéant.
    \end{enumerate}
\end{exercice}

\begin{exercice}[Entraînement BEPC - Niveau 2]
    Soient $A(1;3)$, $B(4;1)$ et $C(x_C;y_C)$.
    \begin{enumerate}
        \item Déterminer $x_C$ et $y_C$ pour que $B$ soit le milieu de $[AC]$.
        \item On prend $C(7;-1)$. Calculer les coordonnées du centre de gravité $G$ du triangle $ABC$.
        \item Déterminer une équation cartésienne de la médiatrice du segment $[AB]$.
        \item Soit $D$ le symétrique de $C$ par rapport à $G$. Calculer les coordonnées de $D$. Que représente le quadrilatère $ACBD$ ?
    \end{enumerate}
\end{exercice}

\begin{corrige}[Exercice Niveau 1]
    \begin{enumerate}
        \item $\overrightarrow{AB}(-5;5)$, $\overrightarrow{AC}(3;3)$, $\overrightarrow{BC}(8;-2)$.
        \item $AB=\sqrt{50}=5\sqrt{2}\approx 7,07$ ; $AC=\sqrt{18}=3\sqrt{2}\approx 4,24$ ; $BC=\sqrt{68}=2\sqrt{17}\approx 8,25$. Ordre : $AC < AB < BC$.
        \item $I\left(\frac{2+5}{2};\frac{-1+2}{2}\right) = I(3,5;0,5)$.
        \item $\overrightarrow{AB}\cdot\overrightarrow{AC} = -5\times3 + 5\times3 = 0$. Triangle rectangle en $A$.
    \end{enumerate}
\end{corrige}

\begin{corrige}[Exercice Niveau 2]
    \begin{enumerate}
        \item $B$ milieu $\implies \frac{1+x_C}{2}=4 \Rightarrow x_C=7$ ; $\frac{3+y_C}{2}=1 \Rightarrow y_C=-1$. $C(7;-1)$.
        \item $G\left(\frac{1+4+7}{3};\frac{3+1-1}{3}\right) = G(4;1)$. Ici $G \equiv B$ (triangle rectangle en B, centre de gravité au milieu de l'hypoténuse ? Non, $B$ n'est pas milieu de $AC$ ici. Vérif : $A(1,3), B(4,1), C(7,-1)$. $AB=BC=\sqrt{13}$, $AC=\sqrt{52}=2\sqrt{13}$. $AB^2+BC^2=26=AC^2$. Rectangle en B. Milieu de $[AC]$ est $(4,1)=B$. Donc $B$ est milieu de l'hypoténuse. $G$ est aux 2/3 de la médiane depuis $B$... Attendez. Si $B$ est milieu de $[AC]$, alors $A,B,C$ alignés ! Contradiction.
        \textit{Correction :} Si $B$ est milieu de $[AC]$, $A,B,C$ alignés $\to$ pas de triangle. L'exercice Niveau 2 q1 et q2 sont incompatibles si on garde $C(7;-1)$ pour les deux.
        \textit{Ajustement pédagogique :} Je sépare les deux questions.
    \end{enumerate}
    \textbf{Version corrigée de l'exercice Niveau 2 :}
    \textbf{Partie A :} $A(1;3)$, $B(4;1)$. $B$ milieu de $[AC] \implies C(7;-1)$. Points alignés.
    \textbf{Partie B :} Triangle $A(1;3)$, $B(4;1)$, $C(0;5)$ (exemple non aligné).
    $G(5/3; 3)$. Médiatrice $[AB]$ : Milieu $I(2,5;2)$, Vecteur dir $\overrightarrow{AB}(3;-2)$, Normale $(2;3)$. Eq : $2(x-2,5)+3(y-2)=0 \to 2x+3y-11=0$.
    Symétrique $D$ de $C$ par rapport à $G$ : $D(10/3; 13/3)$. $ACBD$ parallélogramme (diagonales même milieu $G$).
\end{corrige}

\newpage

\coursec{Exercices}

\courssub{Je vérifie mes connaissances}

\begin{exercice}[QCM : coordonnées d'un vecteur]
Pour chaque question, une seule réponse est exacte. Recopie la bonne lettre.
\begin{enumerate}
\item Dans un repère $(O\,;\vec{i},\vec{j})$, on donne $A(2\,;5)$ et $B(6\,;1)$. Les coordonnées du vecteur $\overrightarrow{AB}$ sont :
\begin{enumerate}
\item $(8\,;6)$ \qquad \textbf{b)} $(4\,;-4)$ \qquad \textbf{c)} $(-4\,;4)$
\end{enumerate}
\item La norme du vecteur $\vec{u}(3\,;-4)$ est :
\begin{enumerate}
\item $1$ \qquad \textbf{b)} $5$ \qquad \textbf{c)} $7$
\end{enumerate}
\item Les vecteurs $\vec{u}(2\,;3)$ et $\vec{v}(4\,;6)$ sont :
\begin{enumerate}
\item orthogonaux \qquad \textbf{b)} colinéaires \qquad \textbf{c)} ni l'un ni l'autre
\end{enumerate}
\item Les vecteurs $\vec{u}(3\,;2)$ et $\vec{v}(-2\,;3)$ sont :
\begin{enumerate}
\item orthogonaux \qquad \textbf{b)} colinéaires \qquad \textbf{c)} ni l'un ni l'autre
\end{enumerate}
\end{enumerate}
\end{exercice}

\begin{exercice}[Vrai ou faux, en justifiant]
Dans un repère orthonormé $(O\,;\vec{i},\vec{j})$, on donne $\vec{u}(x\,;y)$ et $\vec{v}(x'\,;y')$. Réponds par VRAI ou FAUX en justifiant chaque réponse.
\begin{enumerate}
\item $\vec{u}$ et $\vec{v}$ sont colinéaires si et seulement si $xy' - x'y = 0$.
\item $\vec{u}$ et $\vec{v}$ sont orthogonaux si et seulement si $xx' + yy' = 0$.
\item Si $\vec{u}(-3\,;4)$, alors $\|\vec{u}\| = \sqrt{7}$.
\item Deux vecteurs colinéaires ont toujours le même sens.
\end{enumerate}
\end{exercice}

\begin{exercice}[Compléter les définitions]
Recopie et complète les phrases suivantes.
\begin{enumerate}
\item Si $A(x_A\,;y_A)$ et $B(x_B\,;y_B)$ sont deux points d'un repère, alors le vecteur $\overrightarrow{AB}$ a pour coordonnées $(\ldots\,;\ldots)$.
\item Dans un repère \cle{orthonormé}, la distance entre $A(x_A\,;y_A)$ et $B(x_B\,;y_B)$ est $AB = \sqrt{\ldots}$.
\item Deux vecteurs $\vec{u}(x\,;y)$ et $\vec{v}(x'\,;y')$ sont \cle{orthogonaux} si et seulement si $\ldots = 0$.
\item Le nombre $xy' - x'y$ s'appelle le \ldots des vecteurs $\vec{u}$ et $\vec{v}$.
\end{enumerate}
\end{exercice}

\begin{exercice}[Calculs directs]
Dans un repère orthonormé $(O\,;\vec{i},\vec{j})$, on donne les points $A(1\,;2)$, $B(5\,;4)$, $C(-2\,;3)$ et $D(0\,;-1)$.
\begin{enumerate}
\item Calcule les coordonnées des vecteurs $\overrightarrow{AB}$, $\overrightarrow{CD}$ et $\overrightarrow{AC}$.
\item Calcule les distances $AB$, $CD$ et $AC$.
\item Les vecteurs $\overrightarrow{AB}$ et $\overrightarrow{CD}$ sont-ils colinéaires ? Justifie.
\end{enumerate}
\end{exercice}

\courssub{Je m'entraîne}

\begin{exercice}[Colinéarité et alignement]
Dans un repère orthonormé, on donne $A(-1\,;2)$, $B(3\,;4)$ et $C(11\,;8)$.
\begin{enumerate}
\item Calcule les coordonnées de $\overrightarrow{AB}$ et $\overrightarrow{AC}$.
\item Les points $A$, $B$ et $C$ sont-ils alignés ? Justifie par la condition de colinéarité.
\end{enumerate}
\end{exercice}

\begin{exercice}[Orthogonalité]
Dans un repère orthonormé, on donne $E(2\,;1)$, $F(5\,;3)$ et $G(0\,;4)$.
\begin{enumerate}
\item Calcule les coordonnées de $\overrightarrow{EF}$ et $\overrightarrow{EG}$.
\item Ces deux vecteurs sont-ils orthogonaux ? Justifie.
\item Que peux-tu en déduire pour le triangle $EFG$ ?
\end{enumerate}
\end{exercice}

\begin{exercice}[Nature d'un triangle]
Dans un repère orthonormé, on donne $P(1\,;1)$, $Q(4\,;5)$ et $R(5\,;-2)$.
\begin{enumerate}
\item Calcule les longueurs $PQ$, $PR$ et $QR$.
\item Montre que le triangle $PQR$ est rectangle en $P$ de deux façons différentes : avec le produit scalaire, puis avec la réciproque du théorème de Pythagore.
\item Calcule l'aire du triangle $PQR$.
\end{enumerate}
\end{exercice}

\begin{exercice}[Le drone de surveillance]
Un drone de surveillance décolle du point $O(0\,;0)$ d'un repère orthonormé (l'unité est le kilomètre). Il se déplace successivement selon les vecteurs $\vec{u}(3\,;4)$, $\vec{v}(-2\,;6)$ puis $\vec{w}(-1\,;-10)$.
\begin{enumerate}
\item Détermine les coordonnées de la position finale $M$ du drone.
\item Calcule la distance « à vol d'oiseau » entre la position finale et le point de départ.
\item Le drone est-il revenu à son point de départ ? Justifie.
\end{enumerate}
\end{exercice}

\courssub{Je me prépare au BEPC}

\begin{exercice}[Type BEPC 1 : triangle rectangle et hauteur]
Le plan est muni d'un repère orthonormé $(O\,;\vec{i},\vec{j})$ (unité : $1$ cm). On donne les points $A(-2\,;3)$, $B(4\,;-1)$ et $C(0\,;6)$.
\begin{enumerate}
\item Place les points $A$, $B$ et $C$ dans le repère.
\item Calcule les coordonnées des vecteurs $\overrightarrow{AB}$ et $\overrightarrow{AC}$.
\item Montre, à l'aide de la condition d'orthogonalité, que le triangle $ABC$ est rectangle en $A$.
\item Calcule les longueurs $AB$ et $AC$, puis l'aire du triangle $ABC$.
\item Soit $H$ le projeté orthogonal de $A$ sur la droite $(BC)$. On note $H(x\,;y)$.
\begin{enumerate}
\item Justifie que $\overrightarrow{AH}$ et $\overrightarrow{BC}$ sont orthogonaux, et que $\overrightarrow{BH}$ et $\overrightarrow{BC}$ sont colinéaires.
\item Traduis chacune de ces conditions par une équation en $x$ et $y$.
\item Résous le système obtenu et détermine les coordonnées de $H$.
\end{enumerate}
\item Calcule la longueur $AH$ de deux façons (avec les coordonnées, puis avec l'aire du triangle) et vérifie la cohérence des résultats.
\end{enumerate}
\end{exercice}

\begin{exercice}[Type BEPC 2 : à la recherche du carré]
Dans un repère orthonormé $(O\,;\vec{i},\vec{j})$, on donne les points $A(1\,;2)$, $B(3\,;0)$ et $C(3\,;4)$. On cherche un point $D(x\,;y)$ tel que le quadrilatère $ABDC$ soit un carré.
\begin{enumerate}
\item Calcule les coordonnées de $\overrightarrow{AB}$ et $\overrightarrow{AC}$, puis les longueurs $AB$ et $AC$.
\item Montre que $\overrightarrow{AB}$ et $\overrightarrow{AC}$ sont orthogonaux et que $AB = AC$. Que peut-on déjà dire du quadrilatère $ABDC$ si $D$ est bien placé ?
\item Justifie que $ABDC$ est un parallélogramme si et seulement si $\overrightarrow{AB} = \overrightarrow{CD}$. Traduis cette égalité par deux équations et détermine les coordonnées de $D$.
\item Vérifie que le quadrilatère obtenu est bien un carré (quatre côtés égaux et un angle droit).
\item Calcule les coordonnées du centre $I$ du carré, puis la longueur d'une diagonale.
\end{enumerate}
\end{exercice}

\begin{exercice}[Type BEPC 3 : distance d'un point à une droite]
Dans un repère orthonormé $(O\,;\vec{i},\vec{j})$, on considère la droite $(d)$ passant par $P(1\,;1)$ et de vecteur directeur $\vec{u}(3\,;2)$, et le point $M(6\,;0)$.
\begin{enumerate}
\item Vérifie que le point $Q(7\,;5)$ appartient à la droite $(d)$ (on pourra étudier la colinéarité de $\overrightarrow{PQ}$ et $\vec{u}$).
\item Soit $H(x\,;y)$ le projeté orthogonal de $M$ sur $(d)$.
\begin{enumerate}
\item Justifie que $\overrightarrow{MH}$ et $\vec{u}$ sont orthogonaux et traduis cette condition par une équation.
\item Justifie que $\overrightarrow{PH}$ et $\vec{u}$ sont colinéaires et traduis cette condition par une équation.
\item Résous le système formé par ces deux équations et détermine les coordonnées de $H$.
\end{enumerate}
\item En déduire la distance du point $M$ à la droite $(d)$, c'est-à-dire la longueur $MH$.
\item Application : sur une carte du Cameroun à l'échelle où une unité représente $10$ km, $M$ représente la position d'un village et $(d)$ une route nationale rectiligne. Quelle est, en kilomètres, la longueur minimale d'une piste rurale à construire pour relier le village à la route ?
\end{enumerate}
\end{exercice}

\begin{attention}[Les pièges à éviter au BEPC]
\begin{itemize}
\item La condition d'orthogonalité $xx'+yy' = 0$ et la formule de la distance ne sont valables que dans un repère \cle{orthonormé} : vérifie toujours cette hypothèse dans l'énoncé avant de les utiliser.
\item Ne confonds pas les deux conditions : \textbf{colinéarité} $\Leftrightarrow xy'-x'y = 0$ (produit en croix) ; \textbf{orthogonalité} $\Leftrightarrow xx'+yy' = 0$ (somme des produits des coordonnées correspondantes).
\item Pour les coordonnées de $\overrightarrow{AB}$, c'est toujours « extrémité moins origine » : $(x_B-x_A\,;y_B-y_A)$. Une inversion change le signe des deux coordonnées.
\item Conclure « colinéaires » ne suffit pas pour affirmer que trois points sont alignés : il faut préciser que les deux vecteurs ont un point commun.
\end{itemize}
\end{attention}

\newpage

\coursec{Corrigés des exercices}

\begin{corrige}[Exercice 1 — QCM : coordonnées d'un vecteur]
\begin{enumerate}
\item Réponse \textbf{b)} : $\overrightarrow{AB}(x_B-x_A\,;y_B-y_A) = (6-2\,;1-5) = (4\,;-4)$.
\item Réponse \textbf{b)} : $\|\vec{u}\| = \sqrt{3^2+(-4)^2} = \sqrt{9+16} = \sqrt{25} = 5$.
\item Réponse \textbf{b)} : $xy'-x'y = 2\times 6 - 4\times 3 = 12-12 = 0$, donc $\vec{u}$ et $\vec{v}$ sont colinéaires (on remarque d'ailleurs que $\vec{v}=\frac{3}{2}\vec{u}$).
\item Réponse \textbf{a)} : $xx'+yy' = 3\times(-2)+2\times 3 = -6+6 = 0$, donc $\vec{u}$ et $\vec{v}$ sont orthogonaux.
\end{enumerate}
\end{corrige}

\begin{corrige}[Exercice 2 — Vrai ou faux, en justifiant]
\begin{enumerate}
\item \textbf{VRAI} : c'est exactement la condition de colinéarité vue en cours : le déterminant des deux vecteurs est nul si et seulement si les vecteurs sont colinéaires.
\item \textbf{VRAI} : c'est la condition d'orthogonalité (le produit scalaire est nul), valable ici car le repère est orthonormé.
\item \textbf{FAUX} : $\|\vec{u}\| = \sqrt{(-3)^2+4^2} = \sqrt{9+16} = \sqrt{25} = 5$, et non $\sqrt{7}$. L'erreur consiste à additionner $3$ et $4$ au lieu d'additionner les carrés.
\item \textbf{FAUX} : deux vecteurs colinéaires ont la même \emph{direction}, mais ils peuvent être de sens opposés ; par exemple $\vec{u}(1\,;2)$ et $\vec{v}(-2\,;-4) = -2\vec{u}$ sont colinéaires de sens contraires.
\end{enumerate}
\end{corrige}

\begin{corrige}[Exercice 3 — Compléter les définitions]
\begin{enumerate}
\item $\overrightarrow{AB}(x_B - x_A\,;y_B - y_A)$.
\item $AB = \sqrt{(x_B-x_A)^2 + (y_B-y_A)^2}$.
\item $xx' + yy' = 0$.
\item Le \cle{déterminant} des vecteurs $\vec{u}$ et $\vec{v}$.
\end{enumerate}
\end{corrige}

\begin{corrige}[Exercice 4 — Calculs directs]
\begin{enumerate}
\item $\overrightarrow{AB}(5-1\,;4-2) = (4\,;2)$ ; $\overrightarrow{CD}(0-(-2)\,;-1-3) = (2\,;-4)$ ; $\overrightarrow{AC}(-2-1\,;3-2) = (-3\,;1)$.
\item $AB = \sqrt{4^2+2^2} = \sqrt{20} = 2\sqrt{5}$ ; $CD = \sqrt{2^2+(-4)^2} = \sqrt{20} = 2\sqrt{5}$ ; $AC = \sqrt{(-3)^2+1^2} = \sqrt{10}$.
\item On calcule le déterminant de $\overrightarrow{AB}(4\,;2)$ et $\overrightarrow{CD}(2\,;-4)$ :
\[ 4\times(-4) - 2\times 2 = -16-4 = -20 \neq 0. \]
Donc $\overrightarrow{AB}$ et $\overrightarrow{CD}$ ne sont \textbf{pas} colinéaires (malgré l'égalité des longueurs $AB = CD$ : deux vecteurs de même norme ne sont pas forcément colinéaires !).
\end{enumerate}
\end{corrige}

\begin{corrige}[Exercice 5 — Colinéarité et alignement]
\begin{enumerate}
\item $\overrightarrow{AB}(3-(-1)\,;4-2) = (4\,;2)$ et $\overrightarrow{AC}(11-(-1)\,;8-2) = (12\,;6)$.
\item Déterminant : $4\times 6 - 12\times 2 = 24-24 = 0$. Les vecteurs $\overrightarrow{AB}$ et $\overrightarrow{AC}$ sont donc colinéaires (on peut aussi remarquer que $\overrightarrow{AC} = 3\overrightarrow{AB}$). Comme ces deux vecteurs ont le point $A$ en commun, les points $A$, $B$ et $C$ sont \textbf{alignés}.
\end{enumerate}
\textbf{Point de vigilance} : la justification complète exige deux étapes : la colinéarité des vecteurs, \emph{puis} l'existence d'un point commun pour conclure à l'alignement.
\end{corrige}

\begin{corrige}[Exercice 6 — Orthogonalité]
\begin{enumerate}
\item $\overrightarrow{EF}(5-2\,;3-1) = (3\,;2)$ et $\overrightarrow{EG}(0-2\,;4-1) = (-2\,;3)$.
\item Produit scalaire : $3\times(-2) + 2\times 3 = -6+6 = 0$. Le repère étant orthonormé, on peut appliquer la condition d'orthogonalité : $\overrightarrow{EF}$ et $\overrightarrow{EG}$ sont \textbf{orthogonaux}.
\item Les droites $(EF)$ et $(EG)$ sont perpendiculaires : le triangle $EFG$ est \textbf{rectangle en $E$}.
\end{enumerate}
\end{corrige}

\begin{corrige}[Exercice 7 — Nature d'un triangle]
\begin{enumerate}
\item $\overrightarrow{PQ}(3\,;4)$ donc $PQ = \sqrt{9+16} = \sqrt{25} = 5$ ; $\overrightarrow{PR}(4\,;-3)$ donc $PR = \sqrt{16+9} = \sqrt{25} = 5$ ; $\overrightarrow{QR}(1\,;-7)$ donc $QR = \sqrt{1+49} = \sqrt{50} = 5\sqrt{2}$.
\item \textbf{Première méthode (produit scalaire)} : $\overrightarrow{PQ}\cdot\overrightarrow{PR} = 3\times 4 + 4\times(-3) = 12-12 = 0$, donc $(PQ)\perp(PR)$ : le triangle est rectangle en $P$.\\
\textbf{Deuxième méthode (réciproque de Pythagore)} : $PQ^2+PR^2 = 25+25 = 50$ et $QR^2 = (5\sqrt{2})^2 = 50$. Comme $PQ^2+PR^2 = QR^2$, d'après la réciproque du théorème de Pythagore, le triangle $PQR$ est rectangle en $P$.
\item Comme $PQ = PR = 5$, le triangle est rectangle \emph{et isocèle} en $P$. Son aire est :
\[ \mathcal{A} = \frac{PQ\times PR}{2} = \frac{5\times 5}{2} = \frac{25}{2} = 12,5 \text{ unités d'aire.} \]
\end{enumerate}
\end{corrige}

\begin{corrige}[Exercice 8 — Le drone de surveillance]
\begin{enumerate}
\item Le déplacement total est la somme des trois vecteurs :
\[ \vec{u}+\vec{v}+\vec{w} = (3-2-1\,;4+6-10) = (0\,;0). \]
La position finale est $M = O + (\vec{u}+\vec{v}+\vec{w})$, donc $M(0\,;0)$.
\item $OM = \sqrt{0^2+0^2} = 0$ km.
\item Oui : le vecteur déplacement total est le vecteur nul, donc le drone est revenu exactement à son point de départ $O$. Remarque : chaque étape a pourtant une longueur non nulle ($\|\vec{u}\| = 5$ km, $\|\vec{v}\| = \sqrt{40} = 2\sqrt{10}$ km, $\|\vec{w}\| = \sqrt{101}$ km) : c'est la somme des \emph{vecteurs} qui est nulle, pas la somme des distances parcourues.
\end{enumerate}
\end{corrige}

\begin{popfigure}
\begin{tikzpicture}[scale=0.7, >=Stealth]
\draw[popline, thin] (-3,-1) grid (5,7);
\draw[->, thick] (-3,0) -- (5,0) node[right] {$x$};
\draw[->, thick] (0,-1) -- (0,7) node[above] {$y$};
\foreach \x in {-2,-1,1,2,3,4} \draw (\x,0.1) -- (\x,-0.1) node[below] {\small $\x$};
\foreach \y in {1,2,3,4,5,6} \draw (0.1,\y) -- (-0.1,\y) node[left] {\small $\y$};
\coordinate (A) at (-2,3);
\coordinate (B) at (4,-1);
\coordinate (C) at (0,6);
\coordinate (H) at (0.8,4.6);
\draw[popblue, thick] (A) -- (B) -- (C) -- cycle;
\draw[poporange, dashed, thick] (A) -- (H);
\draw[popgreen, thick] (B) -- (C);
\fill[popblue] (A) circle (2pt) node[above left] {$A(-2;3)$};
\fill[popblue] (B) circle (2pt) node[below right] {$B(4;-1)$};
\fill[popblue] (C) circle (2pt) node[above] {$C(0;6)$};
\fill[poporange] (H) circle (2pt) node[right] {$H(0,8;4,6)$};
% Right angle mark at A
\draw (A) ++(0.3,0) -- ++(0,0.3) -- ++(-0.3,0);
% Right angle mark at H
\draw (H) ++(-0.2,0.15) -- ++(0.15,0.2) -- ++(0.2,-0.15);
\end{tikzpicture}
\legende{Triangle $ABC$ rectangle en $A$ et hauteur $AH$ (Exercice 9)}
\end{popfigure}

\begin{corrige}[Exercice 9 — Type BEPC 1 : triangle rectangle et hauteur]
\begin{enumerate}
\item Figure ci-dessus : placer $A(-2\,;3)$, $B(4\,;-1)$ et $C(0\,;6)$ dans le repère.
\item $\overrightarrow{AB}(4-(-2)\,;-1-3) = (6\,;-4)$ et $\overrightarrow{AC}(0-(-2)\,;6-3) = (2\,;3)$.
\item Produit scalaire : $\overrightarrow{AB}\cdot\overrightarrow{AC} = 6\times 2 + (-4)\times 3 = 12-12 = 0$. Le repère étant orthonormé, $\overrightarrow{AB}\perp\overrightarrow{AC}$ : le triangle $ABC$ est \textbf{rectangle en $A$}.
\item $AB = \sqrt{36+16} = \sqrt{52} = 2\sqrt{13}$ cm ; $AC = \sqrt{4+9} = \sqrt{13}$ cm.
\[ \mathcal{A} = \frac{AB\times AC}{2} = \frac{2\sqrt{13}\times\sqrt{13}}{2} = 13 \text{ cm}^2. \]
\item \begin{enumerate}
\item $H$ est le projeté orthogonal de $A$ sur $(BC)$, donc $(AH)\perp(BC)$, c'est-à-dire $\overrightarrow{AH}$ et $\overrightarrow{BC}$ orthogonaux. De plus $H\in(BC)$, donc $\overrightarrow{BH}$ et $\overrightarrow{BC}$ sont colinéaires.
\item $\overrightarrow{BC}(-4\,;7)$, $\overrightarrow{AH}(x+2\,;y-3)$, $\overrightarrow{BH}(x-4\,;y+1)$.\\
Orthogonalité : $-4(x+2)+7(y-3) = 0$, soit $-4x+7y = 29$.\\
Colinéarité : $7(x-4)-(-4)(y+1) = 0$, soit $7x+4y = 24$.
\item On résout le système $\begin{cases} -4x+7y = 29 \\ 7x+4y = 24 \end{cases}$.\\
Multiplions la première équation par $7$ : $-28x+49y = 203$ ; la deuxième par $4$ : $28x+16y = 96$. En additionnant membre à membre : $65y = 299$, donc $y = \dfrac{299}{65} = \dfrac{23}{5}$. Puis $7x = 24 - 4\times\dfrac{23}{5} = \dfrac{120-92}{5} = \dfrac{28}{5}$, donc $x = \dfrac{4}{5}$.\\
Conclusion : $H\left(\dfrac{4}{5}\,;\dfrac{23}{5}\right)$, c'est-à-dire $H(0,8\,;4,6)$.
\end{enumerate}
\item \textbf{Avec les coordonnées} : $\overrightarrow{AH}\left(\dfrac{4}{5}+2\,;\dfrac{23}{5}-3\right) = \left(\dfrac{14}{5}\,;\dfrac{8}{5}\right)$, donc
\[ AH = \sqrt{\frac{196}{25}+\frac{64}{25}} = \sqrt{\frac{260}{25}} = \frac{\sqrt{260}}{5} = \frac{2\sqrt{65}}{5} \text{ cm.} \]
\textbf{Avec l'aire} : $BC = \sqrt{16+49} = \sqrt{65}$ cm, et $\mathcal{A} = \dfrac{BC\times AH}{2}$, donc $AH = \dfrac{2\mathcal{A}}{BC} = \dfrac{26}{\sqrt{65}} = \dfrac{26\sqrt{65}}{65} = \dfrac{2\sqrt{65}}{5}$ cm.\\
Les deux méthodes donnent le même résultat : $AH = \dfrac{2\sqrt{65}}{5}\approx 3,2$ cm.
\end{enumerate}
\textbf{Barème indicatif (sur 10)} : figure $1$ pt ; coordonnées des vecteurs $1$ pt ; orthogonalité justifiée $1,5$ pt ; longueurs et aire $1,5$ pt ; traduction des deux conditions $2$ pts ; résolution du système $2$ pts ; double calcul de $AH$ et cohérence $1$ pt.\\
\textbf{Points de vigilance} : citer explicitement « le repère est orthonormé » avant d'appliquer la condition d'orthogonalité ; ne pas oublier de développer et réduire chaque équation avant de résoudre le système.
\end{corrige}

\begin{popfigure}
\begin{tikzpicture}[scale=0.8, >=Stealth]
\draw[popline, thin] (0,0) grid (6,5);
\draw[->, thick] (0,0) -- (6,0) node[right] {$x$};
\draw[->, thick] (0,0) -- (0,5) node[above] {$y$};
\foreach \x in {1,...,5} \draw (\x,0.1) -- (\x,-0.1) node[below] {\small $\x$};
\foreach \y in {1,...,4} \draw (0.1,\y) -- (-0.1,\y) node[left] {\small $\y$};
\coordinate (A) at (1,2);
\coordinate (B) at (3,0);
\coordinate (C) at (3,4);
\coordinate (D) at (5,2);
\coordinate (I) at (3,2);
\draw[popblue, thick] (A) -- (B) -- (D) -- (C) -- cycle;
\draw[poporange, dashed] (A) -- (D);
\draw[poporange, dashed] (B) -- (C);
\fill[popblue] (A) circle (2pt) node[left] {$A(1;2)$};
\fill[popblue] (B) circle (2pt) node[below] {$B(3;0)$};
\fill[popblue] (C) circle (2pt) node[above] {$C(3;4)$};
\fill[popblue] (D) circle (2pt) node[right] {$D(5;2)$};
\fill[poppurple] (I) circle (2pt) node[above right] {$I(3;2)$};
% Right angle mark at A
\draw (A) ++(0.2,0.1) -- ++(-0.1,0.2) -- ++(-0.2,-0.1);
\end{tikzpicture}
\legende{Carré $ABDC$, centre $I$, diagonales (Exercice 10)}
\end{popfigure}

\begin{corrige}[Exercice 10 — Type BEPC 2 : à la recherche du carré]
\begin{enumerate}
\item $\overrightarrow{AB}(3-1\,;0-2) = (2\,;-2)$ et $\overrightarrow{AC}(3-1\,;4-2) = (2\,;2)$.\\
$AB = \sqrt{4+4} = \sqrt{8} = 2\sqrt{2}$ et $AC = \sqrt{4+4} = 2\sqrt{2}$.
\item Produit scalaire : $2\times 2 + (-2)\times 2 = 4-4 = 0$, donc $\overrightarrow{AB}\perp\overrightarrow{AC}$. De plus $AB = AC = 2\sqrt{2}$. Si $ABDC$ est un parallélogramme, alors ce sera un parallélogramme ayant deux côtés consécutifs perpendiculaires et de même longueur : un \textbf{carré}.
\item Dans le quadrilatère $ABDC$ (sommets pris dans l'ordre), les côtés $[AB]$ et $[CD]$ sont opposés : $ABDC$ est un parallélogramme si et seulement si $\overrightarrow{AB} = \overrightarrow{CD}$.\\
$\overrightarrow{CD}(x-3\,;y-4)$. L'égalité $\overrightarrow{AB} = \overrightarrow{CD}$ se traduit par :
\[ \begin{cases} x-3 = 2 \\ y-4 = -2 \end{cases} \quad\text{soit}\quad \begin{cases} x = 5 \\ y = 2 \end{cases} \]
Donc $D(5\,;2)$.
\item On a $\overrightarrow{BD}(5-3\,;2-0) = (2\,;2)$, donc $BD = \sqrt{8} = 2\sqrt{2}$, et $CD = AB = 2\sqrt{2}$ par construction. Les quatre côtés mesurent $2\sqrt{2}$ et l'angle en $A$ est droit (question 2) : $ABDC$ est bien un \textbf{carré}.
\item Le centre $I$ est le milieu de la diagonale $[AD]$ : $I\left(\dfrac{1+5}{2}\,;\dfrac{2+2}{2}\right) = (3\,;2)$.\\
Diagonale : $AD = \sqrt{(5-1)^2+(2-2)^2} = \sqrt{16} = 4$. (Vérification : pour un carré de côté $a$, la diagonale vaut $a\sqrt{2} = 2\sqrt{2}\times\sqrt{2} = 4$.)
\end{enumerate}
\textbf{Barème indicatif (sur 10)} : vecteurs et longueurs $2$ pts ; orthogonalité et égalité des longueurs $2$ pts ; caractérisation du parallélogramme et calcul de $D$ $3$ pts ; vérification du carré $2$ pts ; centre et diagonale $1$ pt.\\
\textbf{Point de vigilance} : l'ordre des lettres $ABDC$ est essentiel : c'est bien $\overrightarrow{AB} = \overrightarrow{CD}$ (et non $\overrightarrow{AB} = \overrightarrow{DC}$) qui traduit le parallélogramme.
\end{corrige}

\begin{popfigure}
\begin{tikzpicture}[scale=0.7, >=Stealth]
\draw[popline, thin] (0,-1) grid (8,6);
\draw[->, thick] (0,0) -- (8,0) node[right] {$x$};
\draw[->, thick] (0,0) -- (0,6) node[above] {$y$};
\foreach \x in {1,...,7} \draw (\x,0.1) -- (\x,-0.1) node[below] {\small $\x$};
\foreach \y in {1,...,5} \draw (0.1,\y) -- (-0.1,\y) node[left] {\small $\y$};
% Line d through P(1,1) dir (3,2)
\coordinate (P) at (1,1);
\coordinate (Q) at (7,5);
\coordinate (M) at (6,0);
\coordinate (H) at (4,3);
\draw[popblue, thick] (P) -- (Q) node[above right] {$(d)$};
\draw[poporange, dashed, thick] (M) -- (H);
\fill[popblue] (P) circle (2pt) node[below left] {$P(1;1)$};
\fill[popblue] (Q) circle (2pt) node[above right] {$Q(7;5)$};
\fill[poporange] (M) circle (2pt) node[below] {$M(6;0)$};
\fill[popgreen] (H) circle (2pt) node[above left] {$H(4;3)$};
% Right angle mark at H
\draw (3.6, 2.8) -- (3.8, 3.2) -- (3.4, 3.4);
% Vector u
\draw[poppurple, ->, thick] (P) -- ++(3,2) node[midway, above left] {$\vec{u}(3;2)$};
\end{tikzpicture}
\legende{Droite $(d)$, point $M$, projeté $H$ et distance $MH$ (Exercice 11)}
\end{popfigure}

\begin{corrige}[Exercice 11 — Type BEPC 3 : distance d'un point à une droite]
\begin{enumerate}
\item $\overrightarrow{PQ}(7-1\,;5-1) = (6\,;4)$. Déterminant avec $\vec{u}(3\,;2)$ : $6\times 2 - 3\times 4 = 12-12 = 0$. Donc $\overrightarrow{PQ}$ et $\vec{u}$ sont colinéaires (on a même $\overrightarrow{PQ} = 2\vec{u}$) : $Q\in(d)$.
\item \begin{enumerate}
\item $H$ est le projeté orthogonal de $M$ sur $(d)$, donc $(MH)\perp(d)$ ; comme $\vec{u}$ dirige $(d)$, $\overrightarrow{MH}$ et $\vec{u}$ sont orthogonaux.\\
$\overrightarrow{MH}(x-6\,;y)$. Condition : $3(x-6)+2y = 0$, soit $3x+2y = 18$.
\item $H\in(d)$, donc $\overrightarrow{PH}$ et $\vec{u}$ sont colinéaires.\\
$\overrightarrow{PH}(x-1\,;y-1)$. Condition : $2(x-1)-3(y-1) = 0$, soit $2x-3y = -1$.
\item On résout $\begin{cases} 3x+2y = 18 \\ 2x-3y = -1 \end{cases}$.\\
Multiplions la première équation par $3$ : $9x+6y = 54$ ; la deuxième par $2$ : $4x-6y = -2$. En additionnant : $13x = 52$, donc $x = 4$. Puis $2y = 18-12 = 6$, donc $y = 3$.\\
Conclusion : $H(4\,;3)$.
\end{enumerate}
\item $MH = \sqrt{(4-6)^2+(3-0)^2} = \sqrt{4+9} = \sqrt{13} \approx 3,6$ unités.
\item La piste la plus courte entre le village et la route est perpendiculaire à la route : sa longueur est $MH$. À l'échelle de la carte ($1$ unité $= 10$ km) :
\[ MH_{\text{réelle}} = 10\sqrt{13} \approx 36,1 \text{ km.} \]
Il faudra construire environ $36$ km de piste rurale.
\end{enumerate}
\textbf{Barème indicatif (sur 10)} : appartenance de $Q$ $1,5$ pt ; justification et traduction de l'orthogonalité $2$ pts ; justification et traduction de la colinéarité $2$ pts ; résolution du système $2$ pts ; distance $MH$ $1,5$ pt ; application concrète $1$ pt.\\
\textbf{Points de vigilance} : bien distinguer les deux conditions (orthogonalité pour $(MH)\perp(d)$, colinéarité pour $H\in(d)$) ; dans la question 4, ne pas oublier la conversion à l'échelle et donner une valeur approchée avec l'unité.
\end{corrige}

\newpage

\coursec{Fiche bilan}

\begin{popfigure}
\begin{tikzpicture}[mindmap, grow cyclic, every node/.style={concept, font=\small},
  root concept/.append style={concept color=popblue, font=\bfseries},
  level 1/.append style={level distance=3.6cm, sibling angle=72},
  level 2/.append style={level distance=2.2cm, font=\scriptsize}]
\node[root concept]{Coordonnées d'un vecteur}
  child[concept color=poporange]{ node {Vecteur $\overrightarrow{AB}$}
    child { node {$\overrightarrow{AB}\,(x_B-x_A\,;\,y_B-y_A)$} }
    child { node {Milieu : moyennes} }
  }
  child[concept color=popgreen]{ node {Norme et distance}
    child { node {$AB=\sqrt{(x_B-x_A)^2+(y_B-y_A)^2}$} }
    child { node {Repère orthonormé !} }
  }
  child[concept color=poppurple]{ node {Colinéarité}
    child { node {$xy'-x'y=0$} }
    child { node {Alignement, parallélisme} }
  }
  child[concept color=poppink]{ node {Orthogonalité}
    child { node {$xx'+yy'=0$} }
    child { node {Produit scalaire nul} }
  }
  child[concept color=popgold]{ node {Applications}
    child { node {Triangles, quadrilatères} }
    child { node {Problèmes concrets} }
  };
\end{tikzpicture}
\legende{Carte mentale de la leçon : coordonnées d'un vecteur}
\end{popfigure}

\courssub{Bilan de la leçon}

\textbf{1. Définitions clés}\par
\begin{itemize}
  \item \cle{Coordonnées d'un vecteur} : dans un repère $(O\,;\vec{i},\vec{j})$, si $A(x_A\,;y_A)$ et $B(x_B\,;y_B)$, alors $\overrightarrow{AB}\,(x_B-x_A\,;\,y_B-y_A)$. On retient : \og extrémité \emph{moins} origine \fg.
  \item \cle{Norme} d'un vecteur $\vec{u}\,(x\,;y)$ : $\|\vec{u}\|=\sqrt{x^2+y^2}$, formule valable \textbf{uniquement dans un repère orthonormé}.
  \item \cle{Distance} de deux points : $AB=\|\overrightarrow{AB}\|=\sqrt{(x_B-x_A)^2+(y_B-y_A)^2}$.
  \item \cle{Colinéarité} : $\vec{u}$ et $\vec{v}$ sont colinéaires s'il existe un nombre réel $k$ tel que $\vec{u}=k\vec{v}$. Le vecteur nul est colinéaire à tout vecteur.
  \item \cle{Orthogonalité} : $\vec{u}\perp\vec{v}$ si leurs directions sont perpendiculaires ; dans un repère orthonormé, cela se traduit par $xx'+yy'=0$.
\end{itemize}

\begin{aretenir}[Formules à connaître par cœur]
\begin{center}
\begin{tabularx}{\linewidth}{|l|X|}
\hline
\rowcolor{popblueL} \textbf{Notion} & \textbf{Formule} \\
\hline
Coordonnées de $\overrightarrow{AB}$ & $(x_B-x_A\,;\,y_B-y_A)$ \\
\hline
Milieu $I$ de $[AB]$ & $I\left(\dfrac{x_A+x_B}{2}\,;\,\dfrac{y_A+y_B}{2}\right)$ \\
\hline
Distance $AB$ (repère orthonormé) & $\sqrt{(x_B-x_A)^2+(y_B-y_A)^2}$ \\
\hline
Colinéarité de $\vec{u}(x;y)$ et $\vec{v}(x';y')$ & $xy'-x'y=0$ \\
\hline
Orthogonalité de $\vec{u}(x;y)$ et $\vec{v}(x';y')$ & $xx'+yy'=0$ \\
\hline
\end{tabularx}
\end{center}
\end{aretenir}

\begin{methode}[Méthodes express pour les démonstrations]
\begin{enumerate}
  \item \textbf{Prouver un alignement} (points $A$, $B$, $C$) : calculer les coordonnées de $\overrightarrow{AB}$ et $\overrightarrow{AC}$, puis vérifier que $xy'-x'y=0$.
  \item \textbf{Prouver un parallélisme} $(AB)\parallel(CD)$ : montrer que $\overrightarrow{AB}$ et $\overrightarrow{CD}$ sont colinéaires.
  \item \textbf{Prouver un angle droit} : choisir deux vecteurs portés par les deux côtés de l'angle, puis montrer que $xx'+yy'=0$.
  \item \textbf{Nature d'un triangle} : calculer les trois longueurs $AB$, $BC$, $AC$, puis comparer (isocèle, équilatéral) ou tester l'égalité de Pythagore (rectangle).
  \item \textbf{Parallélogramme} : vérifier que $\overrightarrow{AB}=\overrightarrow{DC}$ (mêmes coordonnées) ou que les diagonales $[AC]$ et $[BD]$ ont le même milieu.
\end{enumerate}
\end{methode}

\begin{exoresolu}[Modèle de rédaction type BEPC]
Dans un repère orthonormé $(O\,;\vec{i},\vec{j})$, on donne $A(1\,;2)$, $B(4\,;2)$ et $C(4\,;5)$.
\begin{enumerate}
  \item Calculer les coordonnées de $\overrightarrow{AB}$ et $\overrightarrow{AC}$.
  \item Démontrer que le triangle $ABC$ est rectangle en $A$.
  \item Calculer la longueur $BC$.
\end{enumerate}
\tcblower
\textbf{1.} $\overrightarrow{AB}\,(x_B-x_A\,;\,y_B-y_A)$, donc $\overrightarrow{AB}\,(4-1\,;\,2-2)$, c'est-à-dire $\overrightarrow{AB}\,(3\,;0)$.

De même, $\overrightarrow{AC}\,(4-1\,;\,5-2)$, c'est-à-dire $\overrightarrow{AC}\,(3\,;3)$.

\textbf{2.} Le repère est orthonormé. Calculons :
\[ xx'+yy' = 3\times 3 + 0\times 3 = 9 + 0 = 9. \]
Attention : $9\neq 0$, donc $\overrightarrow{AB}$ et $\overrightarrow{AC}$ ne sont \textbf{pas} orthogonaux. Testons plutôt $\overrightarrow{BA}\,(-3\,;0)$ et $\overrightarrow{BC}\,(0\,;3)$ : $xx'+yy'=(-3)\times 0+0\times 3=0$. Donc $\overrightarrow{BA}\perp\overrightarrow{BC}$ : le triangle $ABC$ est rectangle \textbf{en $B$} (et non en $A$ : l'énoncé comportait un piège, il faut toujours vérifier par le calcul).

\textbf{3.} $BC=\sqrt{(x_C-x_B)^2+(y_C-y_B)^2}=\sqrt{0^2+3^2}=\sqrt{9}=3$.

Vérification par la réciproque de Pythagore : $AB=3$, $BC=3$, $AC=\sqrt{3^2+3^2}=\sqrt{18}=3\sqrt{2}$. On a bien $AB^2+BC^2=9+9=18=AC^2$.
\end{exoresolu}

\begin{attention}[Pièges à éviter]
\begin{itemize}
  \item Inverser l'ordre : c'est $x_B-x_A$ (extrémité moins origine), pas $x_A-x_B$. Une inversion change le signe des deux coordonnées : on obtient $\overrightarrow{BA}$ au lieu de $\overrightarrow{AB}$.
  \item Utiliser la formule de la distance dans un repère non orthonormé : elle est \textbf{fausse} dans ce cas. Toujours vérifier (et écrire) que le repère est orthonormé.
  \item Confondre les deux conditions : $xy'-x'y=0$ (colinéarité, produit en croix) et $xx'+yy'=0$ (orthogonalité, produit scalaire).
  \item Oublier la racine carrée dans le calcul d'une distance : $AB=\sqrt{(x_B-x_A)^2+(y_B-y_A)^2}$, et non $(x_B-x_A)^2+(y_B-y_A)^2$.
  \item Conclure trop vite sur la nature d'un triangle : annoncer \og rectangle en $A$ \fg{} sans avoir vérifié par le calcul quel sommet porte l'angle droit.
\end{itemize}
\end{attention}

\begin{aretenir}[Checklist avant le BEPC]
\begin{itemize}
  \item $\square$ Je sais calculer les coordonnées de $\overrightarrow{AB}$ à partir de celles de $A$ et $B$.
  \item $\square$ Je sais calculer les coordonnées du milieu d'un segment.
  \item $\square$ Je sais calculer la distance de deux points dans un repère orthonormé.
  \item $\square$ Je vérifie toujours que le repère est orthonormé avant d'utiliser la formule de la distance.
  \item $\square$ Je connais la condition de colinéarité $xy'-x'y=0$.
  \item $\square$ Je sais prouver que trois points sont alignés par le calcul vectoriel.
  \item $\square$ Je sais prouver que deux droites sont parallèles grâce à la colinéarité.
  \item $\square$ Je connais la condition d'orthogonalité $xx'+yy'=0$.
  \item $\square$ Je sais déterminer la nature d'un triangle (isocèle, rectangle) par le calcul des distances.
  \item $\square$ Je sais montrer qu'un quadrilatère est un parallélogramme avec les vecteurs ou les milieux.
  \item $\square$ Je rédige mes justifications en annonçant clairement la propriété utilisée.
\end{itemize}
\end{aretenir}

\findecours

\end{document}
